% Dissertation : produire le plan détaillé d’une dissertation littéraire — Français Tle Tech — Cours signé OnBuch+
% Programme officiel MINESEC. Compiler avec : tectonic N01-dissertation-produire-plan-detaille-dissertation-litter.tex
\newcommand{\DOCMATIERE}{Français}
\newcommand{\DOCNIVEAU}{Tle Tech}
\newcommand{\DOCMODULE}{Français – classe de Terminale technique}
\newcommand{\DOCLECON}{N01}
\newcommand{\DOCTITRE}{Dissertation : produire le plan détaillé d’une dissertation littéraire}
% OnBuch+ — préambule « cours signés » ENRICHI (compilé avec tectonic / XeLaTeX)
% Système visuel « Pop » d'OnBuch+ : fond crème (jamais blanc), bordures encre
% épaisses, ombres « dures » décalées, titres Archivo Black en capitales.
% Version MATHÉMATIQUES : graphes (pgfplots/tikz), boîtes pédagogiques
% (définition, propriété, méthode, attention, le savais-tu, exercice résolu,
% exercice, TP, bilan), page de garde et signature OnBuch+ ancrée en bas.
%
% Macros attendues dans main.tex AVANT \input{preamble} :
%   \DOCMATIERE  \DOCNIVEAU  \DOCMODULE  \DOCLECON (numéro)  \DOCTITRE
\documentclass[11pt]{article}

\usepackage{fontspec}
\usepackage[french]{babel}
\usepackage[a4paper,margin=2cm,top=3.4cm,bottom=2.8cm,headheight=1.4cm,footskip=1.3cm]{geometry}
\usepackage{amsmath,amssymb,amsfonts}
\usepackage{mathtools} % \xmapsto, \coloneqq... souvent utilisés par les modèles
\usepackage{cancel} % simplifications barrées (bilans de forces)
% \abs{z} (module, valeur absolue) et \norm{u} (norme), utilisés par les modèles
\providecommand{\abs}{}\let\abs\relax\DeclarePairedDelimiter{\abs}{\lvert}{\rvert}
\providecommand{\norm}{}\let\norm\relax\DeclarePairedDelimiter{\norm}{\lVert}{\rVert}
\usepackage[table]{xcolor} % [table] : fournit aussi \rowcolors (alternance de lignes)
\usepackage{pagecolor}
\usepackage{tikz}
\usetikzlibrary{arrows.meta,positioning,calc,shapes.geometric,shapes.misc,shapes.arrows,decorations.pathmorphing,decorations.pathreplacing,decorations.markings,patterns,shadows,mindmap,trees,fit,backgrounds,matrix,intersections,angles,quotes}
\usepackage{pgfplots}
\usepackage[version=4]{mhchem} % équations nucléaires \ce{^{235}_{92}U}
\usepackage[european,siunitx]{circuitikz} % schémas électriques
\pgfplotsset{compat=1.18}
\usepgfplotslibrary{fillbetween,groupplots} % aires entre courbes (calcul intégral)
\usepackage{enumitem}
\usepackage{fancyhdr}
% [most] charge toutes les bibliothèques tcolorbox (listings, minted, raster,
% documentation...) et gonfle inutilement la mémoire TeX sur un document long
% (~300 boîtes) : on ne charge que ce qui est réellement utilisé.
\usepackage[skins,breakable]{tcolorbox}
\usepackage{graphicx}
\usepackage{colortbl}
\usepackage{booktabs}
\usepackage{tabularx}
\usepackage{multirow}
\usepackage{diagbox} % case d'en-tête diagonale des tableaux à double entrée
% Colonnes extensibles centrée / gauche / droite (C, L, R), souvent utilisées par les modèles
\newcolumntype{C}{>{\centering\arraybackslash}X}
\newcolumntype{L}{>{\raggedright\arraybackslash}X}
\newcolumntype{R}{>{\raggedleft\arraybackslash}X}
\usepackage{array}
\usepackage{multicol}
\usepackage{siunitx}
% parse-numbers=false : accepte aussi \SI{2,5 \times 10^{-3}}{...}
\sisetup{locale=FR,output-decimal-marker={,},group-minimum-digits=4,parse-numbers=false}
\DeclareSIUnit\ohm{\ensuremath{\symup{\Omega}}} % U+2126 absent de la police
\DeclareSIUnit\division{div} % réglages d'oscilloscope (ms/div, V/div)
\DeclareSIUnit\tr{tr} % tours (vitesse de rotation, ex. \SI{1500}{\tr\per\minute})
\DeclareSIUnit\molar{M} % concentration molaire (ex. \SI{3}{\molar}, \SI{1}{\micro\molar})
\DeclareSIUnit\an{an} % année (le modèle confond parfois avec \year, un compteur TeX) — ex. \SI{5}{\kilo\meter\per\an}
\DeclareSIUnit\FCFA{FCFA} % franc CFA (ex. \SI{500}{\FCFA\per\kilo\gram})
\DeclareSIUnit\degreeBrix{\textdegree Brix} % teneur en sucre d'un sirop/jus (réfractométrie)
\DeclareSIUnit\month{mois} % le modèle utilise parfois \month, un compteur TeX, comme unité de durée
\DeclareSIUnit\microliter{\micro\liter} % le modèle écrit parfois \microliter en un seul token
\DeclareSIUnit\million{millions} % dénombrement (ex. \SI{15}{\million\per\mL} spermatozoïdes)
\usepackage{qrcode}
% \degree : commande fréquemment utilisée par les modèles pour un angle ou une
% température en dehors de \SI{}{} (non fournie par les paquets chargés).
\providecommand{\degree}{^{\circ}}
% \qed (fin de démonstration), utilisé par les modèles sans amsthm
\providecommand{\qed}{\hfill$\square$}
% \barre : double barre des valeurs interdites dans un tableau de variations (tkz-tab)
\providecommand{\barre}{\ensuremath{\|}}
% environnement proof (amsthm non chargé) utilisé par les modèles
\ifdefined\proof\else\newenvironment{proof}[1][Démonstration]{\par\noindent\textit{#1.}\ }{\qed\par}\fi
\usepackage{lastpage}
\usepackage{hyperref}
\hypersetup{hidelinks}

% ============================================================
% Palette « Pop » OnBuch+ — mêmes hex que la classe P (plus_ui.dart)
% ============================================================
\definecolor{poporange}{HTML}{F59321}
\definecolor{popdark}{HTML}{E07A0C}
\definecolor{popcream}{HTML}{FFF6E8}
\definecolor{popink}{HTML}{1C1714}
\definecolor{poppurple}{HTML}{7A5AE0}
\definecolor{popgreen}{HTML}{2E9E6A}
\definecolor{poppink}{HTML}{E0457B}
\definecolor{popblue}{HTML}{2F7DE1}
\definecolor{popgold}{HTML}{C98A12}
\definecolor{popmuted}{HTML}{8A7F76}
\definecolor{popline}{HTML}{EADFD2}
\definecolor{popgray}{HTML}{8A7F76}
\colorlet{popgrey}{popgray}
% Alias tolérants (le modèle nomme parfois les couleurs différemment)
\colorlet{popwhite}{white}
\colorlet{popblack}{popink}
\colorlet{popred}{poppink}
\colorlet{popyellow}{popgold}
% Teintes claires (fonds de tableaux/encadrés) — le modèle nomme parfois
% les couleurs avec un suffixe L (ex. popblueL) sans les avoir définies.
\colorlet{poporangeL}{poporange!20}
\colorlet{popdarkL}{popdark!20}
\colorlet{poppurpleL}{poppurple!20}
\colorlet{popgreenL}{popgreen!20}
\colorlet{poppinkL}{poppink!20}
\colorlet{popblueL}{popblue!20}
\colorlet{popgoldL}{popgold!20}
\colorlet{popredL}{popred!20}
\colorlet{popgrayL}{popgray!20}
% Teintes claires pour fonds de tableaux / zones
\colorlet{popblueL}{popblue!10!white}
\colorlet{poporangeL}{poporange!14!white}
\colorlet{popgreenL}{popgreen!12!white}
\colorlet{poppurpleL}{poppurple!12!white}
\colorlet{poppinkL}{poppink!12!white}
\colorlet{popgoldL}{popgold!14!white}

\pagecolor{popcream}

% ============================================================
% Typographie
% ============================================================
\defaultfontfeatures{Ligatures=TeX}
\setmainfont{PlusJakartaSans-Medium.ttf}[Path=../../../fonts/,BoldFont=PlusJakartaSans-Bold.ttf]
% Maths en sans-serif (Fira Math) avec chiffres et lettres droites de Plus
% Jakarta, pour que les formules se fondent dans le texte.
\usepackage[math-style=ISO,bold-style=ISO]{unicode-math}
\setmathfont{FiraMath-Regular.otf}[Path=../../../fonts/]
\setmathfont{PlusJakartaSans-Medium.ttf}[Path=../../../fonts/,range={up/num,up/{Latin,latin},\mathup}]
\usepackage{icomma} % virgule décimale sans espace en mode math
% Caractères mathématiques Unicode absents des polices de texte (Plus Jakarta,
% Archivo Black) : rendus par leur équivalent mathématique, en texte comme en
% formule — sinon ils s'affichent en carré vide (ex. « Arithmétique dans ℤ »).
\usepackage{newunicodechar}
\newunicodechar{ℝ}{\ensuremath{\mathbb{R}}}
\newunicodechar{ℤ}{\ensuremath{\mathbb{Z}}}
\newunicodechar{ℂ}{\ensuremath{\mathbb{C}}}
\newunicodechar{ℕ}{\ensuremath{\mathbb{N}}}
\newunicodechar{ℚ}{\ensuremath{\mathbb{Q}}}
\newunicodechar{𝔻}{\ensuremath{\mathbb{D}}}
\newunicodechar{α}{\ensuremath{\alpha}}
\newunicodechar{β}{\ensuremath{\beta}}
\newunicodechar{γ}{\ensuremath{\gamma}}
\newunicodechar{δ}{\ensuremath{\delta}}
\newunicodechar{ε}{\ensuremath{\varepsilon}}
\newunicodechar{θ}{\ensuremath{\theta}}
\newunicodechar{λ}{\ensuremath{\lambda}}
\newunicodechar{μ}{\ensuremath{\mu}}
\newunicodechar{σ}{\ensuremath{\sigma}}
\newunicodechar{τ}{\ensuremath{\tau}}
\newunicodechar{φ}{\ensuremath{\varphi}}
\newunicodechar{ψ}{\ensuremath{\psi}}
\newunicodechar{ω}{\ensuremath{\omega}}
\newunicodechar{Σ}{\ensuremath{\Sigma}}
% majuscules produites par \MakeUppercase dans les titres (α-aminés -> Α)
\newunicodechar{Α}{\ensuremath{\alpha}}
\newunicodechar{Β}{\ensuremath{\beta}}
\newunicodechar{Γ}{\ensuremath{\Gamma}}
\newunicodechar{Δ}{\ensuremath{\Delta}}
\newunicodechar{Ε}{\ensuremath{\varepsilon}}
\newunicodechar{Θ}{\ensuremath{\Theta}}
\newunicodechar{Λ}{\ensuremath{\Lambda}}
\newunicodechar{Μ}{\ensuremath{\mu}}
\newunicodechar{Τ}{\ensuremath{\tau}}
\newunicodechar{Φ}{\ensuremath{\Phi}}
\newunicodechar{Ψ}{\ensuremath{\Psi}}
\newunicodechar{Ω}{\ensuremath{\Omega}}
\newunicodechar{↦}{\ensuremath{\mapsto}}
\newunicodechar{∈}{\ensuremath{\in}}
\newunicodechar{∉}{\ensuremath{\notin}}
\newunicodechar{∧}{\ensuremath{\wedge}}
\newunicodechar{⇔}{\ensuremath{\Leftrightarrow}}
\newunicodechar{⟺}{\ensuremath{\Longleftrightarrow}}
\newunicodechar{⇒}{\ensuremath{\Rightarrow}}
\newunicodechar{⟹}{\ensuremath{\Longrightarrow}}
\newunicodechar{∘}{\ensuremath{\circ}}
\newunicodechar{∪}{\ensuremath{\cup}}
\newunicodechar{∩}{\ensuremath{\cap}}
\newunicodechar{≡}{\ensuremath{\equiv}}
\newunicodechar{⊂}{\ensuremath{\subset}}
\newunicodechar{⊥}{\ensuremath{\perp}}
\newunicodechar{∃}{\ensuremath{\exists}}
\newunicodechar{∀}{\ensuremath{\forall}}
\newunicodechar{∛}{\ensuremath{\sqrt[3]{\,}}}
\newunicodechar{∜}{\ensuremath{\sqrt[4]{\,}}}
\newunicodechar{⃗}{\ensuremath{{}^{\to}}}
\newunicodechar{ⁿ}{\ifmmode^{n}\else\textsuperscript{\ensuremath{n}}\fi}
\newunicodechar{ˣ}{\ifmmode^{x}\else\textsuperscript{\ensuremath{x}}\fi}
\newunicodechar{ᵇ}{\ifmmode^{b}\else\textsuperscript{\ensuremath{b}}\fi}
\newunicodechar{ʳ}{\ifmmode^{r}\else\textsuperscript{\ensuremath{r}}\fi}
\newunicodechar{ᵘ}{\ifmmode^{u}\else\textsuperscript{\ensuremath{u}}\fi}
\newunicodechar{ʸ}{\ifmmode^{y}\else\textsuperscript{\ensuremath{y}}\fi}
\newunicodechar{ᵃ}{\ifmmode^{a}\else\textsuperscript{\ensuremath{a}}\fi}
\newunicodechar{ᵉ}{\ifmmode^{e}\else\textsuperscript{\ensuremath{e}}\fi}
\newunicodechar{ᵏ}{\ifmmode^{k}\else\textsuperscript{\ensuremath{k}}\fi}
\newunicodechar{ᵖ}{\ifmmode^{p}\else\textsuperscript{\ensuremath{p}}\fi}
\newunicodechar{ᴺ}{\ifmmode^{N}\else\textsuperscript{\ensuremath{N}}\fi}
\newunicodechar{ᶜ}{\ifmmode^{c}\else\textsuperscript{\ensuremath{c}}\fi}
\newunicodechar{ʰ}{\ifmmode^{h}\else\textsuperscript{\ensuremath{h}}\fi}
\newunicodechar{ᵠ}{\ifmmode^{q}\else\textsuperscript{\ensuremath{q}}\fi}
\newunicodechar{ₙ}{\ifmmode_{n}\else\textsubscript{\ensuremath{n}}\fi}
\newunicodechar{ₐ}{\ifmmode_{a}\else\textsubscript{\ensuremath{a}}\fi}
\newunicodechar{ᵢ}{\ifmmode_{i}\else\textsubscript{\ensuremath{i}}\fi}
\newunicodechar{ᵣ}{\ifmmode_{r}\else\textsubscript{\ensuremath{r}}\fi}
\newunicodechar{ₖ}{\ifmmode_{k}\else\textsubscript{\ensuremath{k}}\fi}
\newunicodechar{ᵦ}{\ifmmode_{\beta}\else\textsubscript{\ensuremath{\beta}}\fi}
\newunicodechar{ₓ}{\ifmmode_{x}\else\textsubscript{\ensuremath{x}}\fi}
\newfontfamily\popdisplay{ArchivoBlack-Regular.ttf}[Path=../../../fonts/]
\newfontfamily\poplogo{SpaceGrotesk-Bold.ttf}[Path=../../../fonts/]
\newfontfamily\popsemi{PlusJakartaSans-SemiBold.ttf}[Path=../../../fonts/]
\newfontfamily\popxbold{PlusJakartaSans-ExtraBold.ttf}[Path=../../../fonts/]
\newfontfamily\popxboldls{PlusJakartaSans-ExtraBold.ttf}[Path=../../../fonts/,LetterSpace=3.0]

\setlist[enumerate]{leftmargin=1.7em,itemsep=5pt,topsep=4pt}
\setlist[itemize]{leftmargin=1.7em,itemsep=4pt,topsep=4pt,label={\color{poporange}\textbullet}}
\setlength{\parindent}{0pt}
\setlength{\parskip}{5pt}
\renewcommand{\arraystretch}{1.25}

% pgfplots : style Pop par défaut
\pgfplotsset{
  every axis/.append style={axis line style={popink,line width=1pt},tick style={popink},
    grid style={popline,line width=0.5pt},label style={font=\small},tick label style={font=\footnotesize}},
  popcurve/.style={poporange,line width=1.8pt,no markers,smooth},
}
% Même style disponible dans un \draw TikZ simple (hors pgfplots)
\tikzset{popcurve/.style={poporange,line width=1.8pt}}
\tikzset{popgrid/.style={popline,very thin}}
\colorlet{popcurve}{poporange} % le nom du style est parfois utilisé comme couleur

% ============================================================
% Pill / squiggle
% ============================================================
\newcommand{\poppill}[3]{%
  \tikz[baseline=-0.55ex]\node[draw=popink,line width=1.2pt,rounded corners=2.6pt,%
    fill=#2,inner xsep=6.5pt,inner ysep=3pt]{%
    \popxboldls\fontsize{7}{7}\selectfont\color{#3}\MakeUppercase{#1}};%
}
\newcommand{\popsquiggle}[1]{%
  \tikz[baseline=-0.4ex]{
    \draw[#1,line width=2.6pt,line cap=round]
      plot [smooth] coordinates {(0,0.09) (0.34,0.01) (0.68,0.21) (1.02,0.01) (1.36,0.10)};
  }%
}
% Mot-clé surligné (marqueur orange sous le texte)
\newcommand{\cle}[1]{{\popxbold\color{popdark}#1}}

% ============================================================
% En-tête / pied
% ============================================================
\pagestyle{fancy}
\fancyhf{}
\renewcommand{\headrulewidth}{0pt}
\renewcommand{\footrulewidth}{0pt}
\lhead{\raisebox{-0.15ex}{\poplogo\fontsize{13}{13}\selectfont\color{popink}OnBuch\color{poporange}+}}
\rhead{\poppill{\DOCMATIERE\ \textbullet\ \DOCNIVEAU\ \textbullet\ Leçon \DOCLECON}{white}{popink}}
\lfoot{\popxbold\fontsize{7.6}{7.6}\selectfont\color{popmuted}\MakeUppercase{Cours signé OnBuch+}\ \textbullet\ {\popsemi\DOCTITRE}}
\rfoot{\tikz[baseline=-0.6ex]\node[rounded corners=8.5pt,fill=popink,minimum height=17pt,minimum width=17pt,inner xsep=5pt,inner sep=0]{\popxbold\fontsize{8}{8}\selectfont\color{popcream}\thepage\,/\,\pageref*{LastPage}};}

% ============================================================
% Titres
% ============================================================
\newcommand{\chaptitre}[2]{%
  \begin{tcolorbox}[enhanced,colback=poporange,colframe=popink,boxrule=2.6pt,arc=11pt,
      drop shadow=popink,left=15pt,right=15pt,top=12pt,bottom=11pt,before skip=3pt,after skip=18pt]
    \poppill{#2}{popink}{popcream}\par\vspace{9pt}
    {\raggedright\hyphenpenalty=10000\exhyphenpenalty=10000\popdisplay\fontsize{22}{24}\selectfont\color{popcream}\MakeUppercase{#1}\par}
  \end{tcolorbox}
}
% Section numérotée : pastille orange avec le numéro + titre Archivo
\newcounter{popsec}
\newcommand{\coursec}[1]{%
  \stepcounter{popsec}\setcounter{popsub}{0}%
  \par\vspace{16pt}\noindent
  \begin{minipage}{\linewidth}
  \tikz[baseline=-0.7ex]\node[draw=popink,line width=1.6pt,fill=poporange,rounded corners=4pt,minimum size=22pt,inner sep=0]{\popdisplay\fontsize{12}{12}\selectfont\color{popcream}\thepopsec};%
  \hspace{8pt}{\popdisplay\fontsize{15}{17}\selectfont\color{popink}\MakeUppercase{#1}}\par
  \vspace{4pt}\noindent{\color{poporange}\rule{36pt}{3.6pt}}
  \end{minipage}\par\nopagebreak\vspace{8pt}
}
\newcounter{popsub}
\newcommand{\courssub}[1]{%
  \stepcounter{popsub}%
  \par\vspace{9pt}\noindent{\popxbold\fontsize{11.5}{13}\selectfont\color{popink}{\color{poporange}\thepopsec.\thepopsub}\hspace{6pt}#1}\par\nopagebreak\vspace{4pt}
}

% ============================================================
% Boîtes pédagogiques — même recette : fond blanc, bordure épaisse colorée,
% ombre dure encre, badge plein. Titre optionnel [#1].
% ============================================================
\tcbset{popbase/.style={breakable,enhanced,colback=white,boxrule=2.3pt,arc=9pt,
  shadow={3pt}{-3pt}{0pt}{popink},left=14pt,right=14pt,top=11pt,bottom=11pt,before skip=11pt,after skip=15pt,
  fontupper=\popsemi\fontsize{10.6}{14.5}\selectfont\color{popink},
  fontlower=\popsemi\fontsize{10.6}{14.5}\selectfont\color{popink}}}
% Coche dessinée (le glyphe ✓ n'existe pas dans les polices du document)
\newcommand{\popcheck}[1]{\tikz[baseline=-0.5ex]\draw[#1,line width=1.6pt,line cap=round,line join=round](-0.13,0.0)--(-0.04,-0.09)--(0.13,0.1);}
\providecommand{\coche}{\popcheck{popgreen}}
\providecommand{\resultat}[1]{\par\smallskip\noindent\fcolorbox{popgreen}{popgreenL}{\parbox{\dimexpr\linewidth-2\fboxsep-2\fboxrule\relax}{\textbf{Résultat.} #1}}\par\smallskip}
\newcommand{\popboxhead}[3]{\poppill{#1}{#2}{white}\ifx\relax#3\relax\else\hspace{6pt}{\popxbold\fontsize{10.6}{12}\selectfont\color{popink}#3}\fi\par\vspace{7pt}}

\newtcolorbox{aretenir}[1][]{popbase,colframe=popgreen,before upper={\popboxhead{À retenir}{popgreen}{#1}}}
\newtcolorbox{exemplebox}[1][]{popbase,colframe=poporange,before upper={\popboxhead{Exemple}{poporange}{#1}}}
\newtcolorbox{definition}[1][]{popbase,colframe=popblue,before upper={\popboxhead{Définition}{popblue}{#1}}}
\newtcolorbox{propriete}[1][]{popbase,colframe=poppurple,before upper={\popboxhead{Propriété}{poppurple}{#1}}}
\newtcolorbox{methode}[1][]{popbase,colframe=popgold,before upper={\popboxhead{Méthode}{popgold}{#1}}}
\newtcolorbox{attention}[1][]{popbase,colframe=poppink,before upper={\popboxhead{Attention}{poppink}{#1}}}
\newtcolorbox{experience}[1][]{popbase,colframe=popblue,colback=popblueL,before upper={\popboxhead{Expérience}{popblue}{#1}}}
% « Le savais-tu ? » : carte violette pleine (contexte, culture, Cameroun)
\newtcolorbox{savaistu}[1][]{popbase,colback=poppurple,colframe=popink,
  fontupper=\popsemi\fontsize{10.4}{14.2}\selectfont\color{white},
  before upper={\poppill{Le savais-tu\,?}{white}{poppurple}\ifx\relax#1\relax\else\hspace{6pt}{\popxbold\color{white}#1}\fi\par\vspace{7pt}}}
% Exercice résolu : énoncé en haut, \tcblower puis la solution sur fond crème
\newcounter{popexo}\newcounter{popexr}
\newtcolorbox{exoresolu}[1][]{popbase,colframe=poporange,colbacklower=popcream,
  segmentation style={popink,line width=1.2pt,dashed},
  before upper={\stepcounter{popexr}\popboxhead{Exercice résolu \thepopexr}{poporange}{#1}},
  before lower={\poppill{Solution}{popink}{popcream}\par\vspace{6pt}}}
% Exercice (non résolu en place) — la correction est dans la partie Corrigés
\newtcolorbox{exercice}[1][]{popbase,colframe=popink,
  before upper={\stepcounter{popexo}\popboxhead{Exercice \thepopexo}{popink}{#1}}}
% Corrigé
\newtcolorbox{corrige}[1][]{popbase,colframe=popgreen,colback=popgreenL,
  before upper={\popboxhead{Corrigé}{popgreen}{#1}}}
% Objectifs / prérequis / bilan
\newtcolorbox{objectifs}{popbase,colframe=popink,colback=poporangeL,
  before upper={\popboxhead{Objectifs de la leçon}{popink}{}}}
\newtcolorbox{prerequis}{popbase,colframe=popink,colback=popblueL,
  before upper={\popboxhead{Prérequis}{popblue}{}}}
\newtcolorbox{bilan}{popbase,colframe=popink,colback=white,boxrule=2.8pt,
  before upper={\popboxhead{Fiche bilan}{poporange}{L'essentiel à connaître pour l'examen}}}
% Figure encadrée avec légende
\newtcolorbox{popfigure}[1][]{enhanced,breakable=false,colback=white,colframe=popline,boxrule=1.4pt,arc=8pt,
  left=10pt,right=10pt,top=10pt,bottom=8pt,before skip=10pt,after skip=12pt,halign=center,
  lower separated=false,halign lower=center,
  fontlower=\popsemi\fontsize{9}{11.5}\selectfont\color{popmuted},
  #1}
\newcommand{\legende}[1]{\par\vspace{6pt}{\popsemi\fontsize{9}{11.5}\selectfont\color{popmuted}#1}}

% ============================================================
% Signature OnBuch+
% ============================================================
\newcommand{\poplogoinline}[1]{{\poplogo\fontsize{#1}{#1}\selectfont\color{popink}OnBuch\color{poporange}+}}
\newcommand{\signatureonbuch}{%
  \begin{tcolorbox}[enhanced,colback=popink,colframe=popink,boxrule=2.2pt,arc=10pt,
      drop shadow=poporange,left=14pt,right=14pt,top=10pt,bottom=10pt,before skip=0pt,after skip=0pt]
    \tikz[baseline=-0.7ex]\node[circle,fill=popgreen,draw=popcream,line width=1.3pt,minimum size=22pt,inner sep=0]{\popcheck{white}};%
    \hspace{9pt}\begin{minipage}[c]{0.62\linewidth}
      {\popxbold\fontsize{12}{13}\selectfont\color{popcream}Signé\ \textbullet\ {\poplogo OnBuch\color{poporange}+}}\par\vspace{2pt}
      {\raggedright\popsemi\fontsize{8}{10.5}\selectfont\color{popline}\MakeUppercase{Équipe pédagogique OnBuch+}\\\MakeUppercase{Conforme au programme officiel MINESEC}\par}
    \end{minipage}\hfill
    {\poplogo\fontsize{22}{22}\selectfont\color{popcream}OnBuch\color{poporange}+}
  \end{tcolorbox}%
}

% ============================================================
% Page de garde
% #1 titre · #2 matière · #3 niveau · #4 module · #5 n° leçon · #6 durée ·
% #7 description / objectifs (texte ou itemize)
% ============================================================
\newcommand{\pagegarde}[7]{%
  \thispagestyle{empty}
  \newgeometry{a4paper,margin=1.7cm,top=1.6cm,bottom=1.4cm}%
  \begin{tikzpicture}[remember picture,overlay]
    \fill[poporange!18] ([xshift=-3.2cm,yshift=-3.6cm]current page.north east) circle (4.6cm);
    \fill[poppurple!14] ([xshift=2.4cm,yshift=7.6cm]current page.south west) circle (3.8cm);
  \end{tikzpicture}%
  {\poplogo\fontsize{30}{30}\selectfont\color{popink}OnBuch\color{poporange}+}\hfill
  \raisebox{0.6ex}{\poppill{Cours signé \textbullet\ Premium}{popink}{popcream}}\par
  \vspace{1pt}\popsquiggle{poppurple}\par
  \vspace{8pt}
  \begin{tcolorbox}[enhanced,colback=poporange,colframe=popink,boxrule=2.8pt,arc=13pt,
      drop shadow=popink,left=18pt,right=18pt,top=12pt,bottom=13pt,after skip=10pt]
    \poppill{#2 \textbullet\ #3}{popink}{popcream}\hspace{5pt}\poppill{#4}{white}{popink}\par\vspace{8pt}
    {\popdisplay\fontsize{14}{14}\selectfont\color{popink}LEÇON #5}\par\vspace{4pt}
    {\raggedright\hyphenpenalty=10000\exhyphenpenalty=10000\popdisplay\fontsize{25}{27}\selectfont\color{popcream}\MakeUppercase{#1}\par}
    \vspace{8pt}
    {\popxbold\fontsize{9.5}{9.5}\selectfont\color{popink}\MakeUppercase{Durée conseillée : #6}}
  \end{tcolorbox}
  \begin{tcolorbox}[enhanced,colback=white,colframe=popink,boxrule=2.2pt,arc=10pt,
      drop shadow=popink,left=16pt,right=16pt,top=10pt,bottom=10pt,after skip=8pt]
    \poppill{Dans cette leçon}{poppurple}{white}\par\vspace{5pt}
    \popsemi\fontsize{9.3}{12.6}\selectfont\color{popink}#7
  \end{tcolorbox}
  \noindent\begin{minipage}[t]{0.487\linewidth}
    \begin{tcolorbox}[enhanced,colback=popgreen,colframe=popink,boxrule=2.2pt,arc=10pt,
        drop shadow=popink,left=11pt,right=11pt,top=10pt,bottom=10pt,height=3.1cm,valign=center]
      \tcbox[colback=white,colframe=popink,boxrule=1.3pt,arc=5pt,boxsep=0pt,nobeforeafter,
        left=3pt,right=3pt,top=3pt,bottom=3pt,valign=center]{\qrcode[height=1.9cm]{https://play.google.com/store/apps/details?id=com.luvvix.onbuch&hl=fr}}%
      \hspace{8pt}\begin{minipage}[c]{0.5\linewidth}
        {\popdisplay\fontsize{11}{12.5}\selectfont\color{white}\MakeUppercase{Télécharge OnBuch}}\par\vspace{4pt}
        {\raggedright\popsemi\fontsize{8.4}{11}\selectfont\color{white}Gratuit \textbullet\ Google Play\\Tuteur IA, annales, concours, résultats\par}
      \end{minipage}
    \end{tcolorbox}
  \end{minipage}\hfill
  \begin{minipage}[t]{0.487\linewidth}
    \begin{tcolorbox}[enhanced,colback=poppurple,colframe=popink,boxrule=2.2pt,arc=10pt,
        drop shadow=popink,left=11pt,right=11pt,top=10pt,bottom=10pt,height=3.1cm,valign=center]
      \tikz[baseline=-0.7ex]\node[circle,fill=white,draw=popink,line width=1.3pt,minimum size=1.6cm,inner sep=0]{\popdisplay\fontsize{24}{24}\selectfont\color{poppurple}+};%
      \hspace{8pt}\begin{minipage}[c]{0.6\linewidth}
        {\popdisplay\fontsize{11}{12.5}\selectfont\color{white}\MakeUppercase{Passe à OnBuch+}}\par\vspace{4pt}
        {\raggedright\popsemi\fontsize{8.4}{11}\selectfont\color{white}Tous les cours signés, Tuteur IA illimité, fiches PDF — onglet OnBuch+ de l'app\par}
      \end{minipage}
    \end{tcolorbox}
  \end{minipage}
  \vspace{8pt}
  \popcontenu
  \vfill
  \signatureonbuch
  \restoregeometry
  \newpage
}

% Rangée « Ce cours contient » (page de garde)
\newcommand{\poptile}[3]{% #1 couleur #2 gros texte #3 légende
  \begin{tcolorbox}[enhanced,colback=#1,colframe=popink,boxrule=2pt,arc=9pt,shadow={2.5pt}{-2.5pt}{0pt}{popink},
      left=6pt,right=6pt,top=9pt,bottom=9pt,halign=center,valign=center,height=1.75cm,width=\linewidth,nobeforeafter]
    {\popdisplay\fontsize{12.5}{13}\selectfont\color{white}\MakeUppercase{#2}}\par\vspace{3pt}
    {\centering\popsemi\fontsize{7.8}{9.5}\selectfont\color{white}#3\par}
  \end{tcolorbox}}
\newcommand{\popcontenu}{%
  \par\noindent{\popxboldls\fontsize{7.5}{7.5}\selectfont\color{popmuted}CE COURS SIGNÉ CONTIENT}\par\vspace{7pt}
  \noindent\begin{minipage}[t]{0.235\linewidth}\poptile{popblue}{Cours}{complet, illustré, pas à pas}\end{minipage}\hfill
  \begin{minipage}[t]{0.235\linewidth}\poptile{poporange}{Méthodes}{et exercices résolus}\end{minipage}\hfill
  \begin{minipage}[t]{0.235\linewidth}\poptile{poppink}{Exercices}{type examen + corrigés}\end{minipage}\hfill
  \begin{minipage}[t]{0.235\linewidth}\poptile{popgreen}{Fiche bilan}{l'essentiel à retenir}\end{minipage}\par}

% Sommaire
\newtcolorbox{sommaire}{enhanced,colback=white,colframe=popink,boxrule=2.2pt,arc=10pt,
  drop shadow=popink,left=18pt,right=18pt,top=13pt,bottom=13pt,before skip=6pt,after skip=20pt,
  before upper={\poppill{Sommaire}{poppurple}{white}\par\vspace{9pt}\popsemi\fontsize{10.6}{14.5}\selectfont\color{popink}}}

% Signature de fin de cours — poussée tout en bas de la dernière page
\newcommand{\findecours}{%
  \par\vfill\nopagebreak
  \begin{center}\popsquiggle{poporange}\end{center}\vspace{4pt}
  \signatureonbuch
}
\AtBeginDocument{\def\llbracket{\mathopen{[\mkern-3.5mu[}}\def\rrbracket{\mathclose{]\mkern-3.5mu]}}} % intervalles d'entiers (glyphe ⟦ absent de la police)
% Glyphes absents de Fira Math : redéfinis à partir de glyphes disponibles
\AtBeginDocument{%
  \def\setminus{\mathbin{\backslash}}\let\smallsetminus\setminus
  \def\vdots{\mathord{\vbox{\baselineskip3.2pt\lineskiplimit0pt\kern4pt\hbox{.}\hbox{.}\hbox{.}}}}%
  \def\ddots{\mathinner{\mkern1mu\raise6.5pt\hbox{.}\mkern2mu\raise3.6pt\hbox{.}\mkern2mu\raise0.7pt\hbox{.}\mkern1mu}}%
  \def\triangle{\mathord{\scalebox{0.9}{$\symup{\Delta}$}}}%
  \def\nabla{\mathord{\raisebox{\depth}{\scalebox{1}[-1]{$\symup{\Delta}$}}}}%
  \def\checkmark{\popcheck{popdark}}%
  \def\star{\mathbin{\ast}}\def\bigstar{\mathord{\ast}}%
  \def\diamond{\mathbin{\scalebox{0.6}{$\Diamond$}}}%
  \def\bigcup{\mathop{\mathchoice{\raisebox{-0.3em}{\scalebox{1.7}{$\cup$}}}{\scalebox{1.25}{$\cup$}}{\cup}{\cup}}\displaylimits}%
  \def\bigcap{\mathop{\mathchoice{\raisebox{-0.3em}{\scalebox{1.7}{$\cap$}}}{\scalebox{1.25}{$\cap$}}{\cap}{\cap}}\displaylimits}%
  \def\lesseqqgtr{\mathrel{\lessgtr}}\def\bigoplus{\mathop{\oplus}\displaylimits}%
}
\providecommand{\ssi}{si et seulement si} % abréviation fréquente
\AtBeginDocument{\providecommand{\wideparen}{\overparen}} % arc de cercle

%% FIGFIX-BEGIN v5 — correctifs globaux des figures (docs/onbuch-generation/tools/figfix)
\usepackage{adjustbox}\usepackage{etoolbox}\tcbuselibrary{raster}
% toute figure TikZ/pgfplots plus large que la ligne est réduite à la largeur disponible
\BeforeBeginEnvironment{tikzpicture}{\begin{adjustbox}{max width=\linewidth}}
\AfterEndEnvironment{tikzpicture}{\end{adjustbox}}
% légendes pgfplots lisibles par-dessus les courbes
\pgfplotsset{every axis/.append style={legend style={fill=white,fill opacity=0.92,text opacity=1,draw=black!15,inner sep=3pt}}}
% étiquettes d'arêtes (syntaxe edge["..."]) avec fond blanc
\tikzset{every edge quotes/.append style={fill=white,inner sep=1.2pt}}
%% FIGFIX-END
\begin{document}

\pagegarde{Dissertation : produire le plan détaillé d’une dissertation littéraire}{Français}{Tle Tech}{Français – classe de Terminale technique}{N01}{14 séances (fiche de progression harmonisée)}{Cette leçon outille l'élève de Terminale technique pour construire, étape par étape, le plan détaillé d'une dissertation littéraire : analyse du sujet, formulation de la problématique, recherche et organisation des idées. Elle s'appuie sur la lecture du Vieux nègre et la médaille de Ferdinand Oyono et mobilise les notions de langue (référent et substituts, synonymie/antonymie, hypéronymie/hyponymie, contenus manifestes et latents, texte argumentatif) au service de la production écrite.
\vspace{4pt}
\begin{multicols}{2}\small
\begin{itemize}
\item À la fin de cette leçon, je sais analyser un sujet de dissertation (mots-clés, type de sujet, consigne implicite)
\item À la fin de cette leçon, je sais formuler une problématique claire à partir d'un sujet donné
\item À la fin de cette leçon, je sais mobiliser les contenus manifestes et latents d'un texte pour alimenter ma réflexion
\item À la fin de cette leçon, je sais rechercher, noter et trier des idées et des exemples pertinents
\item À la fin de cette leçon, je sais distinguer et choisir les types de plans (thématique, dialectique, analytique)
\item À la fin de cette leçon, je sais rédiger un plan détaillé complet (introduction, développement en parties et sous-parties, conclusion)
\end{itemize}\end{multicols}}

\begin{sommaire}
\begin{multicols}{2}
\begin{enumerate}
\item Lecture de l'œuvre support : paratexte et lecture ouvroir
\item Outils de langue I : référent, substituts, synonymie et antonymie
\item Outils de langue II : hypéronymie/hyponymie, contenus manifestes et latents, texte argumentatif
\item Analyser le sujet et formuler la problématique
\item Rechercher et organiser les idées
\item Élaborer et produire le plan détaillé
\item Activité d'intégration
\item Méthodes et exercices résolus
\item Exercices
\item Corrigés des exercices
\item Fiche bilan
\end{enumerate}
\end{multicols}
\end{sommaire}

\begin{objectifs}
\begin{itemize}
\item À la fin de cette leçon, je sais analyser un sujet de dissertation (mots-clés, type de sujet, consigne implicite)
\item À la fin de cette leçon, je sais formuler une problématique claire à partir d'un sujet donné
\item À la fin de cette leçon, je sais mobiliser les contenus manifestes et latents d'un texte pour alimenter ma réflexion
\item À la fin de cette leçon, je sais rechercher, noter et trier des idées et des exemples pertinents
\item À la fin de cette leçon, je sais distinguer et choisir les types de plans (thématique, dialectique, analytique)
\item À la fin de cette leçon, je sais rédiger un plan détaillé complet (introduction, développement en parties et sous-parties, conclusion)
\item À la fin de cette leçon, je sais utiliser la synonymie, l'antonymie, l'hypéronymie et les substituts du référent pour éviter les répétitions et enrichir mon expression
\item À la fin de cette leçon, je sais identifier les caractéristiques du texte argumentatif pour structurer ma démonstration
\item À la fin de cette leçon, je sais rédiger et justifier une démarche, une réponse ou une conclusion
\end{itemize}
\end{objectifs}

\begin{prerequis}
\begin{itemize}
\item Notions de Première sur le commentaire et l'essai argumentatif
\item Lecture intégrale du Vieux nègre et la médaille (parcours de l'œuvre au programme)
\item Classes et fonctions grammaticales (nature et fonction des mots)
\item Notion de paragraphe argumenté (thèse, argument, exemple)
\item Connecteurs logiques de base vus en Première
\end{itemize}
\end{prerequis}

\newpage

\coursec{Lecture de l'œuvre support : paratexte et lecture ouvroir}

Au concours d'entrée dans une grande école de Yaoundé, deux candidats traitent le même sujet de dissertation : « L'école camerounaise forme-t-elle encore des citoyens ? ». Le premier écrit dix pages d'idées justes mais désordonnées ; le second, en quatre pages, suit un plan clair où chaque partie répond à une étape du raisonnement. C'est le second qui obtient la meilleure note. Au Baccalauréat technique comme dans la vie professionnelle (rapports, mémoires, plaidoyers), ce n'est pas seulement ce qu'on dit qui compte, mais l'ordre dans lequel on le dit. Comment transformer un sujet et des idées éparses en un plan détaillé solide et convaincant ? C'est tout l'enjeu de cette séquence. Avant d'analyser un sujet de dissertation, nous devons d'abord maîtriser l'œuvre qui servira de réservoir d'exemples : \emph{Le Vieux nègre et la médaille} de Ferdinand Oyono. Cette première section pose les bases de la réception du texte : comprendre ce qu'est le paratexte, formuler des hypothèses de lecture, situer l'auteur et son contexte, et organiser une lecture active grâce au carnet de lecture.

\courssub{Le paratexte : définition et repérage}

\begin{definition}[Paratexte]
Le \cle{paratexte} est l'ensemble des éléments qui entourent le texte principal (le corps du roman) et qui en assurent la présentation, la réception et l'interprétation. Il constitue un \emph{seuil} (Gérard Genette, \emph{Seuils}, 1987) par lequel le lecteur entre dans l'œuvre. On distingue le \textbf{péritexte} (éléments appartenant au livre physique : couverture, titre, préface, dédicace, quatrième de couverture, notes) et l'\textbf{épitexte} (éléments extérieurs au livre : interviews, articles critiques, correspondance de l'auteur).
\end{definition}

Pour \emph{Le Vieux nègre et la médaille}, le paratexte nous livre des indices décisifs avant même la première page du récit. Observons la structure typique d'un roman édité :

\begin{popfigure}
\begin{tikzpicture}[
    node distance=1.2cm and 1.5cm,
    box/.style={draw=popdark, thick, rounded corners, fill=popcream, text width=3.2cm, align=center, minimum height=1.8cm},
    arrow/.style={-Stealth, thick, popblue},
    labelbox/.style={draw=poporange, thick, rounded corners, fill=poporangeL, text width=4.5cm, align=center, minimum height=1.2cm, font=\small}
]
% Central node
\node[box, fill=popblue!20, text width=4cm] (core) {\textbf{TEXTE PRINCIPAL}\\(Roman : chapitres 1 à n)};

% Peritext nodes
\node[box, above=of core, align=center] (cover){\textbf{COUVERTURE}\\(Illustration, titre, auteur, éditeur)};
\node[box, left=of cover, align=center] (title){\textbf{TITRE}\\\emph{Le Vieux nègre et la médaille}};
\node[box, right=of cover, align=center] (author){\textbf{AUTEUR}\\\textsc{Ferdinand Oyono}};
\node[box, below=of core, align=center] (back){\textbf{4\textsuperscript{e} DE COUVERTURE}\\(Résumé, extraits critiques)};
\node[box, left=of back, align=center] (preface){\textbf{PRÉFACE / AVANT-PROPOS}\\(Contexte, intentions, genre)};
\node[box, right=of back, align=center] (dedicace){\textbf{DÉDICACE}\\(À qui ? Pourquoi ?)};

% Arrows from peritext to core
\draw[arrow] (cover.south) -- (core.north);
\draw[arrow] (title.south east) -- (core.north west);
\draw[arrow] (author.south west) -- (core.north east);
\draw[arrow] (back.north) -- (core.south);
\draw[arrow] (preface.north east) -- (core.south west);
\draw[arrow] (dedicace.north west) -- (core.south east);

% Legend
\node[labelbox, below=1.8cm of back, text width=12cm] (legend) {\textbf{Légende :} 1. Couverture : premier contact visuel, attentes génériques. 2. Titre : clé de lecture, thème central. 3. Nom d'auteur : autorité auctoriale, biographie. 4. Quatrième de couverture : promesse de lecture, horizon d'attente. 5. Préface : cadrage interprétatif. 6. Dédicace : dimension affective, engagement personnel.};
\end{tikzpicture}
\legende{Schéma des éléments du péritexte d'un roman et de leur fonction d'orientation de la lecture (d'après Gérard Genette, \emph{Seuils}).}
\end{popfigure}

Appliquons ce schéma à notre œuvre :
\begin{itemize}
    \item \textbf{Le titre} : \emph{Le Vieux nègre et la médaille}. Il juxtapose un personnage défini par l'âge et la race (« Vieux nègre ») et un objet symbolique (« la médaille »). La conjonction « et » annonce une rencontre, un conflit ou une ironie.
    \item \textbf{L'auteur} : Ferdinand Oyono. Le nom signale une littérature africaine d'expression française, souvent engagée contre la colonisation.
    \item \textbf{La couverture} (selon les éditions : Julliard, 10/18, CLE) : souvent une illustration d'un vieil homme africain, parfois une médaille mise en évidence. L'image anticipe l'ironie finale.
    \item \textbf{La quatrième de couverture} : présente Meka, vieux chef traditionnel, décoré de la médaille pour services rendus à la France (il a donné ses fils à la guerre et ses terres à la mission). Elle annonce la chute : la médaille ne lui apporte que déception et humiliation.
    \item \textbf{La préface} (dans les éditions scolaires) : replace l'œuvre dans la littérature de la décolonisation, le réalisme critique, la dénonciation de l'assimilation.
\end{itemize}

\begin{methode}[Analyser le paratexte d'une œuvre au programme]
1. \textbf{Inventorier} tous les éléments du péritexte (couverture, page de titre, préface, 4\textsuperscript{e} de couverture, dédicace, notes).
2. \textbf{Identifier} la fonction de chacun : informer (éditeur, date), séduire (illustration, résumé), guider (préface), engager (dédicace).
3. \textbf{Formuler} des hypothèses de lecture : genre, ton, thèmes, thèse probable, public visé.
4. \textbf{Consigner} ces hypothèses dans le carnet de lecture pour les vérifier ou les infirmer pendant la lecture.
\end{methode}

\courssub{Hypothèses de lecture à partir du titre}

Le titre \emph{Le Vieux nègre et la médaille} est un condensé du programme idéologique du roman. Analysons-le mot à mot pour en extraire les hypothèses de lecture.

\begin{center}
\begin{tabularx}{\linewidth}{|l|X|X|}
\hline
\textbf{Élément du titre} & \textbf{Analyse linguistique / sémantique} & \textbf{Hypothèses de lecture générées} \\
\hline
\textbf{« Le »} (article défini) & Désigne un individu spécifique, connu ou archétypal. & Ce n'est pas « un » vieux nègre quelconque, mais \emph{le} représentant d'une condition collective : un personnage typique. \\
\hline
\textbf{« Vieux »} & Adjectif : âge avancé, expérience, autorité traditionnelle, mais aussi fragilité. & Conflit des générations ? Sagesse ou naïveté ? Respect dû contre traitement infligé. \\
\hline
\textbf{« Nègre »} & Substantif péjoratif à l'époque coloniale (terme imposé par le colonisateur). Marque l'altérité raciale et la domination. & Regard du colonisateur ? Ironie de l'auteur ? Condition de sujet colonial, non de citoyen. Identité imposée de l'extérieur. \\
\hline
\textbf{« et »} & Conjonction de coordination : liaison, addition, mais aussi opposition latente. & Relation entre l'homme et l'objet : possession ? don ? piège ? échange inégal ? \\
\hline
\textbf{« la Médaille »} & Objet métallique, symbole de récompense, de mérite, de distinction honorifique. Matérialité contre symbolique. & Que récompense-t-on ? La fidélité ? La soumission ? Le sacrifice des fils ? Valeur symbolique contre valeur réelle. Ironie : décoré mais toujours « nègre ». \\
\hline
\end{tabularx}
\end{center}

\begin{exemplebox}[Questions directrices pour la lecture]
À partir de ces hypothèses, le lecteur actif se pose les questions suivantes, qu'il notera dans son carnet :
\begin{enumerate}
    \item \textbf{Qui est ce « Vieux nègre » ?} Meka, vieux chef traditionnel de Doum. Quel est son statut réel dans la hiérarchie coloniale ?
    \item \textbf{Que vaut cette médaille ?} Valeur marchande ? Valeur honorifique ? Valeur symbolique pour les Blancs ? Pour les villageois ? Pour Meka lui-même ?
    \item \textbf{Quelle est la nature du lien « et » ?} Est-ce une acquisition fière ou une charge humiliante ? La médaille change-t-elle le statut de Meka aux yeux de l'administration ? De sa famille ? De lui-même ?
    \item \textbf{Où est l'ironie ?} Le titre même ne contient-il pas la chute : un « nègre » décoré par une puissance coloniale qui continue de nier son humanité ?
\end{enumerate}
\end{exemplebox}

\courssub{Ferdinand Oyono et le contexte historique}

\begin{savaistu}[Ferdinand Oyono (1929--2010) : de l'étudiant écrivain au diplomate d'État]
Né le 14 septembre 1929 à Ngoulemakong (région du Centre, Cameroun), Ferdinand Oyono écrit ses trois romans — \emph{Une vie de boy} (1956), \emph{Le Vieux nègre et la médaille} (1956), \emph{Chemin d'Europe} (1960) — alors qu'il est encore \textbf{étudiant en France} (droit et administration). Il a 27 ans quand paraissent les deux premiers, publiés chez Julliard à Paris.
Il devient ensuite une figure majeure de l'État camerounais : ambassadeur du Cameroun dans plusieurs pays (Liberia, France, Royaume-Uni, Algérie, entre autres) et auprès d'organisations internationales, puis ministre des Affaires étrangères (1992--1997), avant d'être nommé ministre d'État chargé de la Culture. Il meurt le 10 juin 2010 à Yaoundé. Son œuvre romanesque, brève mais fondatrice, est écrite « à chaud », dans l'urgence de la décolonisation. Il incarne la génération qui a pensé le Cameroun par la littérature avant de le servir par la politique.
\end{savaistu}

\begin{popfigure}
\begin{tikzpicture}[
    scale=0.9,
    event/.style={circle, fill=popblue, inner sep=2pt},
    eventlabel/.style={above, font=\small, text width=3.5cm, align=center, popdark},
    dateline/.style={thick, popline},
    milstone/.style={diamond, fill=popgold, inner sep=3pt},
    milstonelabel/.style={below, font=\small\bfseries, text width=4cm, align=center, popdark}
]
% Timeline axis
\draw[dateline] (0,0) -- (14,0);
% Ticks and dates
\foreach \x/\year in {0/1929, 2.5/1939, 4.5/1945, 6.5/1956, 9.5/1960, 14/2010} {
    \draw (\x,0.1) -- (\x,-0.1) node[below, font=\small] {\year};
}
% Events
\node[event] (birth) at (0,0) {};
\node[eventlabel, anchor=south, align=center] at (0,0.5){Naissance à\\Ngoulemakong (1929)};

\node[milstone] (warstart) at (2.5,0) {};
\node[milstonelabel, anchor=north, align=center] at (2.5,-0.5){Début de la 2\textsuperscript{e} Guerre\\mondiale (1939)};
\node[event] (warend) at (4.5,0) {};
\node[eventlabel, anchor=south, align=center] at (4.5,0.5){Fin de la guerre --\\Retour des tirailleurs (1945)};

\node[event] (pub) at (6.5,0) {};
\node[eventlabel, anchor=south, align=center] at (6.5,0.5){Publication d'\emph{Une vie de boy}\\et du \emph{Vieux nègre...} (1956)};

\node[milstone] (indep) at (9.5,0) {};
\node[milstonelabel, anchor=north, align=center] at (9.5,-0.5){Indépendance du\\Cameroun (1960)};

\node[event] (death) at (14,0) {};
\node[eventlabel, anchor=south] at (14,0.5) {Décès à Yaoundé (2010)};

% Context box
\node[draw=poppurple, thick, rounded corners, fill=poppurpleL, text width=6cm, align=left, font=\small, anchor=north west] at (0.5, -2.8) {
    \textbf{Contexte colonial français (1919--1960)} :\\
    -- Cameroun sous mandat SDN puis tutelle ONU (France).\\
    -- Code de l'indigénat : travail forcé, impôt de capitation, justice sommaire.\\
    -- « Politique d'assimilation » : école française, citoyenneté très rare.\\
    -- Tirailleurs : enrôlement massif, promesses non tenues.
};
\end{tikzpicture}
\legende{Frise chronologique : vie de Ferdinand Oyono, contexte colonial, Seconde Guerre mondiale et indépendance du Cameroun.}
\end{popfigure}

\textbf{Le Cameroun sous administration française (1919--1960).} Après la défaite allemande pendant la Première Guerre mondiale, le Cameroun est partagé entre la France et la Grande-Bretagne ; la France administre la plus grande partie (le Cameroun oriental) sous mandat de la SDN, puis sous tutelle de l'ONU après 1945. C'est le régime de l'\cle{indigénat} : les Africains sont des « sujets » français, non des citoyens. Ils subissent le travail forcé (les « prestations »), l'impôt de capitation, la justice sommaire des chefs aux ordres de l'administration. L'école est en grande partie confessionnelle (missions catholiques et protestantes) ; l'assimilation n'est promise qu'à une élite infime.

\textbf{La Seconde Guerre mondiale et les tirailleurs.} En 1939, la France enrôle massivement des soldats dans ses colonies d'Afrique : ce sont les \cle{tirailleurs}. On leur promet la citoyenneté, des pensions, la gloire. Ils combattent en Europe et ailleurs. Au retour, beaucoup sont démobilisés sans droits réels, sans pensions égales à celles des soldats français, oubliés. Le massacre de Thiaroye (près de Dakar, 1\textsuperscript{er} décembre 1944), où des tirailleurs réclamant leur solde sont tués par l'armée française, symbolise cette trahison. Ce contexte est le \emph{substrat historique} du roman : Meka a donné ses deux fils, morts à la guerre pour la France, et ses terres à la mission catholique ; la médaille est la « récompense » officielle de ces sacrifices.

\courssub{Lecture ouvroir : mise en place du carnet de lecture}

La « lecture ouvroir » (expression popularisée par Daniel Pennac, \emph{Comme un roman}, 1992) n'est pas une lecture passive de divertissement. C'est une \textbf{lecture active, annotée, interrogée}, qui transforme l'élève en chercheur. L'outil indispensable est le \cle{carnet de lecture} (cahier dédié ou fichier numérique structuré).

\begin{methode}[Tenir un carnet de lecture efficace (standard Bac technique)]
1. \textbf{Fiche d'identité de l'œuvre} : titre complet, auteur, date de publication, éditeur, genre, mouvement littéraire.
2. \textbf{Repérage des personnages} : une fiche par personnage principal (nom, rôle, statut social, traits physiques et psychologiques, évolution, citations clés avec numéros de page).
3. \textbf{Repérage des lieux et du temps} : carte mentale des espaces (village, ville, case, bureau de l'administration, prison) ; chronologie interne (durée de l'action, retours en arrière, ellipses).
4. \textbf{Repérage des temps forts (schéma narratif)} : situation initiale $\rightarrow$ élément déclencheur $\rightarrow$ péripéties $\rightarrow$ climax $\rightarrow$ dénouement $\rightarrow$ situation finale.
5. \textbf{Collecte de citations « à exploiter »} : pour chaque thème repéré (colonisation, dignité, trahison, langue, corps, objet), noter 2 ou 3 citations exactes (page, lignes) prêtes à être insérées dans une dissertation.
6. \textbf{Questions et réactions personnelles} : zones d'ombre, incompréhensions, indignations, liens avec l'actualité ou d'autres œuvres.
7. \textbf{Bilan thématique synthétique} : 5 ou 6 phrases résumant la thèse de l'œuvre (et non l'intrigue !).
\end{methode}

\begin{exemplebox}[Application au roman : premières fiches du carnet]
\textbf{Personnages principaux :}
\begin{itemize}
    \item \textbf{Meka} : vieux chef traditionnel de Doum. Corps usé par l'âge, mais esprit vif. Fidèle à l'administration coloniale, fier de sa loyauté. Il attend la médaille comme une consécration. Naïf sur la nature réelle du pouvoir colonial.
    \item \textbf{Kelara} : épouse de Meka. Pragmatique, lucide, elle gère le quotidien. Méfiante vis-à-vis des Blancs : la médaille, selon elle, ne nourrit personne.
    \item \textbf{Amalia} : autre épouse, plus jeune, proche des missionnaires. Représente l'influence du christianisme et l'aliénation culturelle possible.
    \item \textbf{Engamba} : frère aîné de Meka, voix de la sagesse traditionnelle et de la lucidité. Il met en garde Meka contre les illusions de la récompense coloniale.
    \item \textbf{Le commandant (l'administrateur)} : figure du pouvoir colonial : distant, rituel, manipulateur, méprisant sous les dehors de la courtoisie.
    \item \textbf{Le père Vandermayer} : missionnaire, figure de l'Église coloniale, bénéficiaire des terres données par Meka.
\end{itemize}

\textbf{Lieux clés :}
\begin{itemize}
    \item \textbf{Doum} : le village, espace de la tradition, de la parole, de la communauté (place publique, case de Meka).
    \item \textbf{La ville} : espace du pouvoir blanc, de l'administration, de l'église, du marché, mais aussi de la misère et de la prison. Lieu de l'altérité et de l'humiliation finale.
\end{itemize}

\textbf{Temps forts (schéma narratif simplifié) :}
\begin{enumerate}
    \item \textbf{Situation initiale} : Meka apprend qu'il va recevoir une médaille pour services rendus à la France (ses deux fils morts à la guerre, ses terres données à la mission). Préparatifs, fierté familiale et villageoise.
    \item \textbf{Élément déclencheur} : le voyage vers la ville pour la cérémonie officielle.
    \item \textbf{Péripéties} : séjour en ville ; découvertes décevantes ; tensions familiales ; avertissements d'Engamba et de Kelara ; attente de la cérémonie.
    \item \textbf{Climax} : la cérémonie de remise de la médaille, avec le discours officiel de l'administration ; joie apparente de Meka, fête nocturne.
    \item \textbf{Dénouement} : ivre et égaré après la fête, Meka est arrêté par la police coloniale, passe la nuit en prison parmi les délinquants, nu et humilié ; la médaille est perdue.
    \item \textbf{Situation finale} : Meka, libéré, prend conscience de la vanité de la récompense : la médaille n'a rien changé à sa condition. La lucidité succède à l'illusion.
\end{enumerate}
\end{exemplebox}

\courssub{Premières pistes thématiques : du paratexte au cœur du sujet}

La lecture ouvroir, croisée avec les hypothèses du paratexte, fait émerger les grands thèmes qui nourriront nos futures dissertations. Ce ne sont pas des « sujets de dissertation » tout faits, mais des \textbf{entrées thématiques} (réservoirs d'arguments et d'exemples).

\begin{center}
\begin{tabularx}{\linewidth}{|l|X|X|}
\hline
\textbf{Thème} & \textbf{Manifestation dans l'œuvre (exemples précis)} & \textbf{Angle d'attaque dissertation (problématique possible)} \\
\hline
\textbf{La colonisation comme système de duperie} & Indigénat, impôts, travail forcé, justice inique. L'administration « donne » la médaille mais maintient l'oppression. & \emph{La récompense honorifique peut-elle masquer l'exploitation coloniale ?} \\
\hline
\textbf{L'illusion de la récompense / l'objet symbolique} & La médaille : objet brillant, sans valeur d'usage. Elle ne protège ni de la faim, ni de l'impôt, ni de la prison. & \emph{Dans quelle mesure le symbole (médaille, diplôme, décor) masque-t-il la réalité du pouvoir ?} \\
\hline
\textbf{Le conflit des générations et des valeurs} & Les fils morts à la guerre ; Engamba et Kelara, lucides, contre Meka, crédule ; les jeunes attirés par la ville. & \emph{L'école et l'administration coloniales divisent-elles la communauté africaine ?} \\
\hline
\textbf{La dignité retrouvée / la lucidité tardive} & La nuit en prison, la nudité, la perte de la médaille $\rightarrow$ chute du mythe. Meka redevient un homme, libéré du leurre. & \emph{La perte des illusions est-elle la condition de la reconquête de la dignité ?} \\
\hline
\textbf{Le corps et la langue} & Corps marqué (vieillesse, nudité de la prison). Langue : français administratif (langue du pouvoir) contre langue vernaculaire (langue de la vérité, proverbes). & \emph{Comment l'œuvre met-elle en scène la violence symbolique exercée sur les corps et les langues ?} \\
\hline
\end{tabularx}
\end{center}

\begin{attention}[Piège majeur : résumer au lieu d'exploiter]
\emph{Erreur fréquente au Bac :} traiter un sujet de dissertation (ex. « Le pouvoir colonial dans \emph{Le Vieux nègre et la médaille} ») en racontant l'histoire de Meka chapitre par chapitre. \textbf{Non.} En dissertation, l'œuvre est un \textbf{réservoir d'exemples}. On ne résume pas ; on \textbf{argumente} : « La cérémonie de remise de la médaille illustre le théâtre du pouvoir colonial : l'administrateur célèbre la "fidélité" de Meka tout en maintenant le régime de l'indigénat. » Chaque paragraphe du développement = une idée-force + un exemple précis (citation ou scène) + une analyse du lien entre l'idée et l'exemple.
\end{attention}

\courssub{Négociation du projet d'études avec la classe}

La réussite de cette séquence suppose un contrat didactique clair entre l'enseignant et le groupe. Voici le cadre proposé (à adapter selon le calendrier de l'établissement) :

\begin{center}
\begin{tabularx}{\linewidth}{|c|X|X|}
\hline
\textbf{Étape} & \textbf{Activités / productions attendues} & \textbf{Évaluation / jalons} \\
\hline
\textbf{S1--S2} (cette section) & Paratexte, hypothèses, contexte, mise en place du carnet. & Carnet initialisé (fiche d'identité + fiches personnages + schéma narratif prévisionnel). \\
\hline
\textbf{S3--S4} & Lecture guidée : le village, l'annonce, les préparatifs. & Quiz de lecture (10 questions factuelles) + 3 citations collectées sur le thème « attente / illusion ». \\
\hline
\textbf{S5--S6} & Lecture guidée : la ville, l'administration, les avertissements. & Fiche lieu « La ville » + analyse d'une scène clé (mépris, inégalités). \\
\hline
\textbf{S7--S8} & Lecture guidée : cérémonie, fête, prison, chute. & \textbf{Contrôle de lecture écrit (1 h)} : 5 questions de cours + 1 question d'analyse (type Bac). \\
\hline
\textbf{S9--S10} & Synthèse thématique (tableau ci-dessus). Débat : « Meka est-il responsable de son malheur ? ». & Compte-rendu oral de synthèse (groupe de 4, 5 min) noté sur 20 (expression orale). \\
\hline
\textbf{S11--S14} & \textbf{Transfert vers la dissertation} : analyse de sujets types, recherche des idées, plans détaillés, rédaction de l'introduction et de la conclusion. & \textbf{Devoir à la maison : dissertation complète (sujet au choix parmi 3)} notée sur 20 (coef. 2). \\
\hline
\end{tabularx}
\end{center}

\textbf{Consignes pour le carnet de lecture (obligatoire, vérifié en S4, S8 et S10) :}
\begin{itemize}
    \item Format A4, 50 pages minimum, écriture lisible, sommaire tenu à jour.
    \item Une couleur par thème (ex. : \textcolor{popblue}{Colonisation}, \textcolor{poporange}{Dignité}, \textcolor{popgreen}{Langue et corps}, \textcolor{poppurple}{Objets et symboles}).
    \item Minimum 15 citations exploitées (numéro de page + 2 lignes d'analyse).
\end{itemize}

\courssub{Exercice résolu : analyser le paratexte d'une édition donnée}

\begin{exoresolu}[Analyse d'une quatrième de couverture]
\textbf{Énoncé :} Voici le texte d'une quatrième de couverture du roman : « Meka, vieux chef traditionnel, attend la plus haute distinction honorifique : la médaille. Il a bien servi la France : ses fils sont morts pour elle, ses terres ont été données à la mission. Mais la médaille ne lui apporte que déception et humiliation. Oyono décrit avec une ironie féroce le mécanisme colonial qui broie les hommes. » \textbf{Analysez ce paratexte en montrant comment il oriente la lecture.}

\tcblower
\textbf{Solution détaillée :}
\begin{enumerate}
    \item \textbf{Identification du genre et du ton} : « ironie féroce », « mécanisme colonial qui broie les hommes » $\rightarrow$ le texte signale un roman \textbf{réaliste critique}, \textbf{engagé}, relevant de la \textbf{littérature anticoloniale}. Le lecteur attend une dénonciation, pas une célébration.
    \item \textbf{Présentation du personnage principal} : « Meka, vieux chef traditionnel » $\rightarrow$ statut social précis (intermédiaire de l'administration coloniale), âge (autorité et vulnérabilité). « Il a bien servi la France : ses fils sont morts pour elle » $\rightarrow$ donne la \textbf{cause} de la médaille (le sacrifice filial) et pose le \textbf{contrat implicite} : service rendu = récompense due.
    \item \textbf{Annonce de la chute} : « Mais la médaille ne lui apporte que déception et humiliation » $\rightarrow$ le paratexte lève le suspense sur l'issue \textbf{matérielle} (il recevra la médaille) mais ouvre le suspense \textbf{sémantique} : \emph{pourquoi} la déception ? \emph{comment} l'humiliation ? Cela crée une \textbf{lecture en tension} : le lecteur connaît la fin, il guette le processus.
    \item \textbf{Fonction herméneutique} : ce texte agit comme un \textbf{mode d'emploi}. Il dit : « Lisez ce roman comme une démonstration de l'absurdité coloniale. Ne vous laissez pas duper par la fierté de Meka au début ; l'auteur, lui, sait que c'est un piège. » Il instaure une \textbf{complicité critique} entre l'auteur et le lecteur, au-dessus du personnage.
    \item \textbf{Liens avec le titre} : « plus haute distinction honorifique » explicite « la médaille » du titre ; « vieux chef traditionnel » explicite « le vieux nègre » (le titre reprend le terme raciste du colonisateur, la quatrième de couverture utilise le terme administratif).
\end{enumerate}
\textbf{Bilan :} ce paratexte est \textbf{directif} : il verrouille l'interprétation dans le sens d'une critique radicale du colonialisme et interdit une lecture naïve (« la belle histoire d'un brave homme décoré »). C'est typique des éditions scolaires, qui encadrent la réception des élèves.
\end{exoresolu}

\begin{aretenir}[Synthèse de la section 1]
\begin{itemize}
    \item Le \textbf{paratexte} (titre, couverture, 4\textsuperscript{e} de couverture, préface, dédicace) est un \textbf{seuil herméneutique} : il oriente la lecture avant le texte. Il faut l'analyser \textbf{avant} de lire.
    \item \emph{Le Vieux nègre et la médaille} : titre ironique, choc des registres (« nègre » + « médaille »). Hypothèses : duperie coloniale, valeur nulle du symbole, conflit des générations, dignité blessée.
    \item \textbf{Ferdinand Oyono} (1929--2010) : écrivain camerounais né à Ngoulemakong, romancier à 27 ans (\emph{Une vie de boy} et \emph{Le Vieux nègre et la médaille}, 1956), puis diplomate et ministre. Contexte : Cameroun sous administration française, indigénat, tirailleurs oubliés de 1939--1945.
    \item \textbf{Lecture ouvroir} = lecture active avec \textbf{carnet de lecture} structuré (personnages, lieux, schéma narratif, citations thématisées, questions).
    \item \textbf{Thèmes réservoirs} : colonisation et duperie, illusion de la récompense, fidélité et trahison, dignité et lucidité, corps et langue.
    \item \textbf{Projet d'études} : calendrier de 14 séances, carnet vérifié, contrôle de lecture, compte-rendu oral, dissertation finale.
    \item \textbf{Règle d'or} : en dissertation, l'œuvre = \textbf{réservoir d'exemples analysés}, \textbf{jamais} résumé d'intrigue.
\end{itemize}
\end{aretenir}

\coursec{Outils de langue I : référent, substituts, synonymie et antonymie}

Dans la section précédente, nous avons franchi le seuil du \emph{Vieux nègre et la médaille} de Ferdinand Oyono : nous avons étudié le paratexte (titre, couverture, quatrième de couverture) et mené une lecture \cle{ouvroir}, c'est-à-dire une première lecture qui ouvre le texte et prépare le travail d'analyse. Au cours de cette lecture, une observation s'impose : le héros, Meka, n'est presque jamais désigné deux fois de la même manière. Tantôt « Meka », tantôt « le vieillard », tantôt « il », tantôt « le vieux nègre ». Pourquoi l'écrivain varie-t-il ainsi ses désignations ? Et comment faire de même dans nos propres dissertations ? C'est à ces questions que répondent les outils de langue de cette section : le \cle{référent} et ses \cle{substituts}, la \cle{synonymie} et l'\cle{antonymie}.

\courssub{Le référent : ce dont on parle}

Commençons par une expérience de pensée simple. Un élève dit : « Il a reçu une médaille. » Son camarade, qui n'a pas entendu le début de la conversation, demande : « Qui ? » Le mot \emph{il} renvoie à quelque chose qui existe dans la situation ou dans le discours précédent, mais qui n'est pas dit dans ce mot lui-même. Ce « quelque chose » — la personne réelle désignée —, c'est le \cle{référent}.

\begin{definition}[Le référent]
Le \cle{référent} est la réalité extralinguistique — personne, objet, animal, notion, événement — désignée par un mot ou un groupe de mots dans un énoncé. C'est \emph{ce dont on parle}. Le mot ou le groupe de mots qui désigne le référent s'appelle le \cle{référant} (ou expression référentielle).
\end{definition}

Il faut bien distinguer trois choses :

\begin{itemize}
\item la \textbf{réalité} : Meka, le vieux paysan d'Ekoumdoum, personnage du roman ;
\item le \textbf{mot} qui le désigne : « Meka », « le vieillard », « il » ;
\item le \textbf{sens} que ce mot ajoute : « le vieux nègre » n'est pas neutre, il porte le regard méprisant de l'administration coloniale.
\end{itemize}

\begin{exemplebox}[Un même référent, plusieurs désignations]
Dans la phrase : « \textbf{Meka} attendait depuis longtemps ; \textbf{le vieillard} tremblait d'émotion ; \textbf{il} tenait sa médaille contre sa poitrine », les trois expressions soulignées ont un seul et même \cle{référent} : le personnage de Meka. Mais elles ne disent pas la même chose : le nom propre identifie, « le vieillard » décrit (l'âge), « il » reprend sans rien ajouter.
\end{exemplebox}

\courssub{Les substituts du référent}

Pour éviter de répéter sans cesse le même nom, la langue dispose de \cle{substituts} : des mots ou groupes de mots qui \emph{reprennent} le référent. On distingue quatre grands types de reprises.

\begin{aretenir}[Les quatre types de substituts]
\begin{enumerate}
\item \textbf{Les pronoms personnels} : il, elle, ils, le, lui, en, y... Exemple : « Meka attend. \emph{Il} est fatigué. »
\item \textbf{Les pronoms démonstratifs} : celui, celle, ceux, cela, ce... Exemple : « Les invités arrivèrent ; \emph{celui} qui ouvrait la marche était le commandant. »
\item \textbf{Les reprises nominales} : un nom commun (souvent précédé de \emph{le, ce, cet}) remplace le nom propre : « Meka » devient « \emph{le vieillard} », « \emph{le paysan} », « \emph{ce père de famille} ».
\item \textbf{Les périphrases} : une expression de plusieurs mots qui désigne le référent par un détail : « \emph{le héros d'Ekoumdoum} », « \emph{l'ami du commandant d'un jour} ».
\end{enumerate}
\end{aretenir}

L'ensemble des reprises qui désignent un même référent dans un texte forme une \cle{chaîne de reprise} (ou chaîne anaphorique). Suivre cette chaîne, c'est vérifier la \cle{cohérence référentielle} du texte : le lecteur doit toujours pouvoir répondre à la question « de qui, de quoi parle-t-on ? ».

\begin{popfigure}
\begin{center}
\begin{tikzpicture}[
  box/.style={draw=popblue, thick, rounded corners, fill=popblueL, minimum width=2.6cm, minimum height=0.9cm, align=center, font=\small},
  ref/.style={draw=poporange, very thick, rounded corners, fill=poporangeL, minimum width=3.2cm, minimum height=1.1cm, align=center, font=\small\bfseries},
  fl/.style={-{Stealth[length=3mm]}, thick, popdark}]
\node[ref, align=center] (R) at (0,0){RÉFÉRENT\\ Meka (le personnage)};
\node[box, align=center] (A) at (-5.2,2.2){« Meka »\\ (nom propre)};
\node[box, align=center] (B) at (-1.8,2.6){« le vieillard »\\ (reprise nominale)};
\node[box, align=center] (C) at (1.8,2.6){« il »\\ (pronom personnel)};
\node[box, align=center] (D) at (5.2,2.2){« le vieux nègre »\\ (périphrase connotée)};
\draw[fl] (A.south) -- (R.north west);
\draw[fl] (B.south) -- (R.north);
\draw[fl] (C.south) -- (R.north);
\draw[fl] (D.south) -- (R.north east);
\draw[popmuted, dashed, thick] (A.east) -- (B.west);
\draw[popmuted, dashed, thick] (B.east) -- (C.west);
\draw[popmuted, dashed, thick] (C.east) -- (D.west);
\end{tikzpicture}
\end{center}
\legende{Chaîne de reprise dans \emph{Le Vieux nègre et la médaille} : quatre expressions différentes (flèches pleines) renvoient à un seul référent, Meka. Les pointillés figurent l'enchaînement des reprises dans le texte.}
\end{popfigure}

\begin{exemplebox}[Exemple chiffré tiré du roman]
Dans un extrait de 15 lignes du chapitre consacré à l'annonce de la médaille, le référent « Meka » apparaît 11 fois : \textbf{2 fois} sous la forme du nom propre « Meka » et \textbf{9 fois} par des substituts variés (5 pronoms personnels « il », 3 reprises nominales comme « le vieillard », 1 périphrase). Soit environ 82 \% de reprises ($9/11 \approx 0{,}82$) : sans ces substituts, le texte serait une litanie monotone de « Meka... Meka... Meka... ».
\end{exemplebox}

\courssub{La synonymie : la proximité de sens}

Observons maintenant le champ lexical de la récompense dans le roman : \emph{médaille}, \emph{récompense}, \emph{distinction}, \emph{honneur}, \emph{décoration}. Ces mots ne sont pas identiques, mais ils se ressemblent : ils entretiennent une relation de \cle{synonymie}.

\begin{definition}[Synonymie et antonymie]
La \cle{synonymie} est la relation de \textbf{proximité de sens} entre deux unités lexicales : deux mots sont synonymes lorsqu'ils peuvent désigner la même réalité dans certains contextes (ex. \emph{récompense / distinction}). L'\cle{antonymie} est la relation d'\textbf{opposition de sens} entre deux unités lexicales (ex. \emph{dignité / humiliation}).
\end{definition}

Il faut insister sur une nuance essentielle : les synonymes sont presque toujours des synonymes \textbf{partiels}. Ils partagent un noyau de sens commun, mais chacun possède des nuances propres :

\begin{itemize}
\item la \textbf{médaille} est un objet matériel, concret, que l'on épingle sur la poitrine ;
\item la \textbf{distinction} souligne le fait d'être distingué, séparé du lot ;
\item la \textbf{récompense} insiste sur l'idée de paiement d'un mérite ;
\item l'\textbf{honneur} relève du registre moral et social, plus abstrait.
\end{itemize}

L'emploi d'un synonyme dépend donc de \textbf{conditions} : le contexte, le registre de langue (familier, courant, soutenu) et la nuance voulue. On dit « recevoir une médaille », mais on ne dit pas « épingler une récompense sur la poitrine ».

\courssub{L'antonymie : l'opposition de sens}

L'antonymie est le moteur de l'opposition argumentative. Dans le roman, tout le drame de Meka tient dans une série d'oppositions : \emph{dignité / humiliation}, \emph{honorer / déshonorer}, \emph{fierté / honte}, \emph{espoir / désillusion}.

La langue fabrique aussi des antonymes par \cle{préfixation}, c'est-à-dire en ajoutant un préfixe de négation ou d'inversion :

\begin{itemize}
\item \textbf{il-} : légal $\rightarrow$ \emph{illégal} ; légitime $\rightarrow$ \emph{illégitime} ;
\item \textbf{in-} : justice $\rightarrow$ \emph{injustice} ; digne $\rightarrow$ \emph{indigne} ;
\item \textbf{im-} : possible $\rightarrow$ \emph{impossible} ; poli $\rightarrow$ \emph{impoli} ;
\item \textbf{dé-} (inversion de l'action) : honorer $\rightarrow$ \emph{déshonorer} ; faire $\rightarrow$ \emph{défaire}.
\end{itemize}

\begin{popfigure}
\begin{center}
\begin{tikzpicture}
\node[draw=popdark, thick, fill=popcream, rounded corners, font=\small\bfseries] at (0,3.4) {Champ lexical de l'honneur dans le roman};
\node[draw=popgreen, thick, fill=popgreenL, rounded corners, minimum width=4.6cm, font=\small] (S) at (-3.2,2.2) {\textbf{Synonymes (proximité)}};
\node[draw=poppink, thick, fill=poppinkL, rounded corners, minimum width=4.6cm, font=\small] (A) at (3.2,2.2) {\textbf{Antonymes (opposition)}};
\node[draw=popline, fill=white, rounded corners, font=\small, align=left] at (-3.2,0.6) {honneur $\approx$ gloire $\approx$ fierté\\ récompense $\approx$ distinction\\ médaille $\approx$ décoration};
\node[draw=popline, fill=white, rounded corners, font=\small, align=left] at (3.2,0.6) {dignité $\neq$ humiliation\\ honorer $\neq$ déshonorer\\ fierté $\neq$ honte};
\draw[-{Stealth}, very thick, poppurple] (-1.2,1.4) -- (1.2,1.4) node[midway, above, font=\scriptsize\itshape, poppurple]{opposition};
\end{tikzpicture}
\end{center}
\legende{Tableau à double entrée du champ lexical de l'honneur : à gauche, les mots de sens proche ; à droite, les mots de sens contraire. La flèche centrale marque la relation d'opposition qui structure le récit (la médaille honore puis déshonore Meka).}
\end{popfigure}

\courssub{Rôle stylistique : pourquoi varier les mots ?}

Ces outils de langue ne sont pas des décorations : ils remplissent trois fonctions précises, que vous devrez exploiter dans vos dissertations.

\begin{enumerate}
\item \textbf{Éviter les répétitions} : répéter dix fois « Meka » ou « l'auteur » fatigue le lecteur et donne une impression de pauvreté vocabulaire.
\item \textbf{Nuancer} : chaque substitut ou synonyme ajoute une information. « Le vieillard » insiste sur l'âge et la faiblesse ; « le vieux nègre » fait entendre le mépris colonial. Choisir un mot, c'est choisir un point de vue.
\item \textbf{Créer des oppositions argumentatives} : l'antonymie construit les contrastes qui nourrissent une problématique (« la médaille, censée \emph{honorer} Meka, ne fait-elle pas en réalité le \emph{déshonorer} ? »).
\end{enumerate}

\begin{methode}[Éviter les répétitions dans une dissertation]
\begin{enumerate}
\item \textbf{Repérer} le nom qui revient trop souvent dans votre brouillon (soulignez-le à chaque occurrence).
\item \textbf{Alterner} les trois procédés : un pronom (« il », « celui-ci »), une reprise nominale (« le vieillard », « ce personnage »), un synonyme ou une périphrase (« le héros d'Ekoumdoum »).
\item \textbf{Vérifier} la nuance : le substitut choisi convient-il au contexte et au registre soutenu de la dissertation ?
\item \textbf{Relire} la phrase entière pour contrôler que le référent reste identifiable (pas de pronom ambigu renvoyant à deux personnes possibles).
\end{enumerate}
\end{methode}

\begin{attention}[Deux synonymes ne sont jamais parfaitement interchangeables]
Erreur fréquente : remplacer un mot par un synonyme sans vérifier le contexte. On ne peut pas écrire « Meka reçut une \emph{dignité} » à la place de « Meka reçut une \emph{médaille} » : la médaille est un objet concret, la dignité une valeur abstraite. Avant d'employer un synonyme, interrogez toujours trois critères : la \textbf{nuance de sens}, le \textbf{registre de langue} (on évitera « galette » pour « argent » dans une copie) et la \textbf{construction grammaticale} du mot. De même, un pronom mal rattaché (« il » pouvant désigner deux personnages) brouille la référence et fait perdre des points de cohérence.
\end{attention}

\begin{exoresolu}[Reformuler une phrase maladroite]
Énoncé : reformulez le passage suivant en supprimant les répétitions, grâce aux substituts du référent et aux synonymes : « \emph{Meka reçoit la médaille. Meka est fier de la médaille. Mais la médaille va causer la perte de Meka, car la médaille isole Meka de sa communauté.} »
\tcblower
Solution détaillée : on identifie d'abord le référent répété : « Meka » (4 occurrences) et « la médaille » (4 occurrences). On applique ensuite la méthode en alternant pronom, reprise nominale et synonyme. Reformulation possible : « \textbf{Meka} reçoit la \textbf{médaille}. \textbf{Le vieillard} est fier de \textbf{cette distinction}. Mais \textbf{celle-ci} va causer \textbf{sa} perte, car \textbf{cette récompense} l'isole de sa communauté. » On vérifie : le référent reste clair à chaque reprise ; les synonymes « distinction » et « récompense » conviennent au contexte (il s'agit bien d'un honneur reçu) ; le registre est soutenu, adapté à une dissertation. Bonus stylistique : l'opposition antonymique peut être soulignée en conclusion de phrase : « ce qui devait l'\emph{honorer} le \emph{déshonore} aux yeux des siens ».
\end{exoresolu}

\begin{exercice}[À vous de jouer]
\begin{enumerate}
\item Dans la phrase « Le commandant remet la croix à Meka ; celui-ci la contemple avec émotion », indiquez le référent de « celui-ci » et de « la ». Que se passerait-il si l'on remplaçait « celui-ci » par « il » ?
\item Classez les mots suivants en deux colonnes (synonymes de \emph{honorer} / antonymes de \emph{honorer}) : déshonorer, glorifier, humilier, célébrer, mépriser, distinguer.
\item Formez des antonymes par préfixation à partir de : \emph{légal, justice, honorer, possible, digne}.
\item Reformulez en évitant toute répétition : « La médaille est remise à Meka. La médaille rend Meka orgueilleux. À cause de la médaille, Meka perd l'amitié des villageois. »
\end{enumerate}
\end{exercice}

\begin{corrige}[Exercice : outils de langue I]
\begin{enumerate}
\item « Celui-ci » a pour référent \textbf{Meka} (le dernier nommé, le plus proche) ; « la » reprend « la croix ». Avec « il », la phrase deviendrait ambiguë : le pronom pourrait renvoyer au commandant comme à Meka — la cohérence référentielle serait rompue.
\item Synonymes de \emph{honorer} : glorifier, célébrer, distinguer. Antonymes : déshonorer, humilier, mépriser.
\item Légal $\rightarrow$ illégal ; justice $\rightarrow$ injustice ; honorer $\rightarrow$ déshonorer ; possible $\rightarrow$ impossible ; digne $\rightarrow$ indigne.
\item Reformulation attendue : « La \textbf{médaille} est remise à \textbf{Meka}. \textbf{Cette distinction} rend \textbf{le vieillard} orgueilleux. À cause d'\textbf{elle}, \textbf{il} perd l'amitié des villageois. » Toute autre alternance correcte (pronom, reprise nominale, synonyme) est acceptée, à condition que le référent demeure identifiable.
\end{enumerate}
\end{corrige}

\begin{aretenir}[L'essentiel de la section]
Le \cle{référent} est la réalité désignée par un énoncé ; ses \cle{substituts} (pronoms personnels et démonstratifs, reprises nominales, périphrases) forment des chaînes de reprise qui assurent la cohérence du texte. La \cle{synonymie} (proximité de sens, toujours partielle) et l'\cle{antonymie} (opposition de sens, souvent construite par préfixation : il-, in-, im-, dé-) permettent d'éviter les répétitions, de nuancer le propos et de bâtir des oppositions argumentatives. Ces outils seront indispensables pour rédiger le plan détaillé de votre dissertation : une problématique bien formulée repose souvent sur une opposition antonymique forte.
\end{aretenir}

\coursec{Outils de langue II : hypéronymie/hyponymie, contenus manifestes et latents, texte argumentatif}

Après avoir maîtrisé les mécanismes de la référence (référent et substituts) et les relations de sens proches ou opposés (synonymie, antonymie), nous disposons maintenant d'outils pour \textbf{nommer} et \textbf{varier} le discours. Il nous faut à présent comprendre comment les mots s'organisent en \textbf{réseaux hiérarchiques} (hypéronymie/hyponymie), comment un texte dit souvent \textbf{plus qu'il n'écrit} (manifeste/latent), et comment se construit un \textbf{texte argumentatif} efficace — compétences indispensables pour analyser \emph{Le Vieux nègre et la médaille} et réussir votre dissertation.

\courssub{L'hypéronymie et l'hyponymie : organiser le lexique en arborescence}

\begin{definition}[Relation d'inclusion sémantique]
L'\textbf{hypéronymie} et l'\textbf{hyponymie} décrivent une relation logique d'inclusion entre deux unités lexicales :
\begin{itemize}
    \item L'\textbf{hypéronyme} (ou terme générique, superordonné) est le mot au sens \textbf{plus large}, qui \textbf{englobe} l'autre. Ex. : \emph{récompense}, \emph{arme}, \emph{vêtement}.
    \item L'\textbf{hyponyme} (ou terme spécifique, subordonné) est le mot au sens \textbf{plus étroit}, \textbf{inclus} dans le premier. Ex. : \emph{médaille}, \emph{diplôme}, \emph{prime} sont hyponymes de \emph{récompense}.
\end{itemize}
On note : \textbf{Hyponyme $\subset$ Hypéronyme}. C'est une relation \textbf{transitive} : si \emph{médaille} $\subset$ \emph{récompense} et \emph{récompense} $\subset$ \emph{distinction}, alors \emph{médaille} $\subset$ \emph{distinction}.
\end{definition}

\begin{methode}[Identifier et exploiter une relation hypéronymique]
\begin{enumerate}
    \item \textbf{Tester l'inclusion} : dire « X est un type de Y » ou « Y, comme par exemple X ». Si la phrase est naturelle, X est hyponyme de Y.
    \item \textbf{Construire l'arbre} : placer l'hypéronyme au sommet, les hyponymes au niveau inférieur (co-hyponymes entre eux).
    \item \textbf{Exploiter en dissertation} : 
    \begin{itemize}
        \item \textbf{Définir} : partir de l'hypéronyme (genre) + différence spécifique $\rightarrow$ définition rigoureuse.
        \item \textbf{Classer} : regrouper des exemples sous une même catégorie pour structurer un paragraphe.
        \item \textbf{Exemplifier} : descendre du général au particulier pour illustrer une idée abstraite.
    \end{itemize}
\end{enumerate}
\end{methode}

\begin{exemplebox}[Arbre hypéronymique dans \emph{Le Vieux nègre et la médaille}]
Dans le roman, la \textbf{médaille} offerte à Meka n'est pas un objet isolé ; elle appartient à un système de \textbf{récompenses} coloniales. L'arbre ci-dessous montre comment le terme générique « récompense » structure plusieurs réalités concrètes (hyponymes), dont certaines apparaissent explicitement ou implicitement dans l'œuvre.
\end{exemplebox}

\begin{popfigure}
\begin{tikzpicture}[
    level distance=2.2cm,
    sibling distance=3.5cm,
    every node/.style={font=\small, align=center},
    edge from parent/.style={draw=popdark, thick, -{Stealth[length=3mm]}},
    hyp/.style={fill=popblue, text=white, rounded corners=4pt, inner sep=6pt, font=\bfseries\small},
    hypo/.style={fill=popblueL, draw=popblue, rounded corners=3pt, inner sep=4pt},
    ex/.style={fill=popcream, draw=poporange, rounded corners=2pt, inner sep=3pt, font=\scriptsize\itshape}
]
\node[hyp, align=center]{Récompense\\ (Hypéronyme)}
    child { node[hypo] {Médaille}
        child { node[ex, align=center]{Médaille d'honneur\\ (Meka, ch. 3)} }
        child { node[ex] {Médaille du travail} }
    }
    child { node[hypo] {Diplôme / Certificat}
        child { node[ex] {Certificat d'études} }
        child { node[ex] {Diplôme de chef} }
    }
    child { node[hypo] {Prime / Argent}
        child { node[ex] {Prime de fidélité} }
        child { node[ex] {Indemnité de départ} }
    }
    child { node[hypo] {Citation / Titre honorifique}
        child { node[ex] {Chef de canton} }
        child { node[ex] {Médaille de la Reconnaissance} }
    };
\end{tikzpicture}
\legende{Arbre hypéronymique : la médaille de Meka comme hyponyme de « récompense », parmi d'autres distinctions coloniales.}
\end{popfigure}

\begin{attention}[Piège fréquent : confusion niveau de langue / niveau sémantique]
Ne confondez pas \textbf{registre de langue} (familier, courant, soutenu) et \textbf{niveau sémantique} (générique vs spécifique). « Bagnole » (familier) et « automobile » (soutenu) sont \textbf{synonymes} (même niveau sémantique), pas hypéronyme/hyponyme. En revanche, « véhicule » (hypéronyme) $\supset$ « automobile » (hyponyme).
\end{attention}

\begin{exoresolu}[Analyser une définition par genre et différence spécifique]
\textbf{Énoncé :} Relevez l'hypéronyme et la différence spécifique dans cette phrase (exemple construit d'après le roman) : « La médaille, \emph{cette distinction honorifique réservée aux serviteurs fidèles de l'Administration}, brillait sur la poitrine de Meka. »

\textbf{Solution détaillée :}
\begin{itemize}
    \item \textbf{Hypéronyme (genre)} : « distinction honorifique » (terme plus large que médaille).
    \item \textbf{Différence spécifique} : « réservée aux serviteurs fidèles de l'Administration » (critère qui singularise cette médaille parmi les autres distinctions).
    \item \textbf{Intérêt} : L'auteur utilise une définition formelle (genre + différence) pour installer l'ironie : le lecteur sait que Meka n'est pas « fidèle » par conviction mais par contrainte, ce qui crée un \textbf{décalage latent} (voir partie suivante).
\end{itemize}
\end{exoresolu}

\courssub{Contenus manifestes et contenus latents : lire entre les lignes}

\begin{definition}[Double niveau de signification]
Tout énoncé littéraire (et tout discours) porte deux strates de sens :
\begin{itemize}
    \item Les \textbf{contenus manifestes} (ou \textbf{explicites}) : ce que le texte \textbf{dit ouvertement}, l'information littérale, le « discours de surface ». Ex. : « Le commandant remet la médaille à Meka. »
    \item Les \textbf{contenus latents} (ou \textbf{implicites}) : ce que le texte \textbf{suggère, sous-entend, tait}, le « discours de profondeur » accessible par l'inférence du lecteur. Ex. : « Le commandant humilie Meka en lui donnant une breloque pour une vie de servitude. »
\end{itemize}
Le passage du manifeste au latent relève de l'\textbf{interprétation} : le lecteur \textbf{actif} comble les blancs, décrypte les non-dits.
\end{definition}

\begin{methode}[Décoder le latent : démarche en 4 étapes]
\begin{enumerate}
    \item \textbf{Repérer le manifeste} : citer la phrase exacte, décrire l'action/l'énoncé littéral.
    \item \textbf{Identifier les indices de décalage} : 
    \begin{itemize}
        \item \textbf{Ironie} : écart entre ce qui est dit et ce qui est pensé (souvent marqueurs : « bien sûr », « naturellement », guillemets ironiques, hyperboles).
        \item \textbf{Connotations} : charge affective/culturelle des mots (ex. « médaille » $\rightarrow$ connotation honneur/valeur ; mais contexte $\rightarrow$ connotation mépris/dérision).
        \item \textbf{Non-dits / Silences} : ce qui n'est pas dit mais attendu (ex. Meka ne remercie pas, le texte ne donne pas son discours intérieur).
        \item \textbf{Contexte énonciatif} : qui parle ? à qui ? dans quelle intention ? (ici : narrateur omniscient ironique, focalisation sur Meka).
    \end{itemize}
    \item \textbf{Formuler l'implicite} : écrire une phrase affirmative qui explicite le sens caché (hypothèse de lecture).
    \item \textbf{Vérifier la cohérence} : l'implicite doit s'accorder avec l'ensemble de l'œuvre (thèmes, ton, finale).
\end{enumerate}
\end{methode}

\begin{exemplebox}[Manifeste vs Latent : la scène de la remise (ch. 3-4)]
\begin{center}
\begin{tabularx}{\linewidth}{|l|X|X|}
\hline
\rowcolor{popblueL}
\textbf{Extrait / Situation} & \textbf{Contenu Manifeste (Explicite)} & \textbf{Contenu Latent (Implicite)} \\ \hline
Le commandant épingle la médaille. & Cérémonie officielle, protocole respecté, reconnaissance formelle. & Geste condescendant ; le commandant « joue » son rôle, il ne respecte pas Meka. \\ \hline
Meka reçoit la médaille « la tête haute ». & Attitude digne, fierté apparente. & Meka subit ; « la tête haute » masque la honte et la lucidité naissante. \\ \hline
Le discours : « \emph{Pour services rendus à la France} ». & Formule administrative standard, justification officielle. & Ironie cruelle : quels « services » ? La soumission, la perte d'identité, la trahison de ses pairs. \\ \hline
La foule applaudit. & Approbation collective, joie partagée. & Applaudissements forcés, peur de l'autorité, incompréhension du drame intime. \\ \hline
\end{tabularx}
\end{center}
\end{exemplebox}

\begin{popfigure}
\begin{tikzpicture}[
    scale=0.9,
    every node/.style={font=\small},
    layer/.style={draw=popdark, thick, rounded corners=6pt, inner sep=8pt, minimum width=12cm},
    arrow/.style={->, >=Stealth, thick, poporange, shorten <=2pt, shorten >=2pt},
    labelstyle/.style={font=\bfseries\small, text=popdark}
]
% Couche Manifeste (Surface)
\node[layer, fill=popblue!10, label=above left:\textcolor{popblue}{\textbf{SURFACE : Contenus Manifestes}}, align=center] (manifeste) at (0, 2){
    \begin{tabular}{@{}p{5.5cm}@{}}
    \textbf{Énoncés littéraux :} « Le commandant remet la médaille à Meka pour services rendus. » \\
    \textbf{Actions visibles :} Cérémonie, discours, applaudissements, médaille qui brille. \\
    \textbf{Registre :} Administratif, solennel, positif.
    \end{tabular}
};

% Couche Latent (Profondeur)
\node[layer, fill=poporange!10, label=above left:\textcolor{poporange}{\textbf{PROFONDEUR : Contenus Latents}}, align=center] (latent) at (0, -2){
    \begin{tabular}{@{}p{5.5cm}@{}}
    \textbf{Sens implicites :} Humiliation déguisée, récompense de la soumission, ironie du sort. \\
    \textbf{Non-dits :} Colère rentrée de Meka, mépris du colon, vide de la distinction. \\
    \textbf{Registre :} Tragique, critique, dérisoire.
    \end{tabular}
};

% Flèches d'interprétation
\draw[arrow] (manifeste.south) -- node[right, labelstyle, align=center]{Inférence \\ Lecteur actif} (latent.north);
\draw[arrow, dashed, popgreen] (latent.east) to[bend right=20] node[right, labelstyle, align=center]{Retour au texte \\ Validation} (manifeste.east);
\end{tikzpicture}
\legende{Schéma à deux couches : le lecteur actif infère le latent à partir des indices du manifeste, puis valide son hypothèse par un retour au texte.}
\end{popfigure}

\begin{savaistu}[La médaille : symbole ironique au cœur du roman]
La médaille offerte à Meka condense tout le sens latent de l'œuvre de Ferdinand Oyono. \textbf{Manifeste} : c'est une décoration honorifique, un objet métallique brillant, symbole de mérite. \textbf{Latent} : elle devient le \textbf{symbole de l'aliénation}. Elle récompense non le mérite, mais \textbf{la servilité} ; non l'homme, mais \textbf{l'instrument} docile. Le renversement final (Meka jette la médaille dans les latrines) est la \textbf{révélation violente du latent} : ce que le texte suggérait tout au long devient manifeste par l'acte ultime de refus. Tout le roman tient dans ce basculement du manifeste au latent.
\end{savaistu}

\begin{attention}[Erreur de lecture : prendre l'ironie au premier degré]
L'erreur majeure en Terminale est de \textbf{confondre le point de vue du personnage (ou du discours officiel) et celui de l'auteur/narrateur}. Si vous écrivez : « Meka est fier de sa médaille car le texte dit qu'il la reçoit la tête haute », vous restez au \textbf{manifeste} et ratez l'\textbf{ironie} d'Oyono. Toujours demander : \emph{« Qui parle ? Quelle est la distance ironique du narrateur ? »}.
\end{attention}

\courssub{Le texte argumentatif : structure, stratégies et connecteurs}

\begin{definition}[Texte argumentatif]
Un \textbf{texte argumentatif} a pour fonction principale de \textbf{convaincre} (visée \textbf{persuasive} : toucher l'affect, l'adhésion) ou de \textbf{persuader} (visée \textbf{délibérative} : amener à agir, à décider) un destinataire. Il repose sur une architecture rigoureuse :
\begin{itemize}
    \item Une \textbf{thèse} (position défendue, réponse à la problématique).
    \item Des \textbf{arguments} (raisons logiques qui soutiennent la thèse), souvent hiérarchisés (du plus fort au plus faible, ou progression dialectique).
    \item Des \textbf{exemples / illustrations} (faits précis, citations, statistiques, anecdotes) qui \textbf{concrétisent} chaque argument.
    \item Des \textbf{connecteurs logiques} (articulateurs) qui assurent la cohérence et guident le raisonnement.
\end{itemize}
\end{definition}

\begin{propriete}[Structure canonique du paragraphe argumentatif -- Méthodologie exigible au Bac]
Tout paragraphe argumentatif bien formé suit le schéma \textbf{A.E.C.} (ou \textbf{A.E.I.}) :
\begin{center}
\textbf{Argument (Idée directrice)} $\rightarrow$ \textbf{Explication / Développement} $\rightarrow$ \textbf{Exemple / Illustration}
\end{center}
\begin{itemize}
    \item L'\textbf{argument} est une \textbf{affirmation générale} (ex. « La colonisation déshumanise l'indigène »).
    \item L'\textbf{exemple} est un \textbf{cas particulier} (ex. « Meka, après 30 ans de service, reçoit une médaille dérisoire »).
    \item \textbf{Règle d'or} : \textbf{L'exemple ne remplace jamais l'argument}. Il le \textbf{prouve}, l'\textbf{ancre} dans le réel.
\end{itemize}
\end{propriete}

\begin{methode}[Identifier les stratégies argumentatives (inventaire Bac)]
\begin{enumerate}
    \item \textbf{Argument d'autorité} : s'appuyer sur un expert, un texte sacré, une tradition. Ex. « Comme l'écrit Aimé Césaire... »
    \item \textbf{Argument de l'expérience / du fait} : invoquer un fait observable, une statistique, un témoignage. Ex. « Dans le roman, Meka perd sa terre... »
    \item \textbf{Argument \emph{ad absurdum} (par l'absurde)} : montrer que la thèse adverse mène à une contradiction. Ex. « Si la médaille honorait vraiment Meka, il ne la jetterait pas. »
    \item \textbf{La concession (ou argument concédant)} : admettre un point à l'adversaire pour mieux le réfuter ensuite. Ex. « Certes, Meka a reçu une médaille (\emph{concession}), mais cette médaille est une insulte (\emph{réfutation}). »
    \item \textbf{Argument analogique} : comparer deux situations similaires. Ex. « Comme le vieux nègre, tout colonisé décoré est en réalité dépossédé. »
\end{enumerate}
\end{methode}

\begin{popfigure}
\begin{tikzpicture}[
    scale=0.85,
    every node/.style={font=\small, align=center},
    thesis/.style={fill=poppurple, text=white, rounded corners=6pt, inner sep=8pt, font=\bfseries},
    arg/.style={fill=poppurpleL, draw=poppurple, rounded corners=4pt, inner sep=6pt},
    ex/.style={fill=popcream, draw=poporange, rounded corners=3pt, inner sep=4pt, font=\scriptsize},
    conn/.style={fill=popgreenL, draw=popgreen, rounded corners=3pt, inner sep=3pt, font=\scriptsize\itshape},
    arrow/.style={->, >=Stealth, thick, popdark},
    doublearrow/.style={<->, >=Stealth, dashed, popmuted}
]
% Thèse
\node[thesis] (thesis) at (0, 5) {THÈSE : \emph{La médaille coloniale est un outil d'aliénation, non de reconnaissance.}};

% Argument 1
\node[arg] (arg1) at (-5, 2.5) {ARGUMENT 1 : \emph{Inversion des valeurs : l'honneur masque le mépris.}};
\node[ex] (ex1a) at (-7, 0.5) {Ex. : Discours du commandant : « Services rendus à la France »};
\node[ex] (ex1b) at (-3, 0.5) {Ex. : Médaille « brillante » sur « poitrine noire » (contraste visuel)};

% Argument 2
\node[arg] (arg2) at (5, 2.5) {ARGUMENT 2 : \emph{La récompense crée une dépendance psychologique.}};
\node[ex] (ex2a) at (3, 0.5) {Ex. : Meka « fier » puis « honteux » (évolution ch. 3 $\rightarrow$ 5)};
\node[ex] (ex2b) at (7, 0.5) {Ex. : Kelara (femme) pousse Meka à l'accepter (intériorisation)};

% Conclusion
\node[thesis, fill=popblue] (concl) at (0, -2) {CONCLUSION : \emph{Le rejet final (latrines) = reprise de dignité / rupture.}};

% Connecteurs
\node[conn] (c1) at (-5, 3.8) {D'abord / En premier lieu};
\node[conn] (c2) at (5, 3.8) {Ensuite / Par ailleurs};
\node[conn] (c3) at (0, 1.2) {Enfin / Surtout};
\node[conn] (c4) at (0, -0.5) {Ainsi / En définitive};

% Flèches structurelles
\draw[arrow] (thesis.south) -- (c1.north);
\draw[arrow] (thesis.south) -- (c2.north);
\draw[arrow] (c1.south) -- (arg1.north);
\draw[arrow] (c2.south) -- (arg2.north);
\draw[arrow] (arg1.south) -- (ex1a.north);
\draw[arrow] (arg1.south) -- (ex1b.north);
\draw[arrow] (arg2.south) -- (ex2a.north);
\draw[arrow] (arg2.south) -- (ex2b.north);
\draw[arrow] (arg1.south east) -- (c3.north west);
\draw[arrow] (arg2.south west) -- (c3.north east);
\draw[arrow] (c3.south) -- (c4.north);
\draw[arrow] (c4.south) -- (concl.north);

% Liens exemples-arguments (illustration)
\draw[doublearrow] (ex1a.east) -- (ex1b.west);
\draw[doublearrow] (ex2a.east) -- (ex2b.west);
\end{tikzpicture}
\legende{Architecture d'un texte argumentatif : thèse, arguments hiérarchisés (connecteurs), exemples illustrant chaque argument, conclusion.}
\end{popfigure}

\begin{exemplebox}[Analyse d'un paragraphe argumentatif (extrait type dissertation)]
\textit{« \textbf{En premier lieu}, la médaille offerte à Meka n'honore pas son mérite, elle sanctionne sa soumission. \textbf{En effet}, le texte montre que le commandant la remet « pour services rendus à la France » (p. 45), formule administrative vide qui occulte la réalité : Meka a servi l'oppresseur contre son peuple. \textbf{Par exemple}, la scène où Meka découvre que sa médaille ne lui assure aucun privilège concret (pas de terre, pas de respect) prouve que l'objet est symboliquement nul. \textbf{Ainsi}, cette distinction est un leurre. »}

\begin{itemize}
    \item \textbf{Connecteurs} : \emph{En premier lieu} (énumération), \emph{En effet} (explication), \emph{Par exemple} (illustration), \emph{Ainsi} (conséquence/conclusion partielle).
    \item \textbf{Structure A.E.I.} : Argument (phrase 1) $\rightarrow$ Explication + référence texte (phrase 2) $\rightarrow$ Exemple concret (phrase 3) $\rightarrow$ Mini-conclusion (phrase 4).
\end{itemize}
\end{exemplebox}

\begin{exoresolu}[Distinguer argument et exemple dans un sujet de dissertation]
\textbf{Sujet :} « La récompense est-elle toujours une marque d'estime ? » (Sujet type Bac).

\textbf{Proposition d'élèves :}
\begin{enumerate}
    \item \textit{Argument :} « Meka reçoit une médaille. »
    \item \textit{Argument :} « Les médailles sont en métal. »
    \item \textit{Argument :} « Dans le roman, Meka jette la médaille aux latrines. »
\end{enumerate}

\textbf{Correction et analyse :}
\begin{itemize}
    \item \textbf{1 et 3 sont des EXEMPLES} (faits précis, événements du roman). Ils ne sont pas des arguments car ils n'expriment pas une \textbf{raison générale}.
    \item \textbf{2 est une propriété physique} (non pertinente ici).
    \item \textbf{Vrais arguments possibles :}
    \begin{itemize}
        \item \textbf{Arg. 1 (Valeur symbolique) :} « La récompense perd sa valeur d'estime quand elle est décernée par un pouvoir illégitime. » $\rightarrow$ \emph{Exemple :} Médaille de Meka.
        \item \textbf{Arg. 2 (Intention du donneur) :} « L'estime suppose une réciprocité ; la médaille coloniale est unilatérale et condescendante. » $\rightarrow$ \emph{Exemple :} Discours du commandant vs silence de Meka.
        \item \textbf{Arg. 3 (Conséquence sur le receveur) :} « Une fausse récompense provoque la honte puis la révolte. » $\rightarrow$ \emph{Exemple :} Jeter la médaille aux latrines (ch. final).
    \end{itemize}
\end{itemize}
\end{exoresolu}

\begin{aretenir}[Synthèse des 3 outils pour la dissertation]
\begin{center}
\begin{tabularx}{\linewidth}{|l|X|X|}
\hline
\rowcolor{popblueL}
\textbf{Outil} & \textbf{Apport pour la DISSERTATION} & \textbf{Application à \emph{Le Vieux nègre et la médaille}} \\ \hline
\textbf{Hypéronymie / Hyponymie} & \textbf{Définir, Classer, Exemplifier} avec rigueur lexicale. & Définir « récompense » (hypéronyme) $\rightarrow$ classer médaille/diplôme/prime (hyponymes) $\rightarrow$ montrer que la médaille est un \emph{mauvais} hyponyme (ironie). \\ \hline
\textbf{Manifeste / Latent} & \textbf{Problématiser} (écart dit/non-dit), \textbf{Analyser} la portée ironique, \textbf{Nuancer} la thèse. & Thèse manifeste : « Meka est honoré. » Thèse latente : « Meka est humilié. » La dissertation doit articuler les deux niveaux. \\ \hline
\textbf{Texte argumentatif} & \textbf{Structurer} le développement (Thèse / Arguments / Exemples / Connecteurs), \textbf{Convaincre} le correcteur. & Votre devoir \emph{est} un texte argumentatif : thèse = votre réponse au sujet ; arguments = vos idées organisées ; exemples = citations du roman / faits précis. \\ \hline
\end{tabularx}
\end{center}
\end{aretenir}

\begin{exercice}[Entraînement : Manifeste / Latent \& Argumentation]
\begin{enumerate}
    \item Relevez dans le chapitre 1 (départ de Meka pour la cérémonie) \textbf{deux contenus manifestes} et formulez pour chacun \textbf{le contenu latent correspondant}.
    \item Pour le sujet « L'obéissance est-elle une vertu ? », proposez \textbf{trois arguments distincts} (généraux) et, pour chacun, \textbf{un exemple précis} tiré du roman.
    \item Reliez les arguments ci-dessous par les \textbf{connecteurs logiques} adaptés : 
    \begin{itemize}
        \item[Arg. A :] La médaille aliène Meka.
        \item[Arg. B :] Elle lui fait perdre l'estime des siens.
        \item[Arg. C :] Elle provoque sa révolte finale.
    \end{itemize}
\end{enumerate}
\end{exercice}

\begin{corrige}[Exercice 1 -- Éléments de réponse]
\begin{enumerate}
    \item \textbf{Ex. 1 :} Manifeste : « Meka endosse son meilleur pagne. » / Latent : « Meka se prépare pour un piège, il joue un rôle qu'il ne maîtrise pas. » \textbf{Ex. 2 :} Manifeste : « Kelara l'admire. » / Latent : « Kelara est aveuglée par la valeur d'échange coloniale (honneur apparent), elle ne voit pas la perte de substance. »
    \item \textbf{Arg. 1 (Vertu civique) :} L'obéissance à des lois justes fonde le lien social. \textit{Ex. :} Meka obéit aux ordres de plantation (travail) $\rightarrow$ cohésion villageoise (au début). \textbf{Arg. 2 (Aliénation) :} L'obéissance aveugle à un pouvoir injuste détruit l'autonomie. \textit{Ex. :} Meka accepte le poste de chef-canton $\rightarrow$ devient auxiliaire de l'oppression. \textbf{Arg. 3 (Révolte) :} La désobéissance finale restaure la dignité. \textit{Ex. :} Jeter la médaille = refus d'obéir au script colonial.
    \item \textbf{Connecteurs :} \emph{D'abord / En premier lieu}, la médaille aliène Meka. \emph{Par conséquent / Du fait que}, elle lui fait perdre l'estime des siens. \emph{Finalement / C'est pourquoi}, elle provoque sa révolte finale.
\end{enumerate}
\end{corrige}

\coursec{Analyser le sujet et formuler la problématique}

Dans la section précédente, nous avons appris à décoder les textes argumentatifs : repérer les hypéronymes et les hyponymes qui organisent le vocabulaire, distinguer les contenus manifestes (ce qui est dit explicitement) des contenus latents (ce qui est suggéré), et reconnaître les caractéristiques du texte argumentatif. Ces outils de lecture vont maintenant nous servir à l'attaque du devoir lui-même. Car avant d'écrire une seule ligne de dissertation, il faut accomplir un travail invisible mais décisif : \cle{analyser le sujet} et en tirer une \cle{problématique}. C'est l'acte de naissance du devoir. Une problématique solide annonce une copie solide ; une problématique bâclée condamne la copie, aussi belle soit l'écriture.

\courssub{Les trois types de sujets de dissertation}

Au Probatoire et au Baccalauréat technique, les sujets de dissertation littéraire se rangent en trois grandes familles. Les reconnaître immédiatement permet d'adopter la bonne attitude intellectuelle.

\begin{definition}[Les trois types de sujets]
\begin{itemize}
\item Le \cle{sujet de réflexion} : une \textbf{question ouverte} posée directement au candidat. Il invite à explorer un problème et à construire une réponse personnelle. Exemple : « Le roman doit-il se contenter de divertir le lecteur ? »
\item Le \cle{sujet de discussion} : une \textbf{citation} d'un auteur, d'un critique ou d'un personnage, que le candidat est invité à \emph{discuter}, c'est-à-dire à examiner en pesant le pour et le contre. Exemple : « "Le romancier est un témoin de son temps." Discutez cette affirmation. »
\item Le \cle{sujet de commentaire} : une \textbf{affirmation} ou une appréciation portée sur une œuvre, un genre ou un phénomène littéraire, que le candidat doit \emph{apprécier}, c'est-à-dire vérifier, nuancer et justifier. Exemple : « Vous apprécierez le réalisme de la peinture de la société coloniale dans \emph{Le Vieux nègre et la médaille} de Ferdinand Oyono. »
\end{itemize}
\end{definition}

\begin{attention}[Ne pas confondre les consignes]
« Discutez » exige une confrontation d'arguments (thèse, antithèse, dépassement). « Appréciez » exige un jugement argumenté et nuancé. « Commentez » exige une explication approfondie. Employer la démarche d'un type de sujet pour un autre type est une faute de méthode sanctionnée au Bac.
\end{attention}

\begin{popfigure}
\begin{tikzpicture}[font=\small]
\draw[fill=popblueL, draw=popblue, thick, rounded corners=2pt] (0,4.6) rectangle (3.9,6.4);
\node[align=center, text width=3.6cm] at (1.95,5.5) {\textbf{Sujet de réflexion}\\ question ouverte\\ \emph{« Le roman doit-il instruire ? »}};
\draw[fill=poporangeL, draw=poporange, thick, rounded corners=2pt] (4.3,4.6) rectangle (8.2,6.4);
\node[align=center, text width=3.6cm] at (6.25,5.5) {\textbf{Sujet de discussion}\\ citation à discuter\\ \emph{« "Le roman est un miroir." Discutez. »}};
\draw[fill=popgreenL, draw=popgreen, thick, rounded corners=2pt] (8.6,4.6) rectangle (12.5,6.4);
\node[align=center, text width=3.6cm] at (10.55,5.5) {\textbf{Sujet de commentaire}\\ affirmation à apprécier\\ \emph{« Appréciez le comique dans l'œuvre d'Oyono. »}};
\draw[fill=popcream, draw=popdark, thick, rounded corners=2pt] (3.2,2.6) rectangle (9.3,3.8);
\node[align=center, text width=5.6cm] at (6.25,3.2) {\textbf{Attitude commune : analyser,\\ problématiser, argumenter}};
\draw[-{Stealth}, thick, popblue] (1.95,4.6) -- (4.6,3.8);
\draw[-{Stealth}, thick, poporange] (6.25,4.6) -- (6.25,3.8);
\draw[-{Stealth}, thick, popgreen] (10.55,4.6) -- (7.9,3.8);
\end{tikzpicture}
\legende{Les trois types de sujets de dissertation littéraire, avec un exemple pour chacun. Quel que soit le type, la démarche d'analyse est la même.}
\end{popfigure}

\courssub{La méthode d'analyse du sujet en quatre étapes}

Analyser un sujet, c'est le démonter comme un mécanicien démonte un moteur : on identifie chaque pièce, on comprend son rôle, puis on repère ce qui fait tourner l'ensemble. Cette analyse demande \textbf{au minimum dix à quinze minutes} au brouillon. C'est un investissement rentable : tout le reste du devoir en dépend.

\begin{methode}[Analyser un sujet en 4 étapes]
\begin{enumerate}
\item \textbf{Lire et souligner les mots-clés.} Lire le sujet au moins trois fois, à voix basse si possible. Souligner les termes porteurs de sens : noms, verbes, adjectifs qui constituent le champ lexical principal du sujet. Ne garder que l'essentiel : en général deux à quatre mots-clés.
\item \textbf{Définir chaque mot-clé.} Pour chaque terme, noter au brouillon : son sens propre (définition de dictionnaire), son sens précis dans le sujet (le contexte le restreint souvent), ses synonymes (pour éviter les répétitions et ouvrir la réflexion) et ses antonymes (qui dessinent par contraste les limites du concept). Les notions de synonymie et d'antonymie étudiées en classe trouvent ici leur application directe.
\item \textbf{Repérer les limites du sujet.} Déterminer ce que le sujet \emph{inclut} (le domaine à traiter : un genre, une époque, une œuvre) et ce qu'il \emph{exclut} (ce qui relèverait du hors-sujet). Les mots de restriction (« seulement », « surtout », « toujours », le nom d'une œuvre précise) sont des bornes à respecter.
\item \textbf{Dégager la tension ou le paradoxe.} Chercher le « problème » : l'opposition, la contradiction apparente ou la difficulté que le sujet met en scène. Un bon sujet de dissertation contient toujours une tension entre deux valeurs, deux exigences ou deux points de vue. C'est cette tension qui donnera naissance à la problématique.
\end{enumerate}
\end{methode}

\begin{popfigure}
\begin{tikzpicture}[font=\small,
etape/.style={draw=popblue, thick, fill=popblueL, rounded corners=3pt, text width=3.1cm, align=center, minimum height=1.5cm},
fleche/.style={-{Stealth}, very thick, poporange}]
\node[etape, align=center] (e1) at (0,0){\textbf{Étape 1}\\ Souligner les mots-clés};
\node[etape, align=center] (e2) at (4.1,0){\textbf{Étape 2}\\ Définir chaque mot-clé};
\node[etape, align=center] (e3) at (8.2,0){\textbf{Étape 3}\\ Repérer les limites};
\node[etape, align=center] (e4) at (12.3,0){\textbf{Étape 4}\\ Dégager la tension};
\node[draw=popdark, very thick, fill=popgoldL, rounded corners=3pt, text width=4.6cm, align=center, minimum height=1.3cm] (pb) at (6.15,-2.6) {\textbf{PROBLÉMATIQUE}\\ la question qui engage tout le devoir};
\draw[fleche] (e1) -- (e2);
\draw[fleche] (e2) -- (e3);
\draw[fleche] (e3) -- (e4);
\draw[fleche] (e1.south) .. controls (2,-1.6) .. (pb.west);
\draw[fleche] (e4.south) .. controls (10.3,-1.6) .. (pb.east);
\draw[fleche] (e2.south) -- (5.4,-1.95);
\draw[fleche] (e3.south) -- (6.9,-1.95);
\end{tikzpicture}
\legende{Organigramme de l'analyse du sujet : les quatre étapes successives convergent vers la formulation de la problématique.}
\end{popfigure}

Détaillons chaque étape avec l'exemple fil conducteur de cette leçon, en lien avec notre lecture de \emph{Le Vieux nègre et la médaille} : « La littérature engagée est-elle efficace pour dénoncer les injustices ? »

\textbf{Étape 1 : les mots-clés.} Dans ce sujet, on souligne : \emph{littérature engagée}, \emph{efficace}, \emph{dénoncer}, \emph{injustices}. Quatre termes qui portent tout le sens.

\textbf{Étape 2 : les définitions.} Au brouillon, on construit un petit tableau :

\begin{center}
\begin{tabularx}{\linewidth}{|l|X|X|}
\hline
\rowcolor{popblueL} \textbf{Mot-clé} & \textbf{Sens dans le sujet} & \textbf{Synonymes / antonymes} \\
\hline
Littérature engagée & œuvre qui prend parti, qui défend une cause (Sartre : « l'écrivain est en situation ») & syn. : littérature de combat, de témoignage ; ant. : art pour l'art \\
\hline
Efficace & qui produit l'effet attendu, qui atteint son but & syn. : opérant, utile ; ant. : vain, inefficace \\
\hline
Dénoncer & révéler publiquement, condamner en l'exposant & syn. : stigmatiser, condamner ; ant. : taire, encenser \\
\hline
Injustices & violations du droit et de l'équité (ici : colonialisme, hypocrisie, mépris) & syn. : abus, oppressions ; ant. : équité, justice \\
\hline
\end{tabularx}
\end{center}

\textbf{Étape 3 : les limites.} Le sujet porte sur la littérature \emph{engagée} (et non toute la littérature : les contes pour enfants ou la poésie amoureuse sont hors champ), et sur sa fonction de \emph{dénonciation} (et non sur ses qualités esthétiques). Parler du style d'Oyono sans lien avec la dénonciation serait un hors-sujet partiel.

\textbf{Étape 4 : la tension.} Le problème apparaît : d'un côté, la littérature semble une arme puissante (elle touche un large public, elle émeut, elle survit à son époque) ; de l'autre, un livre peut-il vraiment changer les choses ? Les mots ont-ils prise sur le pouvoir politique ? C'est précisément cette tension entre la \emph{force} et l'\emph{impuissance} supposée de la littérature qui fait problème.

\courssub{Formuler la problématique}

\begin{definition}[Problématique]
La \cle{problématique} est la \textbf{question centrale}, issue de l'analyse du sujet, à laquelle \emph{tout le devoir} va répondre. Elle reformule le sujet sous forme d'un véritable problème intellectuel, annonce l'angle d'attaque du candidat et commande le plan. Elle est en général introduite par « dans quelle mesure », « comment », « faut-il », « pourquoi ».
\end{definition}

La problématique n'est pas une simple répétition du sujet avec un point d'interrogation. Elle en est la \emph{mise en crise} : elle rend explicite la tension dégagée à l'étape 4. Pour notre exemple, on ne se contentera pas de recopier « La littérature engagée est-elle efficace pour dénoncer les injustices ? » ; on formulera par exemple : \emph{« Dans quelle mesure la littérature engagée, arme de mots face aux pouvoirs, parvient-elle réellement à dénoncer les injustices et à éveiller les consciences ? »} ou encore : \emph{« Comment l'écrivain engagé, comme Oyono, transforme-t-il le récit en instrument de dénonciation, et quelles sont les limites de cette efficacité ? »}

\begin{aretenir}[Les trois qualités d'une bonne problématique]
\begin{itemize}
\item Elle est \textbf{fidèle au sujet} : elle reprend ses mots-clés et respecte ses limites (pas de hors-sujet).
\item Elle est \textbf{précise} : elle n'est ni trop large (« la littérature est-elle utile ? » serait trop vague), ni déjà résolue (une question dont la réponse est évidente ne fait pas problème).
\item Elle \textbf{engage tout le devoir} : chaque partie du plan devra apporter un élément de réponse à cette question unique.
\end{itemize}
\end{aretenir}

\begin{exemplebox}[Exemple chiffré : « Faut-il récompenser la fidélité ? »]
Appliquons la méthode complète à ce sujet, qui fait écho à la médaille offerte à Meka pour sa « fidélité » dans \emph{Le Vieux nègre et la médaille}.
\begin{itemize}
\item \textbf{3 mots-clés} : \emph{faut-il} (interrogation sur la légitimité), \emph{récompenser} (donner un bien en échange d'un mérite), \emph{fidélité} (constance dans l'attachement et le service).
\item \textbf{1 tension} : la récompense suppose que la fidélité a une \emph{valeur marchande} qu'on peut payer ; mais la fidélité n'a-t-elle de prix que si elle est \emph{gratuite} et désintéressée ? Récompenser la fidélité, n'est-ce pas la détruire comme fidélité ? Dans le roman d'Oyono, la médaille donnée à Meka apparaît comme une récompense dérisoire et humiliante : elle masque l'exploitation au lieu de la payer.
\item \textbf{1 problématique} : \emph{« La récompense honore-t-elle la fidélité, ou risque-t-elle au contraire de la rabaisser en la transformant en marchandise ? »}
\end{itemize}
\end{exemplebox}

\begin{attention}[Erreur fréquente : la fausse problématique]
Beaucoup de candidats se contentent de \textbf{reformuler le sujet} en changeant l'ordre des mots : « Nous nous demanderons si la littérature engagée est efficace pour dénoncer les injustices. » Ce n'est pas une problématique, c'est une photocopie. La problématique doit poser un \emph{véritable problème} : elle fait apparaître la tension (ici : la force des mots contre l'impuissance des mots) et annonce un angle personnel.
\end{attention}

\begin{attention}[Erreur fréquente au Bac : le hors-sujet]
Le hors-sujet naît presque toujours d'une analyse bâclée en moins de dix minutes. Le candidat pressé lit le sujet une seule fois, retient un seul mot (« littérature »), et rédige tout ce qu'il sait sur la littérature en général. Résultat : une copie parfois bien écrite, mais à côté de la question, lourdement sanctionnée. Consacrez \textbf{quinze minutes} à l'analyse : c'est le meilleur placement de temps de toute l'épreuve.
\end{attention}

\begin{exoresolu}[Analyse complète d'un sujet de discussion]
Énoncé. Soit le sujet : « "Un roman africain ne vaut que s'il dénonce." Discutez cette affirmation. » Effectuez l'analyse du sujet et formulez la problématique.
\tcblower
Solution détaillée.
\textbf{Étape 1 — Mots-clés} : \emph{roman africain}, \emph{ne vaut que} (restriction forte), \emph{dénonce}.
\textbf{Étape 2 — Définitions} : le roman africain désigne la fiction narrative écrite par des auteurs africains ; « ne vaut que » signifie « n'a de valeur que si » : c'est un critère exclusif ; « dénoncer » : condamner publiquement une injustice (colonialisme, corruption, dictature).
\textbf{Étape 3 — Limites} : le sujet porte sur la \emph{valeur} de l'œuvre (critère esthétique et moral), non sur son succès commercial ; il exclut les autres genres (poésie, théâtre) sauf comparaison brève.
\textbf{Étape 4 — Tension} : d'un côté, la littérature africaine est née dans le combat (négritude, romans anticolonialistes comme ceux d'Oyono, de Mongo Beti) ; de l'autre, réduire le roman africain à la dénonciation, n'est-ce pas lui refuser le droit à l'imaginaire, à l'intime, à la beauté pure ? Un roman d'amour ou un conte fantastique africain ne « vaut »-il rien ?
\textbf{Problématique} : \emph{« La dénonciation est-elle la seule source de valeur du roman africain, ou celui-ci peut-il aussi s'accomplir dans d'autres formes et d'autres ambitions ? »}
Cette problématique appelle naturellement un plan dialectique : la dénonciation comme grandeur du roman africain (thèse), ses limites (antithèse), une conception ouverte de la valeur littéraire (synthèse).
\end{exoresolu}

\begin{exercice}[À vous de jouer]
Analysez le sujet suivant selon les quatre étapes, puis formulez une problématique : « L'humour est-il une arme efficace pour critiquer la société ? Vous appuierez votre réflexion sur \emph{Le Vieux nègre et la médaille}. »
\end{exercice}

\begin{corrige}[Exercice 1]
\textbf{Étape 1} : mots-clés : \emph{humour}, \emph{arme efficace}, \emph{critiquer la société}. \textbf{Étape 2} : humour : manière de présenter la réalité en la rendant ridicule ou drôle (ironie, satire) ; arme : moyen de combat ; efficace : qui atteint son but ; critiquer : juger en blâmant. \textbf{Étape 3} : limites : le sujet porte sur l'humour comme procédé de critique sociale, illustré par le roman d'Oyono ; une étude générale du comique sans lien avec la critique serait hors-sujet. \textbf{Étape 4} : tension : l'humour fait rire, donc adoucit et désamorce ; mais ce rire peut aussi frapper plus fort que la colère, car il désarme la censure et marque durablement les esprits. Le rire est-il compatible avec la gravité de la dénonciation ? \textbf{Problématique possible} : « Comment le rire, en apparence léger, devient-il chez Oyono un instrument redoutable de critique sociale, et quelles sont les limites de cette efficacité ? »
\end{corrige}

\begin{savaistu}[Oyono, une plume au service de la dénonciation]
Publié en 1956, \emph{Le Vieux nègre et la médaille} a été écrit alors que le Cameroun était encore sous administration coloniale. Ferdinand Oyono y dénonce, par l'ironie, l'hypocrisie du système colonial : Meka, vieil homme « fidèle », reçoit une médaille pour avoir donné ses terres et ses fils à la France, avant d'être brutalement ramené à sa condition d'opprimé. Le roman illustre parfaitement la problématique de cette leçon : la littérature peut-elle dénoncer efficacement ? L'histoire littéraire africaine répond oui : ce livre, traduit et étudié dans le monde entier, a durablement marqué les consciences.
\end{savaistu}

En résumé, analyser le sujet est la première victoire du candidat : quatre étapes (mots-clés, définitions, limites, tension) conduisent à une problématique précise, fidèle et engageante. C'est sur cette fondation que nous construirons, dans la prochaine section, la recherche et l'organisation des idées.

\coursec{Rechercher et organiser les idées}

Dans la section précédente, nous avons appris à analyser un sujet de dissertation et à formuler une problématique claire. Mais une problématique, aussi pertinente soit-elle, reste une question vide si l'on n'a rien à lui répondre. C'est ici que commence le véritable travail de réflexion : \cle{rechercher les idées}, c'est-à-dire mobiliser tout ce que l'on sait — lectures, expérience, actualité — pour construire une réponse argumentée. Cette étape se déroule entièrement au brouillon : elle est invisible dans la copie finale, mais c'est elle qui détermine la richesse ou la pauvreté de votre devoir. Un candidat qui saute cette étape écrit « au fil de la plume » et produit une dissertation désorganisée, pleine de répétitions et de hors-sujet. Un candidat méthodique, lui, arrive à la rédaction avec un trésor d'idées déjà triées, hiérarchisées et illustrées.

\courssub{Les techniques de mobilisation des idées}

\begin{definition}[La recherche des idées]
La \cle{recherche des idées} est l'étape de la dissertation qui consiste à faire jaillir, noter puis organiser toutes les réflexions susceptibles de répondre à la problématique. Elle comporte trois gestes successifs : \textbf{mobiliser} (faire venir les idées), \textbf{regrouper} (les rassembler par familles) et \textbf{organiser} (les trier et les hiérarchiser).
\end{definition}

\textbf{Le remue-méninges (brainstorming).} L'intuition est simple : le cerveau fonctionne par associations. Si l'on se donne un temps court (cinq à dix minutes) et l'interdiction de juger, les idées affluent. La règle d'or du \cle{remue-méninges} est donc : \emph{noter tout, ne rien censurer}. Une idée qui paraît ridicule au départ peut, une fois reformulée, devenir un excellent argument. Au brouillon, on écrit des mots, des bouts de phrases, des titres d'œuvres, des souvenirs d'actualité, sans se soucier de l'orthographe ni de l'ordre.

\textbf{Les questions-guides : Qui ? Quoi ? Pourquoi ? Comment ?} Pour éviter la page blanche, on interroge méthodiquement le sujet. Prenons le sujet fil conducteur de cette leçon : \emph{« La littérature peut-elle changer la société ? »}

\begin{center}
\begin{tabularx}{\linewidth}{|l|X|}
\hline
\rowcolor{popblueL} \textbf{Question} & \textbf{Exemples de réponses en vrac} \\
\hline
Qui ? & Les écrivains (Oyono, Mongo Beti, Sembène Ousmane), les lecteurs, les gouvernants, la jeunesse scolaire. \\
\hline
Quoi ? & La dénonciation de l'injustice, l'éveil des consciences, la transmission des valeurs, la satire. \\
\hline
Pourquoi ? & Parce que la littérature touche les émotions, rend visible ce qu'on cache, donne des modèles. \\
\hline
Comment ? & Par le rire satirique, par le témoignage, par l'école, par les adaptations au cinéma et au théâtre. \\
\hline
\end{tabularx}
\end{center}

\textbf{Le tableau Pour/Contre.} Quand le sujet appelle un débat (c'est le cas de la plupart des sujets de dissertation), on trace deux colonnes et l'on remplit chacune sans souci d'ordre. Cela garantit que l'on envisage les deux faces de la question — condition d'un plan dialectique équilibré.

\begin{center}
\begin{tabularx}{\linewidth}{|X|X|}
\hline
\rowcolor{popblueL} \textbf{Pour : la littérature change la société} & \textbf{Contre : ses limites} \\
\hline
Elle dénonce l'injustice coloniale (Oyono). & Elle touche peu de lecteurs (faible alphabétisation). \\
\hline
Elle éveille les consciences (indépendances). & Le pouvoir peut censurer ou ignorer les œuvres. \\
\hline
Elle éduque et transmet des valeurs. & Les changements réels exigent l'action politique. \\
\hline
\end{tabularx}
\end{center}

\begin{methode}[Rechercher les idées en cinq étapes]
\begin{enumerate}
\item \textbf{Brainstormer} : pendant 5 à 10 minutes, noter au brouillon toutes les idées qui viennent, sans en rejeter aucune.
\item \textbf{Regrouper} : relire la liste et rassembler par affinités les idées qui se ressemblent (on obtient des familles d'idées, futurs axes du plan).
\item \textbf{Trier} : éliminer les idées hors sujet, les doublons, les idées trop faibles ou invérifiables.
\item \textbf{Hiérarchiser} : ordonner les familles de la plus simple à la plus complexe, du constat à l'approfondissement.
\item \textbf{Illustrer} : associer à chaque idée retenue un argument et un exemple précis (citation, fait historique, situation vécue).
\end{enumerate}
\end{methode}

\courssub{Les sources des idées : où puiser ?}

Un élève de Terminale technique dispose de quatre grands réservoirs d'idées. Les mobiliser tous est le secret d'une copie riche.

\textbf{1. L'œuvre au programme.} \emph{Le Vieux nègre et la médaille} de Ferdinand Oyono est votre première mine d'exemples. Le roman montre précisément une littérature qui dénonce : à travers Meka, vieil homme décoré pour avoir donné ses terres et ses fils à la France, Oyono ridiculise l'hypocrisie coloniale et réveille la conscience du lecteur. La médaille, symbole de reconnaissance, devient symbole de mystification.

\textbf{2. Les autres œuvres lues.} Au collège et au lycée, vous avez rencontré d'autres textes : \emph{Une si longue lettre} de Mariama Bâ (condition de la femme), \emph{L'Enfant noir} de Camara Laye (tradition et éducation), les fables de La Fontaine (critique sociale par le récit animalier), \emph{Les Misérables} de Victor Hugo (dénonciation de la misère). Chacune est une réserve d'exemples.

\textbf{3. L'actualité camerounaise et africaine.} Les débats de société que vous suivez — éducation, justice, corruption, place des jeunes, rôle des artistes engagés — fournissent des exemples vivants qui prouvent au correcteur que vous reliez la littérature au monde réel.

\textbf{4. L'expérience personnelle.} Un livre qui vous a marqué, une pièce de théâtre vue au village, un poème récité à l'école : ces souvenirs, bien formulés, peuvent illustrer une idée, à condition de rester pertinents et brefs.

\begin{savaistu}[Ce que valorisent les correcteurs]
Au Probatoire et au Baccalauréat technique, les correcteurs lisent des centaines de copies. Ce qui distingue une excellente copie, ce ne sont pas les généralités (« la littérature est très importante dans la vie »), mais les \cle{exemples précis et variés} : une œuvre citée avec son auteur, un fait historique daté, une situation d'actualité bien choisie. Une copie qui cite Oyono, Hugo et un fait camerounais contemporain montre une culture large : elle est immédiatement mieux notée.
\end{savaistu}

\courssub{Regrouper, trier, hiérarchiser : l'organisation des idées}

\textbf{Regrouper par affinités.} Une fois la liste en vrac obtenue, on la relit et l'on rassemble les idées qui « vont ensemble ». Par exemple, « dénoncer l'injustice », « ridiculiser les colons » et « éveiller les consciences » forment une même famille : la fonction critique de la littérature. Chaque famille deviendra un axe (une partie) du plan.

\textbf{Trier.} Trois catégories d'idées doivent disparaître :
\begin{itemize}
\item les idées \textbf{hors sujet} (elles sont vraies, mais ne répondent pas à la problématique : par exemple, « la littérature coûte cher à l'école » pour notre sujet) ;
\item les \textbf{doublons} (deux formulations d'une même idée : on garde la meilleure) ;
\item les idées \textbf{trop faibles ou invérifiables} (rumeurs, opinions sans fondement, exemples dont on n'est pas sûr).
\end{itemize}

\textbf{Hiérarchiser.} Les axes retenus s'ordonnent de l'idée la plus simple à la plus complexe, du constat à l'approfondissement. Pour notre sujet : on commencera par montrer que la littérature agit sur les esprits (constat), puis on examinera ses limites (nuance), et l'on finira par définir les conditions de son efficacité (approfondissement). Cette progression donne au devoir un mouvement, une dynamique que le correcteur apprécie.

\begin{popfigure}
\begin{tikzpicture}[mindmap, grow cyclic, every node/.style={concept, text=white, font=\small}, concept color=popblue, level 1/.append style={level distance=3.4cm, sibling angle=90}, level 2/.append style={level distance=2.4cm, sibling angle=50}]
\node[font=\small\bfseries] {La littérature peut-elle changer la société ?}
  child[concept color=popgreen] { node {Dénoncer l'injustice}
    child { node[concept color=popgreen!70] {\emph{Le Vieux nègre et la médaille} : la médaille, symbole de mystification} }
    child { node[concept color=popgreen!70] {Hugo, \emph{Les Misérables} : la misère} }
  }
  child[concept color=poporange] { node {Éveiller et éduquer}
    child { node[concept color=poporange!70] {Éveil des consciences avant les indépendances} }
    child { node[concept color=poporange!70] {L'école transmet les valeurs} }
  }
  child[concept color=poppurple] { node {Ses limites}
    child { node[concept color=poppurple!70] {Peu de lecteurs, censure} }
    child { node[concept color=poppurple!70] {L'action politique reste nécessaire} }
  }
  child[concept color=popgold] { node {Conditions d'efficacité}
    child { node[concept color=popgold!70] {Un public formé et libre} }
    child { node[concept color=popgold!70] {Cinéma, théâtre, réseaux : nouveaux relais} }
  };
\end{tikzpicture}
\legende{Carte mentale de recherche d'idées autour du sujet « La littérature peut-elle changer la société ? ». Chaque branche principale est une famille d'idées (futur axe du plan) ; chaque sous-branche est un exemple précis.}
\end{popfigure}

\courssub{Associer à chaque idée un argument et un exemple}

Une idée seule ne suffit pas : dans une dissertation, chaque idée doit être \textbf{défendue par un argument} (le raisonnement qui la justifie) et \textbf{prouvée par un exemple} (la référence concrète qui l'illustre). La structure type d'un paragraphe est donc : \emph{idée → argument → exemple}.

\begin{exemplebox}[Idée + argument + exemple]
\textbf{Idée} : la littérature dénonce l'injustice. \textbf{Argument} : en mettant en scène des personnages victimes, elle rend visible et insupportable ce que le pouvoir préfère cacher. \textbf{Exemple} : dans \emph{Le Vieux nègre et la médaille}, Oyono montre Meka, décoré après avoir tout sacrifié, humilié puis abandonné par les autorités coloniales : le lecteur comprend, par le rire et la pitié, la profonde injustice de la colonisation.
\end{exemplebox}

\begin{attention}[L'erreur fréquente : l'accumulation d'exemples]
Beaucoup de candidats croient enrichir leur copie en empilant les exemples : « Oyono, mais aussi Beti, Bâ, Hugo, La Fontaine, Achebe... » sans jamais expliquer ce que chacun prouve. C'est une erreur grave. Un exemple sans argument est du remplissage ; le correcteur le sanctionne. Règle absolue : \textbf{chaque exemple doit servir une idée clairement formulée}. Mieux vaut deux exemples bien exploités que six titres jetés en vrac.
\end{attention}

\begin{methode}[Constituer une banque d'exemples par fiches thématiques]
Pour ne jamais être à court d'illustrations, constituez au fil de l'année des \cle{fiches par thème}. Chaque fiche comporte : le thème, une œuvre avec son auteur, une citation ou un passage précis, et un fait d'actualité ou historique. Exemple de fiches :
\begin{itemize}
\item \textbf{Colonialisme} : Oyono, \emph{Le Vieux nègre et la médaille} — la scène de la remise de la médaille — + les indépendances africaines de 1960.
\item \textbf{Justice} : Hugo, \emph{Les Misérables} — la condamnation de Jean Valjean — + les débats actuels sur l'accès à la justice au Cameroun.
\item \textbf{Éducation} : Camara Laye, \emph{L'Enfant noir} — l'école française et la tradition — + l'école gratuite et ses défis aujourd'hui.
\item \textbf{Tradition et modernité} : Mariama Bâ, \emph{Une si longue lettre} — le conflit entre coutume et émancipation — + les débats sur la place de la femme.
\end{itemize}
\end{methode}

\courssub{Vérifier l'équilibre et alimenter la banque collective}

Avant de passer au plan détaillé, un dernier contrôle s'impose : votre ensemble d'idées retenues doit comporter \textbf{au moins deux ou trois axes distincts et complémentaires}. Un seul axe donne un devoir pauvre ; quatre axes ou plus donnent un devoir dispersé. Vérifiez aussi que les axes ne se répètent pas et qu'ils répondent tous à la problématique.

Le \textbf{contrôle de lecture} du \emph{Vieux nègre et la médaille} joue ici un rôle décisif : le compte-rendu collectif en classe (personnages, thèmes, passages marquants, citations retenues) alimente la \cle{banque d'exemples de la classe}, affichée ou photocopiée, que chacun pourra mobiliser le jour de l'examen.

\begin{exoresolu}[De 14 idées en vrac à 3 axes organisés]
\textbf{Énoncé.} Sur le sujet « La littérature peut-elle changer la société ? », un élève a noté 14 idées en vrac : (1) dénoncer l'injustice ; (2) Oyono ridiculise les colons ; (3) les livres coûtent cher ; (4) éveiller les consciences ; (5) Hugo défend les pauvres ; (6) peu de gens lisent au village ; (7) la littérature éduque ; (8) éveiller les esprits ; (9) la censure existe ; (10) le théâtre au village sensibilise ; (11) l'école transmet des valeurs ; (12) j'aime lire ; (13) l'action politique est aussi nécessaire ; (14) les réseaux sociaux relaient les œuvres. Regrouper, trier et organiser ces idées en axes.
\tcblower
\textbf{Solution détaillée.}
\emph{Étape 1 — Tri.} On élimine (3), hors sujet (question économique), et (12), opinion personnelle invérifiable. On fusionne les doublons (4) et (8), qui disent la même chose. Il reste 11 idées utiles, que l'on condense en 6 idées fortes.
\emph{Étape 2 — Regroupement en 3 axes.}
\begin{itemize}
\item \textbf{Axe 1 — La littérature dénonce et éveille} : (1), (2), (4+8), (5). Exemples : Oyono, Hugo.
\item \textbf{Axe 2 — La littérature éduque et sensibilise} : (7), (10), (11). Exemples : l'école, le théâtre communautaire.
\item \textbf{Axe 3 — Ses limites et ses conditions} : (6), (9), (13), (14). Exemples : faible lectorat, censure, rôle de l'action politique et des nouveaux médias.
\end{itemize}
\emph{Étape 3 — Hiérarchisation.} On ordonne du constat à l'approfondissement : Axe 1 (ce que la littérature fait) → Axe 2 (comment elle forme) → Axe 3 (ses limites et conditions). Les 14 idées de départ sont devenues un squelette de plan cohérent : 6 idées, 3 axes, chacun pourvu d'exemples.
\end{exoresolu}

\begin{popfigure}
\begin{tikzpicture}
\draw[fill=popgreenL, draw=popgreen!70!black, rounded corners] (0,0) rectangle (3.9,5.2);
\draw[fill=popgoldL, draw=popgold!70!black, rounded corners] (4.3,0) rectangle (8.2,5.2);
\draw[fill=poppinkL, draw=poppink!70!black, rounded corners] (8.6,0) rectangle (12.5,5.2);
\node[font=\small\bfseries, text=popgreen!50!black] at (1.95,4.8) {À GARDER};
\node[font=\small\bfseries, text=popgold!60!black] at (6.25,4.8) {À FUSIONNER};
\node[font=\small\bfseries, text=poppink!60!black] at (10.55,4.8) {À ÉLIMINER};
\node[font=\scriptsize, text width=3.5cm, align=left] at (1.95,2.3) {%
\textbullet\ Dénoncer l'injustice\\[2pt]
\textbullet\ Oyono ridiculise les colons\\[2pt]
\textbullet\ Hugo défend les pauvres\\[2pt]
\textbullet\ L'école transmet des valeurs\\[2pt]
\textbullet\ La censure existe\\[2pt]
\textbullet\ L'action politique est nécessaire};
\node[font=\scriptsize, text width=3.5cm, align=left] at (6.25,2.9) {%
\textbullet\ « Éveiller les consciences » + « éveiller les esprits » $\rightarrow$ une seule idée\\[6pt]
\textbullet\ « Éduquer » + « le théâtre sensibilise » $\rightarrow$ fonction éducative};
\node[font=\scriptsize, text width=3.5cm, align=left] at (10.55,2.9) {%
\textbullet\ « Les livres coûtent cher » : hors sujet\\[6pt]
\textbullet\ « J'aime lire » : opinion personnelle invérifiable};
\end{tikzpicture}
\legende{Tableau de tri des idées sur l'exemple du sujet « La littérature peut-elle changer la société ? » : chaque idée en vrac est rangée dans l'une des trois colonnes avant le regroupement en axes.}
\end{popfigure}

\begin{exercice}[À vous de jouer]
Sur le sujet « Faut-il préserver les traditions ? », notez au brouillon au moins dix idées en vrac (remue-méninges de huit minutes), puis : 1) éliminez les idées hors sujet et les doublons ; 2) regroupez les idées restantes en deux ou trois axes ; 3) associez à chaque axe un argument et un exemple tiré de \emph{Le Vieux nègre et la médaille}, d'une autre œuvre ou de l'actualité camerounaise.
\end{exercice}

\begin{corrige}[Exercice : Faut-il préserver les traditions ?]
Éléments de correction attendus. \emph{Idées possibles en vrac} : les traditions unissent la communauté ; certaines pratiques sont dépassées ; Meka respecte les coutumes de son village ; la dot ; les langues locales disparaissent ; la modernité apporte le progrès ; les griots transmettent l'histoire ; l'école détourne les jeunes des coutumes ; les festivals culturels au Cameroun ; certaines traditions oppriment les femmes. \emph{Tri} : on élimine les idées vagues, on fusionne « les traditions unissent » et « les griots transmettent » (fonction de cohésion et de mémoire). \emph{Axes possibles} : Axe 1 — les traditions fondent l'identité et la cohésion (exemple : la solidarité villageoise autour de Meka ; les festivals culturels camerounais) ; Axe 2 — mais certaines traditions doivent évoluer (exemple : le sort des femmes dans \emph{Une si longue lettre}) ; Axe 3 — il faut donc un équilibre entre tradition et modernité (exemple : l'école et la coutume chez Camara Laye). Chaque axe est défendu par un argument et illustré par un exemple précis.
\end{corrige}

\begin{aretenir}[L'essentiel]
Rechercher les idées, c'est : \textbf{brainstormer} sans se censurer, \textbf{regrouper} par affinités, \textbf{trier} (hors sujet, doublons, idées faibles), \textbf{hiérarchiser} du simple au complexe, et \textbf{illustrer} chaque idée d'un argument et d'un exemple précis. Les sources d'idées sont l'œuvre au programme, les autres lectures, l'actualité et l'expérience personnelle. Le résultat attendu : deux ou trois axes distincts, équilibrés et complémentaires — la matière première du plan détaillé que nous construirons dans la prochaine section.
\end{aretenir}

\coursec{Élaborer et produire le plan détaillé}

Dans la section précédente, nous avons appris à rechercher les idées, à les trier et à les classer en grandes familles. Mais un paquet d'idées, même bien trié, n'est pas encore une dissertation. Il faut maintenant donner à ce matériau une \cle{architecture} : c'est l'étape décisive de l'élaboration du \cle{plan détaillé}. Un bâtisseur ne commence jamais à monter les murs sans plan ; de même, le candidat au Baccalauréat technique ne rédige jamais sa copie sans avoir d'abord construit, au brouillon, le plan complet de son devoir.

\courssub{Qu'est-ce qu'un plan détaillé ?}

\begin{definition}[Le plan détaillé]
Le \cle{plan détaillé} est l'organisation écrite et hiérarchisée de l'ensemble du devoir, rédigée au brouillon avant la copie. Il comprend : l'\cle{introduction} entièrement rédigée, les grandes \cle{parties} et \cle{sous-parties} formulées chacune en une phrase (avec l'argument et l'exemple prévus), les \cle{transitions} entre les parties, et la \cle{conclusion} entièrement rédigée.
\end{definition}

Autrement dit, le plan détaillé est la \emph{maquette} de la dissertation : il ne reste plus, sur la copie, qu'à développer chaque élément prévu. Le plan détaillé se distingue du simple \emph{plan apparent} (I, II, III) par le fait que tout le raisonnement y est déjà visible.

\begin{propriete}[Les trois qualités d'un bon plan]
Un bon plan de dissertation est :
\begin{itemize}
\item \cle{progressif} : chaque partie fait avancer le raisonnement (on ne répète jamais la même idée sous un autre habillage) ;
\item \cle{équilibré} : les grandes parties ont un poids comparable (pas une partie de trois pages et une autre de trois lignes) ;
\item \cle{homogène} : les parties ne se chevauchent pas (chaque idée a une place unique, et une seule).
\end{itemize}
\end{propriete}

\courssub{Les trois grands types de plans}

Au Baccalauréat technique, trois types de plans couvrent la quasi-totalité des sujets. Il faut les connaître par cœur et savoir lequel choisir.

\textbf{1. Le plan thématique.} Il examine successivement les différents \emph{aspects} d'une notion, d'un personnage ou d'une œuvre. Exemple de sujet : « Le personnage de Meka ». I. Meka, un vieillard naïf ; II. Meka, un homme humilié ; III. Meka, un homme révolté.

\textbf{2. Le plan dialectique.} Il oppose une \cle{thèse} à une \cle{antithèse}, puis les dépasse dans une \cle{synthèse}. C'est le plan de la \emph{discussion}, quand le sujet contient une opposition ou une alternative (« ou », « faut-il », « peut-on »). Exemple : I. Oui... ; II. Mais... ; III. En réalité / Au-delà...

\textbf{3. Le plan analytique.} Il suit le cheminement du réel : \emph{constat} (ou problème) $\rightarrow$ \emph{causes} $\rightarrow$ \emph{conséquences} (ou solutions). Exemple de sujet : « La pauvreté des villageois dans \emph{Le Vieux nègre et la médaille} ». I. Le constat : une misère visible ; II. Les causes : la domination coloniale ; III. Les conséquences : la révolte qui monte.

\begin{popfigure}
\begin{center}
\begin{tikzpicture}[
  grow=down, level distance=1.1cm,
  every node/.style={font=\small},
  racine/.style={rectangle, rounded corners, draw=popdark, fill=popblueL, text width=2.6cm, align=center, inner sep=3pt},
  branche/.style={rectangle, rounded corners, draw=popline, fill=popcream, text width=2.1cm, align=center, inner sep=2pt},
  titre/.style={font=\bfseries\small, text=popdark}
]
\node[titre] at (0,0.6) {Plan thématique};
\node[racine] at (0,0) {Aspects d'une notion}
  child { node[branche] {Aspect 1} }
  child { node[branche] {Aspect 2} }
  child { node[branche] {Aspect 3} };
\node[titre] at (6,0.6) {Plan dialectique};
\node[racine] at (6,0) {Discussion}
  child { node[branche] {Thèse} }
  child { node[branche] {Antithèse} }
  child { node[branche] {Synthèse} };
\node[titre] at (12,0.6) {Plan analytique};
\node[racine] at (12,0) {Explication}
  child { node[branche] {Constat} }
  child { node[branche] {Causes} }
  child { node[branche] {Conséquences} };
\end{tikzpicture}
\end{center}
\legende{Les trois grands types de plans de dissertation : le plan thématique (aspects), le plan dialectique (thèse/antithèse/synthèse), le plan analytique (constat/causes/conséquences).}
\end{popfigure}

\begin{propriete}[Critère de choix du plan]
Le type de plan dépend du \cle{type de sujet} :
\begin{itemize}
\item \cle{sujet de réflexion} (on invite à analyser, à décrire, à expliquer) : plan \textbf{thématique} ou \textbf{analytique} ;
\item \cle{sujet de discussion} (le sujet pose une alternative, contient « ou », « faut-il », « peut-on ») : plan \textbf{dialectique}.
\end{itemize}
Ce critère est exigible au Baccalauréat technique : le correcteur vérifie d'abord que le plan choisi correspond au sujet posé.
\end{propriete}

\begin{attention}[Piège : forcer le dialectique]
Beaucoup de candidats plaquent un plan dialectique sur tous les sujets. Or un sujet comme « Présentez le personnage de Meka » n'appelle aucune discussion : un plan thématique s'impose. À l'inverse, traiter « La médaille de Meka : récompense ou humiliation ? » en simple plan thématique, c'est ignorer l'alternative posée par le sujet. Lisez la \emph{forme} du sujet avant de choisir.
\end{attention}

\courssub{La structure de l'introduction et de la conclusion}

L'introduction fonctionne comme un \emph{entonnoir} : on part du large (une idée générale) pour aboutir au précis (l'annonce du plan). La conclusion fait le mouvement inverse.

\begin{definition}[Les quatre étapes de l'introduction]
\begin{enumerate}
\item L'\cle{amorce} (ou accroche) : une phrase générale liée au thème (vérité partagée, constat, citation brève) ;
\item la \cle{présentation du sujet} : on situe le sujet, on définit ses termes clés, on le rattache à l'œuvre ;
\item la \cle{problématique} : la question directrice, souvent formulée par une interrogation ;
\item l'\cle{annonce du plan} : on présente les grandes parties, sans les numéroter (« Nous verrons d'abord que..., puis que..., enfin que... »).
\end{enumerate}
\end{definition}

\begin{definition}[Les deux étapes de la conclusion]
\begin{enumerate}
\item Le \cle{bilan} : on répond clairement à la problématique en résumant le chemin parcouru ;
\item l'\cle{élargissement} (ou ouverture) : on ouvre vers une question nouvelle, sans y répondre (vers un autre aspect de l'œuvre, un autre texte, une réalité contemporaine).
\end{enumerate}
\end{definition}

\begin{popfigure}
\begin{center}
\begin{tikzpicture}[font=\small]
% Entonnoir introduction
\begin{scope}
\node[font=\bfseries, text=popdark] at (0,4.6) {Introduction (entonnoir)};
\draw[fill=popblueL, draw=popdark] (-2.6,4.0) rectangle (2.6,3.2);
\node at (0,3.6) {1. Amorce (idée générale)};
\draw[fill=popblueL!70, draw=popdark] (-2.0,3.0) rectangle (2.0,2.2);
\node at (0,2.6) {2. Présentation du sujet};
\draw[fill=poporangeL, draw=popdark] (-1.4,2.0) rectangle (1.4,1.2);
\node at (0,1.6) {3. Problématique};
\draw[fill=poporange!40, draw=popdark] (-0.9,1.0) rectangle (0.9,0.2);
\node at (0,0.6) {4. Annonce du plan};
\draw[->, thick, popdark] (3.0,3.6) -- (3.0,0.6) node[midway, right, align=left] {du général\\ vers le précis};
\end{scope}
% Entonnoir inversé conclusion
\begin{scope}[xshift=9.5cm]
\node[font=\bfseries, text=popdark] at (0,4.6) {Conclusion (entonnoir inversé)};
\draw[fill=popgreenL, draw=popdark] (-1.2,4.0) rectangle (1.2,3.2);
\node at (0,3.6) {1. Bilan (réponse)};
\draw[fill=popgreenL!60, draw=popdark] (-2.2,3.0) rectangle (2.2,2.2);
\node at (0,2.6) {2. Élargissement (ouverture)};
\draw[->, thick, popdark] (2.8,3.6) -- (2.8,2.6) node[midway, right, align=left] {du précis\\ vers le large};
\end{scope}
\end{tikzpicture}
\end{center}
\legende{Structure de l'introduction (du général vers le précis, en quatre étapes) et de la conclusion (du précis vers le large, en deux étapes).}
\end{popfigure}

\courssub{La structure du développement}

Le développement comporte \textbf{deux ou trois grandes parties} (jamais une seule, rarement plus de trois). Chaque grande partie est divisée en \textbf{deux ou trois sous-parties}, et chaque sous-partie obéit à un schéma fixe : \cle{idée} (une phrase qui affirme) $+$ \cle{argument} (le raisonnement qui justifie) $+$ \cle{exemple} (une référence précise tirée de l'œuvre ou de la culture).

\begin{methode}[Produire le plan détaillé en six temps]
\begin{enumerate}
\item \textbf{Analyser le sujet} : souligner les mots-clés, définir chaque terme, repérer le type de sujet (réflexion ou discussion) ;
\item \textbf{Problématiser} : transformer le sujet en une question directrice ;
\item \textbf{Brainstormer} : noter au brouillon toutes les idées, arguments et exemples qui viennent, sans les juger ;
\item \textbf{Trier} : barrer les idées hors sujet ou faibles, regrouper les idées voisines ;
\item \textbf{Choisir le type de plan} : thématique, dialectique ou analytique, selon le sujet ;
\item \textbf{Rédiger le plan détaillé} : introduction et conclusion entièrement rédigées, chaque sous-partie formulée en une phrase, transitions prévues.
\end{enumerate}
\end{methode}

\begin{attention}[La cohérence annonce / plan réel]
Erreur fréquente et grave : annoncer dans l'introduction un plan (« Nous verrons d'abord que la médaille est une récompense, puis qu'elle est une humiliation ») et suivre en réalité un autre plan. Le correcteur \emph{vérifie systématiquement} cette cohérence : une annonce non respectée est sanctionnée comme un défaut de construction. Rédigez l'annonce du plan \emph{après} avoir fixé définitivement vos parties, ou vérifiez-la en relisant.
\end{attention}

\courssub{Les transitions}

Entre deux grandes parties, la copie ne doit pas « sauter » brutalement. La \cle{transition} est une phrase (ou deux) qui \emph{boucle} la partie qui finit et \emph{annonce} la suivante. Modèle : « Si [bilan de la partie I], cependant [idée de la partie II] ». Exemple : « Si la médaille apparaît d'abord comme la consécration du mérite de Meka, la cérémonie révèle bientôt qu'elle est aussi un instrument d'humiliation. »

\begin{popfigure}
\begin{center}
\begin{tikzpicture}[font=\small,
  partie/.style={rectangle, rounded corners, draw=popdark, fill=popblueL, text width=3.4cm, align=center, inner sep=4pt},
  sp/.style={rectangle, draw=popline, fill=popcream, text width=2.9cm, align=center, inner sep=2pt, font=\footnotesize},
  trans/.style={font=\footnotesize\itshape, text=poporange, align=center},
  fl/.style={->, thick, poporange}]
\node[partie] (p1) at (0,0) {\textbf{I.} Grande partie 1};
\node[sp] (p1a) at (-2.1,-1.6) {A. idée + argument + exemple};
\node[sp] (p1b) at (2.1,-1.6) {B. idée + argument + exemple};
\draw[popdark] (p1) -- (p1a); \draw[popdark] (p1) -- (p1b);
\node[trans] (t1) at (0,-2.9) {transition 1 : bilan de I + annonce de II};
\draw[fl] (p1b.south) |- (t1.east);
\node[partie] (p2) at (0,-4.1) {\textbf{II.} Grande partie 2};
\node[sp] (p2a) at (-2.1,-5.7) {A. idée + argument + exemple};
\node[sp] (p2b) at (2.1,-5.7) {B. idée + argument + exemple};
\draw[popdark] (p2) -- (p2a); \draw[popdark] (p2) -- (p2b);
\draw[fl] (t1.south) -- (p2.north);
\node[trans] (t2) at (0,-7.0) {transition 2 : bilan de II + annonce de III};
\draw[fl] (p2b.south) |- (t2.east);
\node[partie] (p3) at (0,-8.2) {\textbf{III.} Grande partie 3};
\node[sp] (p3a) at (-2.1,-9.8) {A. idée + argument + exemple};
\node[sp] (p3b) at (2.1,-9.8) {B. idée + argument + exemple};
\draw[popdark] (p3) -- (p3a); \draw[popdark] (p3) -- (p3b);
\draw[fl] (t2.south) -- (p3.north);
\end{tikzpicture}
\end{center}
\legende{Maquette d'un plan détaillé complet : trois grandes parties, chacune divisée en deux sous-parties (idée + argument + exemple), reliées par des transitions fléchées.}
\end{popfigure}

\courssub{Exemple commenté et production guidée}

\begin{exemplebox}[Plan dialectique : « Faut-il récompenser la fidélité ? »]
\textbf{I. La récompense consacre le mérite.}
\quad A. Elle reconnaît publiquement les services rendus (exemple : la médaille officielle).
\quad B. Elle encourage les autres à la fidélité (argument : l'exemple forme la communauté).
\textbf{II. Mais la récompense peut corrompre ou humilier.}
\quad A. Elle peut acheter une fausse soumission (argument : récompenser, c'est parfois faire taire).
\quad B. Elle peut humilier celui qu'elle prétend honorer (exemple : Meka, exhibé puis moqué à la cérémonie de Doum).
\textbf{III. La vraie récompense est la dignité.}
\quad A. La fidélité vaut d'abord par elle-même (argument : elle n'a pas besoin de médaille).
\quad B. Meka retrouve sa dignité en rejetant les valeurs coloniales (exemple : son retour au village et à sa terre).
\end{exemplebox}

\begin{exoresolu}[Production guidée : « La médaille de Meka : récompense ou humiliation ? »]
Énoncé. À partir de votre lecture du \emph{Vieux nègre et la médaille} de Ferdinand Oyono, produisez le plan détaillé complet d'une dissertation sur le sujet : « La médaille de Meka : récompense ou humiliation ? »
\tcblower
\textbf{Solution détaillée.} Le sujet contient une alternative (« ou ») : c'est un sujet de \emph{discussion}, donc un \cle{plan dialectique}.

\emph{Introduction rédigée.} Dans les sociétés coloniales, les décorations officielles passent pour le signe suprême de la reconnaissance de l'État. Dans \emph{Le Vieux nègre et la médaille}, Ferdinand Oyono raconte comment Meka, vieillard du village de Doum, reçoit une médaille pour avoir donné ses terres à la mission et ses fils à la guerre. Mais cette distinction est-elle vraiment un honneur ? La médaille de Meka est-elle une récompense ou une humiliation ? Nous verrons d'abord que la médaille apparaît comme une récompense, puis qu'elle se révèle une humiliation, enfin qu'elle devient pour Meka l'occasion d'une prise de conscience.

\emph{Développement (sous-parties en une phrase).}
I. La médaille est d'abord vécue comme une récompense.
\quad A. Elle couronne officiellement une vie de sacrifices (exemple : terres données, deux fils morts à la guerre).
\quad B. Elle fait de Meka un personnage admiré et envié au village (exemple : la fierté des villageois, le parrainage de Kelara).
Transition : « Si la médaille flatte d'abord l'orgueil de Meka, la cérémonie elle-même en révèle le vrai visage. »
II. La médaille se révèle une humiliation.
\quad A. La cérémonie transforme Meka en objet de spectacle (exemple : ses pieds brûlés par ses chaussures neuves, les rires des Blancs).
\quad B. La nuit qui suit achève sa déchéance (exemple : son arrestation, sa nuit en cellule, les coups).
Transition : « Cette humiliation, toutefois, n'est pas une fin : elle ouvre les yeux du vieillard. »
III. L'humiliation devient prise de conscience.
\quad A. Meka comprend que la médaille ne paie pas ses vraies pertes (argument : ni terres ni fils ne reviendront).
\quad B. Le village partage cette lucidité et renoue avec la solidarité (exemple : l'accueil de Meka à son retour, la pluie qui tombe enfin).

\emph{Conclusion rédigée.} La médaille de Meka, présentée comme une récompense, se révèle donc une humiliation savamment organisée par l'administration coloniale ; mais en déchirant les illusions du vieillard, elle lui rend paradoxalement sa dignité et sa lucidité. On peut alors se demander si, dans \emph{Le Vieux nègre et la médaille}, le rire d'Oyono n'est pas lui-même une arme de libération.
\end{exoresolu}

\begin{savaistu}[Le brouillon, un investissement rentable]
Au Baccalauréat technique, les candidats qui rédigent un plan détaillé complet au brouillon gagnent en moyenne une trentaine de minutes de rédaction efficace sur la copie : ils n'hésitent plus, ne raturent presque pas et évitent la plupart des hors-sujets. Le brouillon n'est pas du temps perdu : c'est le temps où la dissertation se construit vraiment.
\end{savaistu}

\courssub{Remédiation : les erreurs fréquentes et leurs corrections}

\begin{attention}[Les trois maladies du plan]
\begin{itemize}
\item \cle{Plan déséquilibré} : une partie très longue, une autre réduite à rien. \emph{Remède} : vérifier au brouillon que chaque grande partie contient au moins deux sous-parties de poids voisin.
\item \cle{Parties qui se chevauchent} : la même idée revient dans deux parties. \emph{Remède} : au moment du tri, donner à chaque idée une place unique ; si deux titres de parties se ressemblent, les fusionner.
\item \cle{Absence de problématique} : le devoir enchaîne les idées sans répondre à une question. \emph{Remède} : formuler la problématique par écrit \emph{avant} de choisir le plan, et vérifier que chaque partie y répond.
\end{itemize}
\end{attention}

\begin{exercice}[À vous de jouer]
Voici trois sujets. Pour chacun, indiquez le type de sujet (réflexion ou discussion) et le type de plan le mieux adapté : 1. « Présentez le village de Doum dans \emph{Le Vieux nègre et la médaille}. » 2. « Le rire est-il une arme efficace contre l'injustice ? » 3. « Expliquez les causes de la révolte de Meka. »
\end{exercice}

\begin{corrige}[Exercice 1]
1. Sujet de réflexion (on invite à décrire et analyser) : plan \textbf{thématique} (par exemple : I. un village pauvre ; II. un village solidaire ; III. un village lucide). 2. Sujet de discussion (interrogation sur une valeur) : plan \textbf{dialectique} (thèse : le rire dénonce ; antithèse : le rire ne suffit pas ; synthèse : le rire éveille les consciences). 3. Sujet de réflexion portant sur des causes : plan \textbf{analytique} (constat : la révolte ; causes : les humiliations coloniales ; conséquences : la prise de conscience et la solidarité villageoise).
\end{corrige}

\begin{aretenir}[L'essentiel de la section]
Le plan détaillé est la maquette écrite du devoir : introduction et conclusion entièrement rédigées, sous-parties formulées en une phrase, transitions prévues. Trois types de plans : thématique (aspects), dialectique (thèse/antithèse/synthèse, pour les sujets de discussion), analytique (constat/causes/conséquences). Un bon plan est progressif, équilibré et homogène. L'introduction va du général au précis (amorce, sujet, problématique, annonce) ; la conclusion va du précis au large (bilan, élargissement). L'annonce du plan doit correspondre exactement au plan réellement suivi.
\end{aretenir}

\newpage

\coursec{Activité d'intégration}

\coursec{Dissertation : produire le plan détaillé d'une dissertation littéraire}

\courssub{Situation d'apprentissage : La semaine de la dissertation à Nkongsamba}

À l'approche du Baccalauréat technique, le lycée technique de Nkongsamba organise, comme chaque année au mois de mars, une « semaine de la dissertation ». Le proviseur a constaté, à la lecture des résultats des trois dernières sessions, que l'épreuve de français est celle où les candidats perdent le plus de points : sur 120 candidats présentés l'an dernier, 78 ont obtenu une note inférieure à 10/20 en dissertation, soit 65\,\% des candidats. Le principal reproche des correcteurs figure sur les relevés : « sujet mal analysé, problématique absente ou hors sujet, plan non annoncé ou déséquilibré ».

Le club littéraire du lycée, que préside ton groupe de Terminale, décide de réagir. Votre professeur de français vous propose de produire, en groupes de quatre élèves, un \cle{plan détaillé complet} qui servira de modèle affiché dans la salle polyvalente pendant toute la semaine. Le sujet retenu est un sujet de type Baccalauréat :

\begin{exemplebox}[Sujet proposé à la classe]
« Les récompenses officielles suffisent-elles à rendre justice aux hommes ? Vous répondrez en vous appuyant sur \emph{Le Vieux nègre et la médaille} de Ferdinand Oyono et sur votre culture personnelle. »
\end{exemplebox}

Chaque groupe dispose de deux heures de séance et du matériel suivant : le roman étudié en classe, la fiche de lecture établie lors du contrôle de lecture, le dictionnaire de la classe, et une feuille de format A3 pour la carte mentale. La présentation au tableau dure cinq minutes par groupe ; la classe évalue chaque production à l'aide de la grille critériée fournie par l'enseignant (problématique : 4 points ; pertinence des idées : 4 points ; équilibre du plan : 4 points ; qualité des exemples : 4 points ; qualité de la langue : 4 points ; total : 20 points).

\courssub{Objectifs de la leçon}

À l'issue de cette séquence, l'élève de Terminale technique sera capable de :

\begin{itemize}
\item analyser un sujet de dissertation de type Baccalauréat technique : dégager les \cle{mots-clés}, les \cle{limites} du sujet et la \cle{tension} qu'il contient ;
\item formuler une \cle{problématique} pertinente et la justifier ;
\item mobiliser l'œuvre étudiée en classe, \emph{Le Vieux nègre et la médaille} de Ferdinand Oyono, comme réservoir d'arguments et d'exemples ;
\item rechercher, trier et organiser les idées sous forme de carte mentale en distinguant \cle{contenu manifeste} et \cle{contenu latent} ;
\item choisir et justifier un \cle{type de plan} adapté au sujet (dialectique, thématique, analytique) ;
\item rédiger un \cle{plan détaillé} complet : introduction rédigée, parties et sous-parties formulées en phrases, transitions écrites, conclusion rédigée ;
\item évaluer la production d'autrui à l'aide d'une grille critériée, conformément à l'exigence du savoir-faire officiel : « rédiger et justifier une démarche, une réponse ou une conclusion ».
\end{itemize}

\courssub{Étape 1 : Lecture de l'œuvre support — Paratextes et lecture ouvroir}

\begin{savaistu}[Ferdinand Oyono et \emph{Le Vieux nègre et la médaille}]
Ferdinand Oyono (1929-2010), figure majeure des lettres camerounaises, publie \emph{Le Vieux nègre et la médaille} en 1956. Ce roman court, écrit à la veille des indépendances, dénonce avec ironie et lucidité le système colonial français. Le \cle{paratexte} (titre, couverture, quatrième de couverture, dédicace) oriente la lecture : le titre oppose « Vieux nègre » (condition humiliante, essentialisée par le colonisateur) et « médaille » (symbole de la reconnaissance officielle, mais dérisoire). La \cle{lecture ouvroir} (première lecture en classe, guidée par l'enseignant) permet d'identifier les personnages principaux (Meka, le Commandant, Kelara, les fils), le cadre spatio-temporel (Doum, village camerounais, années 1930-40, période coloniale), et le schéma narratif : situation initiale (attente de la médaille), élément déclencheur (cérémonie), péripéties (discours, pluie, fête), dénouement (arrestation, humiliation), situation finale (désillusion).
\end{savaistu}

\begin{attention}[Pièges de la lecture rapide]
Ne pas confondre \emph{récit} et \emph{thèse}. Le roman n'est pas un pamphlet politique explicite ; c'est une fiction qui \emph{montre} l'injustice par l'ironie (le discours du commandant, l'ignorance de l'administration sur l'identité réelle de Meka) et le symbolisme (la pluie, la médaille qui brûle la peau). Pour la dissertation, on ne résume pas l'intrigue : on \emph{exploite} des scènes précises comme arguments.
\end{attention}

\courssub{Étape 2 : Outils de langue I — Le référent et ses substituts, synonymie et antonymie}

\begin{aretenir}[Communication : Le référent et ses substituts]
Dans tout énoncé, le \cle{référent} est l'être ou la chose dont on parle. Pour éviter la répétition, la langue utilise des \cle{substituts} :
\begin{itemize}
\item \textbf{Pronoms personnels, possessifs, démonstratifs, relatifs, indéfinis} (il, le sien, celui-ci, qui, quelqu'un).
\item \textbf{Nominalisations / Périphrases} (le vieux paysan, le décoré, le père de famille, cet homme).
\item \textbf{Ellipses} (omission du référent déjà identifié).
\end{itemize}
\emph{Application au sujet :} « Les récompenses officielles suffisent-elles à rendre justice aux \textbf{hommes} ? » $\rightarrow$ \textbf{hommes} = référent. Substituts possibles dans la dissertation : \emph{les méritants, les récipiendaires, les citoyens, les serviteurs de l'État, ces derniers, ceux-ci}.
\end{aretenir}

\begin{aretenir}[Sémantique : Synonymie et antonymie]
\begin{itemize}
\item \cle{Synonymie} : relation entre mots (synonymes) de sens très proche (ex. : \emph{récompense / distinction / décoration}). Nuances : registre (formel/usuel), connotation (péjoratif/méloratif), intensité.
\item \cle{Antonymie} : relation d'opposition entre mots (antonymes). Types : \emph{contradiction} (vivant/mort), \emph{contrariété} (grand/petit — graduable), \emph{inverses} (acheter/vendre).
\end{itemize}
\emph{Utilité pour la dissertation :} définir les mots-clés du sujet, varier le vocabulaire (cohésion lexicale), préciser la tension (opposition thèse/antithèse).
\end{aretenir}

\courssub{Étape 3 : Outils de langue II — Hypéronymie/hyponymie, contenus manifestes/latents, texte argumentatif}

\begin{aretenir}[Sémantique : Hypéronymie et hyponymie]
\begin{itemize}
\item \cle{Hyperonyme} : terme au sens large (classe supérieure). Ex. : \emph{récompense}.
\item \cle{Hyponyme} : terme au sens plus précis (classe inférieure). Ex. : \emph{médaille, décorations, promotion, prime, titre honorifique}.
\end{itemize}
Relation : \emph{Une médaille EST UNE récompense}. \emph{Une récompense A POUR HYPONYMES : médaille, décorations...}
\emph{Utilité :} organiser les idées par catégories (plan thématique), éviter les répétitions, montrer l'étendue du sujet.
\end{aretenir}

\begin{aretenir}[Communication : Contenus manifestes et contenus latents]
\begin{itemize}
\item \cle{Contenu manifeste} : ce qui est dit explicitement, littéralement (ex. : « Le commandant remet la médaille à Meka »).
\item \cle{Contenu latent} : ce qui est suggéré, implicite, sous-entendu (ex. : « Le commandant ne connaît pas Meka » $\rightarrow$ latent : l'administration est indifférente aux individus ; la médaille est un acte bureaucratique vide).
\end{itemize}
\emph{En dissertation :} les arguments forts naissent souvent de l'analyse du contenu latent des exemples littéraires ou sociaux.
\end{aretenir}

\begin{aretenir}[Stylistique : Le texte argumentatif — Caractéristiques]
Un texte argumentatif vise à \cle{convaincre} (par la raison) ou à \cle{persuader} (par l'émotion). Sa structure type :
\begin{enumerate}
\item \textbf{Thèse} : position défendue (souvent en énoncé initial ou final).
\item \textbf{Arguments} : raisons logiques, faits, chiffres, citations d'autorité, exemples.
\item \textbf{Exemples} : illustrations concrètes qui ancrent l'argument dans le réel.
\item \textbf{Connecteurs logiques} : \emph{donc, car, en effet, cependant, par conséquent, de surcroît}.
\end{enumerate}
\emph{En dissertation :} chaque sous-partie = un argument + un exemple. La transition = opérateur logique entre arguments.
\end{aretenir}

\courssub{Étape 4 : Analyser le sujet et formuler la problématique}

\begin{methode}[Analyse d'un sujet de dissertation — 4 temps]
\begin{enumerate}
\item \textbf{Souligner les mots-clés} : termes porteurs de sens, souvent noms, verbes, adjectifs qualificatifs.
\item \textbf{Définir chaque mot-clé} : à l'aide de \cle{synonymes} (pour préciser le sens) et d'\cle{antonymes} (pour délimiter le champ). Recourir au dictionnaire si besoin.
\item \textbf{Fixer les limites} : identifier ce qui est \textbf{inclus} (ex. : récompenses \emph{officielles}, pas privées) et ce qui est \textbf{exclu} (ex. : on ne fait pas l'histoire des médailles, on évalue leur \emph{suffisance}).
\item \textbf{Dégager la tension} : repérer l'opposition interne au sujet (souvent : affirmation vs nuance, apparence vs réalité, idéal vs réel). La tension \textbf{est} le moteur de la problématique.
\end{enumerate}
\end{methode}

\begin{methode}[Formulation de la problématique]
\begin{enumerate}
\item Transformer la tension en \textbf{question ouverte} (commence par \emph{Comment, Pourquoi, Dans quelle mesure, La... est-elle... ?}).
\item Vérifier : la question \textbf{reprend-elle les mots-clés} ? \textbf{Respecte-t-elle les limites} ? \textbf{Appelle-t-elle une discussion} (pas une réponse oui/non simple) ?
\item Rédiger une \textbf{justification} (2-3 phrases) : « Cette problématique est pertinente car elle oppose... et engage à... »
\end{enumerate}
\end{methode}

\begin{exemplebox}[Application au sujet : « Les récompenses officielles suffisent-elles à rendre justice aux hommes ? »]
\textbf{Mots-clés :} \cle{récompenses officielles}, \cle{suffire}, \cle{rendre justice}, \cle{hommes}.

\begin{center}
\begin{tabularx}{\linewidth}{|l|X|X|}\hline
\rowcolor{popblueL} Mot-clé & Synonymes (au moins 2) & Antonymes (au moins 1) \\ \hline
Récompenses officielles & décorations, distinctions, médailles, honneurs publics, titres honorifiques & sanctions, punitions, oubli, mépris, indifférence \\ \hline
Suffire & être suffisant, contenter, satisfaire, combler, répondre aux attentes & manquer, être insuffisant, décevoir, faillir \\ \hline
Rendre justice & reconnaître le mérite, réparer, honorer à juste titre, réhabiliter, valider & commettre une injustice, méconnaître, humilier, bafouer, nier \\ \hline
Hommes & individus, personnes, citoyens, méritants, serviteurs, récipiendaires & (pas d'antonyme direct ; opposer : \emph{fonctionnaires, numéros, anonymes}) \\ \hline
\end{tabularx}
\end{center}

\textbf{Limites :}
\begin{itemize}
\item \emph{Inclus :} récompenses \textbf{officielles} (État, administration, armée), évaluation de leur \textbf{suffisance} (discussion critique), ancrage sur \emph{Le Vieux nègre et la médaille} + culture personnelle.
\item \emph{Exclus :} gratitude privée, récompenses symboliques non institutionnelles, simple description des cérémonies, biographie d'Oyono.
\end{itemize}

\textbf{Tension :} D'un côté, la récompense officielle \textbf{est} une forme de justice (reconnaissance publique, statut, fierté collective) ; de l'autre, elle \textbf{peut être} insuffisante, injuste, humiliante, ou purement symbolique (cas de Meka). L'opposition : \textbf{Apparence de justice / Réalité de l'injustice}.

\textbf{Problématiques candidates :}
\begin{enumerate}
\item \emph{Pourquoi les chefs aiment-ils les médailles ?} $\rightarrow$ \textbf{Rejetée} : hors sujet (psychologie des chefs), ne reprend pas les mots-clés, pas de tension.
\item \emph{La médaille est-elle toujours une récompense ?} $\rightarrow$ \textbf{Rejetée} : trop large, ne cible pas la « justice » ni la « suffisance ».
\item \textbf{\emph{La récompense officielle reconnaît-elle réellement le mérite de celui qui la reçoit, ou n'est-elle qu'un geste symbolique, parfois injuste et humiliant ?}} $\rightarrow$ \textbf{Retenue} : reprend les 3 mots-clés (récompense, justice/suffisance, hommes/mérite), traduit la tension, appelle une discussion en 3 temps (thèse/antithèse/synthèse).
\end{enumerate}
\end{exemplebox}

\courssub{Étape 5 : Rechercher et organiser les idées}

\begin{methode}[Recherche et organisation des idées — Carte mentale]
\begin{enumerate}
\item \textbf{Brainstorming} : noter toutes les idées, citations, exemples (œuvre + culture personnelle) sans trier.
\item \textbf{Distinction Manifeste / Latent} : pour chaque exemple du roman, écrire ce qui est \emph{vu/lu} (manifeste) et ce que cela \emph{signifie} (latent).
\item \textbf{Regroupement par arguments} : rassembler les exemples qui soutiennent la même idée directrice.
\item \textbf{Hiérarchisation} : ordonner les arguments du plus évident au plus profond (ou chronologique, ou logique).
\item \textbf{Représentation graphique} : carte mentale sur A3 (sujet au centre, 3 branches = 3 parties, sous-branches = sous-parties/exemples).
\end{enumerate}
\end{methode}

\begin{exemplebox}[Carte mentale des idées — Contenu minimum attendu]
\textbf{Branche 1 (Thèse : La récompense rend justice)}
\begin{itemize}
\item Argument : Reconnaissance publique du mérite et du sacrifice.
\item Exemple roman (Manifeste) : Médaille remise à Meka pour « services rendus » (don de la terre à la mission, fils morts à la guerre).
\item Exemple roman (Latent) : Le village de Doum célèbre Meka ; fierté collective ; modèle pour la jeunesse.
\item Exemple extérieur : Décorations nationales (Ordre de la Valeur, Ordre du Mérite) au Cameroun ; inspiration pour les jeunes générations ; statut social acquis.
\end{itemize}

\textbf{Branche 2 (Antithèse : La récompense est insuffisante / injuste)}
\begin{itemize}
\item Argument : Geste vide, sans connaissance réelle du méritant.
\item Exemple roman (Manifeste) : Discours pompeux du commandant (ne connaît pas le prénom de Meka, confond les faits).
\item Exemple roman (Latent) : Administration indifférente, bureaucratie coloniale, Meka réduit à un « numéro ».
\item Argument : Ne change rien à la condition ; masque l'injustice.
\item Exemple roman (Manifeste) : Pluie qui gâche la fête ; arrestation et humiliation de Meka la nuit même (vol de la médaille, prison).
\item Exemple roman (Latent) : La médaille ne protège de rien ; elle devient une cible (vol) et un fardeau.
\item Exemple extérieur : Décorations controversées (personnalités politiques/affairistes) vs oubli des « petits » méritants (enseignants en zone rurale, infirmiers, anciens combattants non reconnus).
\end{itemize}

\textbf{Branche 3 (Synthèse : La récompense ne vaut que par la reconnaissance réelle)}
\begin{itemize}
\item Argument : Vraie justice = connaissance + respect de la personne + actes concrets durables.
\item Exemple roman (Latent) : Ce qu'il manquait à Meka : pas la médaille, mais le respect de sa terre, de sa famille, de sa dignité.
\item Exemple extérieur : Reconnaissance tardive des tirailleurs sénégalais/camerounais (pensions, naturalisation, mémoire) ; politiques de réparation mémorielle.
\end{itemize}
\end{exemplebox}

\courssub{Étape 6 : Élaborer et produire le plan détaillé}

\begin{methode}[Choix du type de plan — Critères de décision]
\begin{itemize}
\item \textbf{Plan dialectique (Thèse / Antithèse / Synthèse)} : Sujet \textbf{question fermée} (« Est-ce que... ? », « Suffit-il... ? », « Faut-il... ? ») $\rightarrow$ tension binaire évidente $\rightarrow$ discussion équilibrée.
\item \textbf{Plan thématique (Aspect 1 / Aspect 2 / Aspect 3)} : Sujet \textbf{question ouverte} (« Quels sont les effets de... ? », « Comment... ? ») ou sujet à \textbf{plusieurs dimensions indépendantes} (social, psychologique, politique).
\item \textbf{Plan analytique (Constat / Causes / Conséquences \textbf{ou} Situation / Problème / Solutions)} : Sujet \textbf{problème à résoudre} ou \textbf{phénomène à expliquer}.
\end{itemize}
\emph{Justification obligatoire :} « J'ai choisi le plan dialectique car le sujet "Les récompenses officielles suffisent-elles... ?" est une question fermée qui oppose une apparence de justice (thèse) à une réalité d'insuffisance (antithèse), appelant une dépassement (synthèse). »
\end{methode}

\begin{methode}[Rédaction du plan détaillé — Structure complète]
\begin{enumerate}
\item \textbf{Introduction rédigée} (1 paragraphe) : \cle{Amorce} (généralité, citation, actualité) $\rightarrow$ \cle{Présentation du sujet} (œuvre, auteur, résumé ciblé) $\rightarrow$ \cle{Problématique} (la question retenue) $\rightarrow$ \cle{Annonce du plan} (les 3 grandes parties en phrases nominales ou verbales).
\item \textbf{Parties (2 ou 3)} : Titre de partie (phrase affirmative résumant l'argument majeur). \textbf{Sous-parties} (2 ou 3 par partie) : \textbf{formulées en phrases complètes} (argument + exemple implicite). Pas de mots seuls.
\item \textbf{Transitions écrites} (entre chaque partie) : 1 phrase qui \textbf{ferme} la partie précédente et \textbf{ouvre} la suivante (reprise de mot-clé, connecteur logique : \emph{Cependant, Néanmoins, Au-delà de cela, Dès lors}).
\item \textbf{Conclusion rédigée} (1 paragraphe) : \cle{Bilan} (réponse synthétique à la problématique) $\rightarrow$ \cle{Ouverture} (autre œuvre, autre contexte, question prospective, citation).
\end{enumerate}
\end{methode}

\begin{attention}[Erreurs fréquentes à éviter dans le plan détaillé]
\begin{itemize}
\item \textbf{Sous-parties en mots-clés} : « 1. La médaille de Meka » $\rightarrow$ \textbf{Corrigé} : « 1. La médaille de Meka couronne officiellement le don de sa terre et la mort de ses fils. »
\item \textbf{Transitions manquantes ou « passerelles » vides} : « Passons à la deuxième partie. » $\rightarrow$ \textbf{Corrigé} : « Mais cette justice apparente résiste-t-elle à l'épreuve des faits vécus par Meka ? »
\item \textbf{Déséquilibre} : 3 sous-parties dans la première partie, 1 seule dans la deuxième. Viser 2 ou 3 sous-parties \textbf{par partie}.
\item \textbf{Introduction/Conclusion non rédigées} : simples schémas. L'examen exige le texte intégral.
\item \textbf{Hors-sujet dans les exemples} : citer un exemple qui ne prouve pas l'argument de la sous-partie.
\end{itemize}
\end{attention}

\courssub{Exercice d'application et corrigé commenté}

\begin{exoresolu}[Correction collective harmonisée : du sujet au plan détaillé]
Énoncé : Produire, à partir du sujet « Les récompenses officielles suffisent-elles à rendre justice aux hommes ? », le plan détaillé complet d'une dissertation littéraire, en s'appuyant sur \emph{Le Vieux nègre et la médaille} et sur la culture personnelle.
\tcblower
\textbf{Étape 1 : Analyse du sujet (rappel synthétique).}

Mots-clés définis par synonymie/antonymie (voir tableau Étape 4). Limites : récompenses \emph{officielles}, évaluation de la \emph{suffisance}. Tension : Apparence de justice (reconnaissance publique) vs Réalité d'injustice (humiliation, indifférence, symbole vide).

\textbf{Étape 2 : Problématique retenue et justifiée.}

\emph{« La récompense officielle reconnaît-elle réellement le mérite de celui qui la reçoit, ou n'est-elle qu'un geste symbolique, parfois injuste et humiliant ? »}

\textbf{Justification :} Cette formulation reprend les termes « récompense officielle », « mérite » (justice) et « homme » (récepteur). Elle traduit l'opposition entre la valeur symbolique attendue (thèse) et la vacuité vécue (antithèse), ouvrant la voie à une redéfinition de la justice (synthèse).

\textbf{Étape 3 : Choix du plan et justification.}

\emph{Choix : Plan dialectique (Thèse / Antithèse / Synthèse).}

\emph{Justification :} Le sujet est une question fermée (« suffisent-elles ? ») qui appelle une discussion contradictoire. La thèse affirme la fonction de justice ; l'antithèse en démontre les limites et les dérives ; la synthèse propose une condition de validité (la reconnaissance réelle).

\textbf{Étape 4 : Plan détaillé rédigé.}

\emph{Introduction} : Dans toutes les sociétés, les autorités décernent des médailles, des décorations et des distinctions pour honorer les mérites. Au Cameroun comme ailleurs, les cérémonies officielles de remise de récompenses rythment la vie nationale et façonnent l'imaginaire collectif. Mais ces honneurs rendent-ils vraiment justice à ceux qui les reçoivent ? C'est toute la question que pose Ferdinand Oyono dans \emph{Le Vieux nègre et la médaille}, où le vieux Meka reçoit une médaille pour ses services rendus à l'administration coloniale, avant d'être humilié la nuit même de sa décoration. On peut alors se demander : la récompense officielle reconnaît-elle réellement le mérite de celui qui la reçoit, ou n'est-elle qu'un geste symbolique, parfois injuste et humiliant ? Nous verrons d'abord que la récompense officielle est une forme de justice, puis qu'elle se révèle souvent insuffisante, avant de montrer qu'elle ne vaut que si elle s'accompagne d'une reconnaissance réelle et durable.

\emph{Première partie — Thèse : Les récompenses officielles rendent justice aux hommes.}
\begin{itemize}
\item \textbf{Sous-partie 1 :} Elles reconnaissent publiquement le mérite et le sacrifice, inscrivant l'individu dans la mémoire collective (exemple : la médaille de Meka couronne officiellement le don de sa terre à la mission et la mort de ses fils au combat).
\item \textbf{Sous-partie 2 :} Elles confèrent un statut social, une fierté légitime et offrent un modèle d'engagement à la communauté (exemple : le village de Doum célèbre Meka ; les décorations nationales inspirent la jeunesse et valident l'idéal civique).
\end{itemize}
\emph{Transition :} Mais cette justice apparente, fondée sur le symbole et le cérémonial, résiste-t-elle à l'épreuve des faits et de la vérité vécue par les récipiendaires ?

\emph{Deuxième partie — Antithèse : Les récompenses officielles sont souvent insuffisantes, voire injustes.}
\begin{itemize}
\item \textbf{Sous-partie 1 :} Elles peuvent n'être qu'un geste bureaucratique vide, sans connaissance réelle ni considération pour la personne honorée (exemple : le commandant prononce un discours convenu, ignore le prénom de Meka et confond ses mérites).
\item \textbf{Sous-partie 2 :} Elles ne changent rien à la condition matérielle et morale du récompensé et peuvent même masquer ou aggraver l'injustice (exemple : la pluie gâche la fête ; Meka est arrêté, dépouillé de sa médaille et humilié la nuit même ; le métal ne le protège pas de l'arbitraire).
\item \textbf{Sous-partie 3 :} Elles sont parfois mal distribuées, récompensant la proximité au pouvoir plutôt que le mérite réel (exemple : décorations accordées à des personnalités contestées pendant que des serviteurs discrets de l'État — enseignants de brousse, agents de santé — restent dans l'ombre).
\end{itemize}
\emph{Transition :} Faut-il alors rejeter toute valeur à la récompense officielle ? Non : il faut en redéfinir le sens et les conditions de légitimité.

\emph{Troisième partie — Synthèse : La récompense ne vaut que par la reconnaissance réelle qui l'accompagne.}
\begin{itemize}
\item \textbf{Sous-partie 1 :} Une vraie justice exige de connaître et de respecter la personne honorée dans sa singularité, au-delà du symbole protocolaire (exemple : ce qui manquait à Meka, ce n'était pas la médaille, mais le respect de sa parole, de sa terre et de sa dignité d'homme libre).
\item \textbf{Sous-partie 2 :} La récompense doit s'accompagner d'actes concrets et durables : amélioration des conditions de vie, mémoire entretenue, transmission aux générations futures (exemple : la reconnaissance tardive mais effective des anciens combattants africains — pensions, naturalisation, lieux de mémoire).
\end{itemize}

\emph{Conclusion} : Les récompenses officielles peuvent rendre justice lorsqu'elles reconnaissent sincèrement le mérite et s'incarnent dans le respect quotidien ; mais l'exemple de Meka montre qu'une médaille sans considération réelle n'est qu'un morceau de métal froid, incapable de réparer l'humiliation. La vraie justice se mesure au respect concret accordé aux hommes, non aux seules fastes des cérémonies. \emph{Ouverture :} N'est-ce pas là une leçon pour toutes les institutions — école, armée, administration — qui honorent leurs serviteurs : le symbole ne doit jamais remplacer la substance ?
\end{exoresolu}

\courssub{Bilan de la leçon et évaluation}

\begin{aretenir}[Acquis essentiels pour le Baccalauréat technique]
\begin{itemize}
\item Un sujet ne se traite qu'après une \cle{analyse rigoureuse} : mots-clés définis par synonymie et antonymie, limites repérées, tension dégagée.
\item La \cle{problématique} n'est pas une question quelconque : elle doit traduire exactement la tension du sujet et en reprendre le vocabulaire.
\item Les idées se construisent en croisant l'\cle{œuvre étudiée} (contenu manifeste + contenu latent) et la \cle{culture personnelle}, organisées par \cle{hypéronymie/hyponymie}.
\item Le \cle{type de plan} se choisit en fonction de la structure du sujet (question fermée $\rightarrow$ dialectique) et se justifie explicitement.
\item Un \cle{plan détaillé} est une maquette complète de la copie : introduction et conclusion \textbf{rédigées}, sous-parties \textbf{formulées en phrases}, transitions \textbf{écrites}.
\item L'évaluation par \cle{grille critériée} (problématique, idées, équilibre, exemples, langue) apprend à justifier ses choix et à juger objectivement une production, compétence directement transposable à la relecture de sa propre copie le jour de l'examen.
\end{itemize}
\end{aretenir}

\begin{exercice}[Entraînement individuel — Devoir maison]
Traitez le sujet suivant en produisant \textbf{uniquement le plan détaillé complet} (introduction, plan avec sous-parties en phrases, transitions, conclusion) :
\begin{quote}
« L'école peut-elle à elle seule assurer l'ascension sociale ? Vous répondrez en vous appuyant sur \emph{Le Vieux nègre et la médaille} (parcours des fils de Meka) et sur votre culture personnelle. »
\end{quote}
\emph{Indications :} Réutilisez la méthode en 6 étapes. Attention : le sujet porte sur l'école et l'ascension sociale, la médaille n'est qu'un contre-exemple ou un point de comparaison. Prévoyez 1h30.
\end{exercice}

\begin{corrige}[Exercice n°1 — Éléments de correction attendus]
\textbf{Analyse :} Mots-clés : École, seule, assurer, ascension sociale. Tension : École = levier théorique d'ascension (méritocratie) vs École = reproduit les inégalités / insuffisante sans capital économique/social/culturel (Bourdieu). Limites : « à elle seule » $\rightarrow$ discuter les facteurs complémentaires.
\textbf{Problématique :} « L'institution scolaire suffit-elle à garantir la mobilité sociale, ou ne fait-elle que légitimer les inégalités d'origine ? »
\textbf{Plan :} Dialectique.
I. Thèse : L'école est le principal levier d'ascension (diplômes, mérite, exemples fils de Meka instruits vs père analphète).
II. Antithèse : L'école seule ne suffit pas (reproduction sociale, coût des études, chômage des diplômés, Meka humilié malgré les sacrifices pour l'école).
III. Synthèse : L'ascension réelle nécessite l'articulation École + Politiques publiques (bourse, emploi) + Capital familial.
\textbf{Introduction/Conclusion rédigées attendues.}
\end{corrige}

\newpage

\coursec{Méthodes et exercices résolus}

\courssub{Lecture de l'œuvre support : paratexte et lecture ouvroir}

\begin{methode}[Analyser les paratextes et effectuer une lecture ouvroir]
\textbf{Étape 1 : L'analyse des paratextes (avant lecture)}
\begin{enumerate}
    \item \textbf{Titre et sous-titre :} Identifier le genre (roman, nouvelle), le ton (ironique, tragique), les mots-clés (\cle{Vieux nègre}, \cle{médaille}). Hypothétiser sur le thème (vieillesse, colonisation, décorations, ironie).
    \item \textbf{Auteur et contexte :} Situer Ferdinand Oyono (Cameroun, 1929–2010), figure de la littérature africaine d'expression française, témoin de la colonisation et de l'indépendance. Noter la date de publication (1956) : période coloniale tardive.
    \item \textbf{Couverture / Illustration :} Décrire l'image (souvent un vieil homme en uniforme, médaille au cou), analyser la symbolique (contraste entre la misère physique et l'apparat militaire).
    \item \textbf{Quatrième de couverture / Préface :} Repérer l'intention annoncée (dénonciation, humour, tragédie), le résumé non spoilant, le public visé.
\end{enumerate}
\textbf{Étape 2 : La lecture ouvroir (première lecture continue)}
\begin{enumerate}
    \item Lire d'une traite sans s'arrêter au vocabulaire difficile pour saisir l'intrigue globale.
    \item Repérer : le \cle{schéma narratif} (situation initiale : Meka attend la médaille ; élément déclencheur : la cérémonie ; péripéties : préparatifs, discours, remise ; dénouement : mort de Meka / ironie finale), le \cle{point de vue} (narrateur omniscient, ironie narrative), les \cle{personnages principaux} (Meka, Kelara, le Commandant, le Prêtre), le \cle{cadre spatio-temporel} (village camerounais, années 50).
    \item Noter ses impressions immédiates (émotion, incompréhension, révolte, rire).
\end{enumerate}
\textbf{Étape 3 : Le compte-rendu de lecture ouvroir}
Rédiger une fiche synthétique : Fiche d'identité (titre, auteur, date, genre) + Résumé de l'intrigue (10 lignes max) + Premières hypothèses d'interprétation (thèmes, visée de l'auteur).
\end{methode}

\begin{exoresolu}[Analyse des paratextes et lecture ouvroir de \emph{Le Vieux nègre et la médaille}]
\textbf{Énoncé :} À partir des éléments de paratexte suivants (simulés), rédigez la fiche de lecture ouvroir type attendue au Probatoire.
\begin{itemize}
    \item Titre : \emph{Le Vieux nègre et la médaille}
    \item Auteur : Ferdinand Oyono
    \item Couverture : Un vieillard noir, chemise déchirée, short usé, coiffé d'un casque colonial trop grand, une médaille clinquante pendue au cou par un ruban défraîchi. Fond : cases en banco, soleil écrasant.
    \item Quatrième de couverture : « Meka a bien servi la France. Aujourd'hui, on lui donne une médaille. Mais à quel prix ? Une satire féroce de l'assimilation. »
\end{itemize}
\tcblower
\textbf{Solution détaillée :}

\textbf{FICHE DE LECTURE OUVROIR}

\textbf{1. Fiche d'identité}
\begin{itemize}
    \item \textbf{Titre :} \emph{Le Vieux nègre et la médaille}
    \item \textbf{Auteur :} Ferdinand Oyono (Cameroun, 1929–2010), haut fonctionnaire, diplomate, écrivain majeur de la négritude critique.
    \item \textbf{Date de publication :} 1956 (Éditions René Julliard, prix Sainte-Beuve).
    \item \textbf{Genre :} Roman court / Nouvelle longue, \cle{roman colonial} / \cle{anti-colonial}, \cle{satire}, \cle{réalisme critique}.
    \item \textbf{Contexte historique :} Cameroun sous tutelle française (après mandat SDN), montée des nationalismes, guerre d'Algérie débutée (1954), loi-cadre Defferre (1956).
\end{itemize}

\textbf{2. Analyse de la couverture (paratexte iconique)}
\begin{itemize}
    \item \textbf{Contraste saisissant :} La misère vestimentaire (chemise déchirée, short usé) vs l'apparat colonial (casque trop grand, médaille clinquante). Symbolise l'\cle{aliénation} : l'habit ne fait pas le moine, la décoration ne masque pas la pauvreté.
    \item \textbf{Le casque colonial :} Symbole de l'autorité coloniale, trop grand pour la tête du vieux -> inadéquation, ridicule, poids de l'histoire.
    \item \textbf{Le ruban défraîchi :} La médaille a déjà servi ? Ou le temps a passé sans que la condition du « vieux nègre » ne change.
    \item \textbf{Cadre :} Cases en banco, soleil -> ancrage réaliste, misère du village africain.
\end{itemize}

\textbf{3. Analyse du titre et de la 4e de couverture}
\begin{itemize}
    \item \textbf{Titre :} « Le Vieux nègre » (déshumanisation, appellation coloniale raciste, âge = expérience/souffrance) + « la médaille » (objet unique, but final, symbole vide). Article défini « la » : médaille spécifique, attendue, unique.
    \item \textbf{4e couv :} « Bien servi la France » -> ironie (servitude). « À quel prix ? » -> annonce du thème du sacrifice (dignité, fils, terre). « Satire féroce » -> ton annoncé : humour noir, dénonciation.
\end{itemize}

\textbf{4. Hypothèses de lecture (avant lecture ouvroir)}
\begin{itemize}
    \item Thème central : L'illusion de l'assimilation / la collaboration vaine.
    \item Structure : Journée de cérémonie (unité de temps/lieu) -> tension dramatique.
    \item Fin prévisible : Mort ou désillusion totale du personnage principal.
\end{itemize}

\textbf{5. Compte-rendu de la lecture ouvroir (après lecture)}
\begin{itemize}
    \item \textbf{Résumé :} Meka, vieux chef de canton loyal, attend la médaille de la Légion d'honneur. Il a sacrifié ses fils (l'un mort à la guerre, l'autre prêtre « blanc »), ses terres. La cérémonie a lieu. Discours hypocrites du Commandant et du Prêtre. Meka reçoit la médaille, mais meurt d'une crise cardiaque en rentrant, réalisant la vanité de ce métal. Kelara, sa femme, récupère la médaille pour la vendre.
    \item \textbf{Impressions :} Malaise croissant. Rire jaune devant les préparatifs ridicules (costume trop grand). Colère contre l'hypocrisie des Blancs (Commandant) et du Prêtre (complice). Tristesse finale : la médaille ne nourrit pas, ne soigne pas.
    \item \textbf{Thèmes confirmés :} Aliénation, collaboration, vanité des honneurs coloniaux, condition de la femme (Kelara), mort.
\end{itemize}

\textbf{Conclusion :} La lecture ouvroir confirme la visée satirique et tragique. L'œuvre dépasse l'anecdote pour devenir une allégorie de l'histoire camerounaise : le « Vieux nègre » est le peuple colonisé, la médaille est l'indépendance factice ou les miettes du colonisateur.
\end{exoresolu}

\courssub{Outils de langue I : référent, substituts, synonymie et antonymie}

\begin{methode}[Maîtriser la référence, la substitution et les rapports de sens lexicaux]
\textbf{A. Le référent et ses substituts (Cohérence textuelle)}
\begin{enumerate}
    \item \textbf{Identifier le référent :} Nommer l'entité réelle ou fictive désignée (personne, objet, idée). Ex : « \emph{Meka} » (référent).
    \item \textbf{Repérer les substituts (coréférence) :} Mots qui reprennent le référent pour éviter la répétition.
    \begin{itemize}
        \item \textbf{Pronoms personnels :} il, lui, le, se, y, en.
        \item \textbf{Pronoms possessifs :} le sien, la sienne.
        \item \textbf{Pronoms démonstratifs :} celui-ci, celui-là, ce dernier.
        \item \textbf{Substituts nominaux (renvois nominaux) :} \cle{le vieux}, \cle{le récipiendaire}, \cle{le vétéran}, \cle{le père de famille}, \cle{le chef de canton}.
        \item \textbf{L'ellipse (omission) :} Sujet supprimé dans une conjonction (ex : « Meka attend \emph{\emptyset} espère »).
    \end{itemize}
    \item \textbf{Vérifier la chaîne de référence :} S'assurer que l'enchaînement est clair (pas d'ambiguïté sur l'antécédent).
\end{enumerate}

\textbf{B. Synonymie et antonymie (Cohésion lexicale et argumentation)}
\begin{enumerate}
    \item \textbf{Synonymie (sens proche) :} Repérer les mots de même famille sémantique pour varier le lexique ou nuancer.
    \begin{itemize}
        \item \textbf{Synonymes parfaits (rares) :} \emph{médaille} / \emph{décoration} (contexte cérémonial).
        \item \textbf{Synonymes contextuels :} \emph{vieux} / \emph{âgé} / \emph{ancien} / \emph{vétéran}.
        \item \textbf{Effet :} Éviter la répétition, créer une progression (euphémisme $\to$ terme cru), construire le portrait.
    \end{itemize}
    \item \textbf{Antonymie (sens opposé) :} Repérer les oppositions pour structurer l'argumentation (thèse/antithèse) ou créer de l'ironie.
    \begin{itemize}
        \item \textbf{Antonymes gradables :} \emph{riche} / \emph{pauvre}, \emph{grand} / \emph{petit}.
        \item \textbf{Antonymes complémentaires :} \emph{vivant} / \emph{mort}, \emph{vrai} / \emph{faux}.
        \item \textbf{Antonymes converses :} \emph{donner} / \emph{recevoir}, \emph{acheter} / \emph{vendre}.
        \item \textbf{Effet :} Mettre en relief le contraste (ex : « \emph{la médaille d'or} » vs « \emph{la misère noire} »).
    \end{itemize}
\end{enumerate}
\textbf{Reflexe Bac :} Dans l'exercice de contraction/résumé, utiliser la substitution et la synonymie pour réduire le texte sans en perdre le sens.
\end{methode}

\begin{exoresolu}[Travail sur la référence et la lexicologie dans un extrait]
\textbf{Énoncé :} Lisez l'extrait suivant (début du roman, adapté) :
« \emph{Meka} était assis sur son siège. \emph{Le vieux chef} attendait. \emph{Il} tournait \emph{sa} médaille entre \emph{ses} doigts. \emph{Le récipiendaire} ne quittait pas des yeux \emph{l'objet} brillant. \emph{Ce dernier} représentait tout \emph{son} passé. Soudain, \emph{Kelara} arriva. \emph{Elle} posa \emph{sa} main sur \emph{l'épaule} de \emph{son mari}. »
\begin{enumerate}
    \item Construisez le tableau des chaînes de référence pour \emph{Meka} et \emph{la médaille}.
    \item Relevez deux paires synonymiques et une paire antonymique (implicite ou contextuelle) dans le texte ou le champ lexical de l'extrait.
    \item Réécrivez le passage en remplaçant les substituts par le référent principal pour mesurer l'alourdissement.
\end{enumerate}
\tcblower
\textbf{Solution détaillée :}

\textbf{1. Tableau des chaînes de référence}

\begin{center}
\begin{tabularx}{\linewidth}{|l|X|X|}\hline
\textbf{Référent} & \textbf{Substituts (occurrences successives)} & \textbf{Type de substitut} \\ \hline
Meka & Meka (1) & Nom propre (référent initial) \\
& Le vieux chef (2) & Substitut nominal (périphrase descriptive) \\
& Il (3) & Pronom personnel sujet (3e pers. sg) \\
& Sa (4) & Adjectif possessif (3e pers. sg) \\
& Ses (5) & Adjectif possessif (3e pers. pl) \\
& Le récipiendaire (6) & Substitut nominal (statut social/fonction) \\
& Son (7) & Adjectif possessif \\
& Son mari (10) & Substitut nominal (relation parentale) \\ \hline
La médaille & Sa médaille (3) & Nom + possessif (introduction) \\
& L'objet (6) & Substitut nominal (désignation générique) \\
& Ce dernier (7) & Pronom démonstratif composé (reprise proche) \\
& L'objet brillant (8) & Groupe nominal répété (insistance visuelle) \\ \hline
Kelara & Kelara (8) & Nom propre \\
& Elle (9) & Pronom personnel sujet \\
& Sa (9) & Adjectif possessif \\ \hline
\end{tabularx}
\end{center}

\textbf{2. Rapports de sens (Synonymie / Antonymie)}
\begin{itemize}
    \item \textbf{Synonymie (contextuelle) :}
    \begin{itemize}
        \item \emph{Le vieux chef} $\approx$ \emph{Le récipiendaire} $\approx$ \emph{Son mari} (désignent Meka sous des angles différents : âge/autorité, fonction cérémoniale, vie privée).
        \item \emph{Sa médaille} $\approx$ \emph{L'objet} $\approx$ \emph{Ce dernier} $\approx$ \emph{L'objet brillant} (désignent la médaille : possession, matière, proximité, aspect).
    \end{itemize}
    \item \textbf{Antonymie (implicite / Champ lexical) :}
    \begin{itemize}
        \item \emph{Brillant} (médaille) vs \emph{Usé / Poussiéreux} (implicite : le siège, les habits, les mains du vieux). Opposition \cle{Apparence / Réalité}, \cle{Valeur symbolique / Valeur d'usage}.
        \item \emph{Attendait} (passivité, temps long) vs \emph{Arrivait / Posait} (action soudaine de Kelara).
    \end{itemize}
\end{itemize}

\textbf{3. Réécriture sans substituts (effet d'alourdissement)}
« \emph{Meka} était assis sur le siège de \emph{Meka}. \emph{Meka} attendait. \emph{Meka} tournait la médaille de \emph{Meka} entre les doigts de \emph{Meka}. \emph{Meka} ne quittait pas des yeux la médaille brillante. La médaille brillante représentait tout le passé de \emph{Meka}. Soudain, \emph{Kelara} arriva. \emph{Kelara} posa la main de \emph{Kelara} sur l'épaule de \emph{Meka}. »
$\to$ Perte de fluidité, effet de répétition fastidieuse, disparition de la nuance (le « vieux chef », le « récipiendaire » apportent de l'information supplémentaire que le pronom « il » gomme).
\end{exoresolu}

\courssub{Outils de langue II : hypéronymie/hyponymie, contenus manifestes et latents, texte argumentatif}

\begin{methode}[Structurer le lexique, lire l'implicite et analyser l'argumentation]
\textbf{A. Hypéronymie et hyponymie (Organisation sémantique)}
\begin{enumerate}
    \item \textbf{Hypéronyme (terme générique, « super-ordre ») :} Mot au sens large qui englobe d'autres mots. Ex : \cle{Arme} (hypéronyme de fusil, lance, machette), \cle{Vêtement} (hypéronyme de pagne, costume, casque), \cle{Cérémonie} (hypéronyme de remise de médaille, fête, enterrement).
    \item \textbf{Hyponymes (termes spécifiques, « sous-ordre ») :} Mots au sens plus restreint, inclus dans l'hypéronyme. Relation : \emph{Hyponyme} « est un » \emph{Hypéronyme}. Ex : \emph{Le casque} est un \emph{vêtement} (ou \emph{coiffure}). \emph{La médaille} est une \emph{décoration}.
    \item \textbf{Utilité :} Construire des champs lexicaux précis, généraliser (monter en généralité) ou préciser (descendre en spécificité), définir un mot (genre + différence spécifique).
\end{enumerate}

\textbf{B. Contenus manifestes et contenus latents (Communication / Pragmatique)}
\begin{enumerate}
    \item \textbf{Contenu manifeste (explicite, dit) :} Ce que le texte énonce clairement, le sens littéral, l'information factuelle. Ex : « Meka a reçu la médaille. »
    \item \textbf{Contenu latent (implicite, non-dit, sous-entendu) :} Ce que le texte suggère, suppose, ou cache. Nécessite une inférence du lecteur (contexte, culture, ironie).
    \begin{itemize}
        \item \textbf{Presupposés :} « Le Commandant a décoré Meka » $\to$ suppose Meka méritant, Commandant légitime.
        \item \textbf{Sous-entendus (ironie) :} « Meka a bien servi la France » (manifeste : loyauté) $\to$ Latent : Meka a été exploité, aliéné, a perdu ses fils pour une cause étrangère.
        \item \textbf{Non-dits sociaux :} Le silence de Kelara pendant la cérémonie (manifeste) $\to$ Latent : réprobation, lucidité, douleur maternelle.
    \end{itemize}
    \item \textbf{Méthode :} Citer le passage manifeste $\to$ Formuler l'inférence (latent) $\to$ Justifier par des indices textuels (ton, contexte, oppositions).
\end{enumerate}

\textbf{C. Le texte argumentatif : caractéristiques et analyse}
\begin{enumerate}
    \item \textbf{Finalité :} Convaincre (raison/logos) ou Persuader (émotion/pathos) un destinataire.
    \item \textbf{Structure canonique :}
    \begin{enumerate}
        \item \textbf{Thèse (position) :} L'opinion défendue.
        \item \textbf{Arguments :} Raisons logiques (faits, statistiques, autorité, analogie) ou émotionnelles (valeurs, peur, espoir).
        \item \textbf{Exemples / Illustrations :} Preuves concrètes (références textuelles précises, citations).
        \item \textbf{Concession / Réfutation :} Prise en compte de l'adversaire (antithèse) pour mieux la renverser.
        \item \textbf{Conclusion :} Récapitulatif + ouverture.
    \end{enumerate}
    \item \textbf{Marqueurs (connecteurs logiques) :} \cle{Donc, car, parce que, en effet, or, mais, cependant, par conséquent, c'est pourquoi, premièrement, enfin}.
    \item \textbf{Modalisation :} Verbes (doit, faut, peut), adverbes (certainement, sans doute), énonciation (je pense que, il est clair que).
\end{enumerate}
\end{methode}

\begin{exoresolu}[Analyse lexicale, implicite et argumentative d'un passage clé]
\textbf{Énoncé :} Extrait (fin du roman, adapté) : « \emph{Le Commandant} épingla \emph{la croix} sur \emph{la poitrine} de \emph{Meka}. "Voilà, \emph{mon vieux}, lui dit-il, \emph{tu es décoré}. \emph{La France} te remercie." \emph{Meka} ne répondit rien. \emph{Il} sentit \emph{le fer froid} contre \emph{sa peau}. \emph{Kelara}, \emph{derrière la porte}, serra \emph{son pagne}. »
\begin{enumerate}
    \item Établir deux arborescences hypéronymie/hyponymie à partir de \emph{la croix} et \emph{le Commandant}.
    \item Identifier un contenu manifeste et le contenu latent correspondant pour les paroles du Commandant et le silence de Meka.
    \item Analyser ce passage comme un micro-texte argumentatif (thèse implicite du Commandant vs thèse latente du récit).
\end{enumerate}
\tcblower
\textbf{Solution détaillée :}

\textbf{1. Arborescences Hypéronymie / Hyponymie}

\begin{center}
\begin{tabularx}{\linewidth}{|l|X|}\hline
\textbf{Hypéronyme (Générique)} & \textbf{Hyponymes (Spécifiques - présents ou potentiels dans l'œuvre)} \\ \hline
\cle{La croix} (ou \cle{La décoration}, \cle{L'insigne}) & \emph{La médaille}, \emph{la Légion d'honneur}, \emph{le ruban}, \emph{le fer}, \emph{l'objet brillant}, \emph{la breloque} (terme dépréciatif latent). \\ \hline
\cle{L'autorité} / \cle{Le représentant du pouvoir} & \emph{Le Commandant}, \emph{l'Administrateur}, \emph{le Chef de poste}, \emph{le Blanc}, \emph{le Colon}, \emph{le Maître}. \\ \hline
\end{tabularx}
\end{center}
\textit{Note :} « La croix » est l'hyperonyme technique (forme de la médaille), « la médaille » est l'hyponyme usuel. « Le Commandant » est l'hyponyme de « L'autorité coloniale ».

\textbf{2. Contenus Manifestes vs Latents}

\begin{center}
\begin{tabularx}{\linewidth}{|l|X|X|}\hline
\textbf{Segment textuel} & \textbf{Contenu Manifeste (Explicite)} & \textbf{Contenu Latent (Implicite / Inféré)} \\ \hline
« \emph{Voilà, mon vieux, tu es décoré. La France te remercie.} » (Commandant) & Acte officiel de remise. Reconnaissance formelle de services rendus. Tutoiement « mon vieux » = familiarité paternaliste. & \textbf{Ironie tragique :} « Mon vieux » = mépris déguisé (âge, race). « La France » = entité abstraite qui ne « remercie » pas concrètement (pas de pension, pas de soins). La médaille est un leurre. \\ \hline
« \emph{Meka ne répondit rien.} » & Silence physique. Absence de discours de remerciement (rupture du protocole). & \textbf{Refus lucide / Honte / Prise de conscience :} Meka mesure le vide de l'honneur. Il ne peut pas mentir. Le silence est sa seule dignité restante. \\ \hline
« \emph{Il sentit le fer froid contre sa peau.} » & Sensation tactile thermique (froid) et matérielle (fer/métal). & \textbf{Choc de réalité :} Le métal froid = la mort, l'inhumanité du colonisateur, le contraste avec le « chaud » du sang/sacrifice. Le fer blesse (poids, froideur). \\ \hline
« \emph{Kelara, derrière la porte, serra son pagne.} » & Action gestuelle : se cacher, serrer son vêtement (pagne = vêtement traditionnel féminin). & \textbf{Protection / Rejet / Lucidité maternelle :} « Derrière la porte » = exclusion de la scène officielle, espace de la vérité. Serrer le pagne = retenir sa douleur, protéger son corps/son intimité face à l'indécence, préparer l'avenir (vendre la médaille). \\ \hline
\end{tabularx}
\end{center}

\textbf{3. Analyse argumentative du passage}
\begin{itemize}
    \item \textbf{Thèse du Commandant (Discours dominant / Propagande) :} « La collaboration paie. L'assimilation est une réussite. L'ordre colonial est juste et généreux. »
    \item \textbf{Arguments du Commandant :} Argument d'autorité (« La France »), argument de récompense (la croix), argument affectif fallacieux (« mon vieux »).
    \item \textbf{Thèse latente du récit (Contre-discours / Vision d'Oyono) :} « La collaboration est une aliénation mortifère. L'honneur colonial est une médaille de plomb (froid). La dignité est dans le silence et la femme (Kelara). »
    \item \textbf{Arguments du récit (montrés, non dits) :}
    \begin{itemize}
        \item \emph{Argument par le réel (sensoriel) :} « fer froid » vs discours chaud.
        \item \emph{Argument par le silence (éthique) :} Meka ne valide pas.
        \item \emph{Argument par le hors-champ (féminin/populaire) :} Kelara, exclue, détient la vérité (elle vendra la médaille).
    \end{itemize}
    \item \textbf{Structure :} Thèse officielle (Commandant) $\to$ Preuve par l'absurde (Silence + Froid + Exclusion) $\to$ Renversement final (Kelara récupère le métal = valeur marchande > valeur symbolique).
\end{itemize}
\end{exoresolu}

\courssub{Analyser le sujet et formuler la problématique}

\begin{methode}[Dissertation : De l'analyse du sujet à la problématique (Méthode en 5 temps)]
\textbf{Temps 1 : Décortiquer le sujet (Exégèse)}
\begin{enumerate}
    \item \textbf{Isoler les mots-clés :} Surligner les termes forts (sujet, verbes, compléments). Définir chacun (sens littéral, sens figuré, connotations).
    \item \textbf{Identifier la citation (si sujet-citation) :} Auteur, œuvre, contexte, sens immédiat dans l'œuvre.
    \item \textbf{Repérer l'instruction (opérateur) :} \cle{Commenter} (expliquer le sens/procédés), \cle{Discuter} (pour/contre, nuances), \cle{Dans quelle mesure} (évaluer une part de vérité), \cle{Qu'en pensez-vous ?} (argumentation personnelle étayée).
\end{enumerate}

\textbf{Temps 2 : Dégager l'opposition / la tension (Le « moteur » de la dissertation)}
\begin{enumerate}
    \item Chercher le \cle{paradoxe}, la \cle{contradiction apparente}, le \cle{débat} implicite.
    \item Formuler sous forme d'alternative : \emph{Thèse A} vs \emph{Thèse B}. Ex : « Honneur vs Aliénation », « Fidélité vs Servilité », « Visible vs Invisible ».
\end{enumerate}

\textbf{Temps 3 : Formuler la problématique (La question directrice)}
\begin{enumerate}
    \item Transformer la tension en \textbf{question ouverte, précise, problématique} (pas une question fermée OUI/NON).
    \item \textbf{Formules types :}
    \begin{itemize}
        \item « En quoi [Sujet] révèle-t-il [Tension] ? »
        \item « Comment [Œuvre/Auteur] articule-t-il [Thème A] et [Thème B] ? »
        \item « Dans quelle mesure [Citation] illustre-t-elle la complexité de [Concept] chez [Auteur] ? »
        \item « [Sujet] : [Thèse] ou [Antithèse] ? » (À éviter si trop binaire, préférer « Comment concilier... »).
    \end{itemize}
    \item \textbf{Test de validation :} La problématique appelle-t-elle un plan en 2 ou 3 parties ? Engage-t-elle l'œuvre étudiée (\emph{Le Vieux nègre...}) ? Est-elle débattable ?
\end{enumerate}

\textbf{Temps 4 : Délimiter le champ d'étude}
Préciser le corpus (ici : \emph{Le Vieux nègre et la médaille} + éventuellement culture générale littéraire africaine/française) et les concepts clés (aliénation, assimilation, satire, dignité, femme, mort).

\textbf{Temps 5 : Rédiger l'introduction (Schéma entonnoir)}
\cle{Amorce} (général, citation, actualité) $\to$ \cle{Présentation sujet/œuvre} $\to$ \cle{Définition termes} $\to$ \cle{Problématique} $\to$ \cle{Annonce du plan} (sommaire : « Nous verrons d'abord... puis... enfin... »).
\end{methode}

\begin{exoresolu}[Analyse de sujet et formulation de problématique]
\textbf{Énoncé :} Sujet de dissertation type Probatoire : 
« \textbf{« Le Vieux nègre et la médaille est une satire de l'assimilation. » Discutez cette affirmation.} »
Appliquez la méthode en 5 temps pour produire : 1. L'exégèse, 2. La tension, 3. La problématique validée, 4. L'annonce de plan.
\tcblower
\textbf{Solution détaillée :}

\textbf{1. Exégèse du sujet}
\begin{itemize}
    \item \textbf{Mots-clés :}
    \begin{itemize}
        \item \cle{Le Vieux nègre et la médaille} : Œuvre de référence unique (roman d'Oyono, 1956). Titre = programme.
        \item \cle{Satire} : Genre littéraire (ton critique, ironie, caricature, rire correcteur). Vise à corriger les mœurs en les ridiculisant. Cibles : les colonisés assimilés, l'administration coloniale, l'Église complice.
        \item \cle{Assimilation} : Politique coloniale française (devenir français par la culture, la langue, la loi, l'armée). Idéal : « Français d'outre-mer ». Réalité : aliénation, perte d'identité, inégalité structurelle.
        \item \cle{Est une satire} : Affirmation forte (thèse imposée par le sujet). Le sujet demande de \cle{Discuter} = ne pas dire seulement « oui », mais nuancer, étendre, contredire partiellement.
    \end{itemize}
    \item \textbf{Instruction :} \textbf{Discuter} $\to$ Examiner les arguments \textbf{POUR} (c'est une satire), les arguments \textbf{CONTRE / NUANCES} (c'est aussi une tragédie, un réalisme social, une critique de la condition féminine, une métaphore de la mort), et \textbf{SYNTHÉTISER}.
\end{itemize}

\textbf{2. Dégager la tension (Le débat)}
\begin{itemize}
    \item \textbf{Thèse (Affirmation du sujet) :} L'œuvre est une \emph{comédie satirique}. Le rire (costume ridicule, discours creux, médaille clinquante) dénonce l'absurdité de l'assimilation. Meka est une marionnette grotesque.
    \item \textbf{Antithèse (Réalité du texte) :} L'œuvre dépasse la satire. C'est une \emph{tragédie humaine} (mort de Meka, sacrifice des fils, pauvreté de Kelara). L'assimilation n'est pas seulement ridicule, elle est \emph{mortifère}. Le rire se fige en horreur. La satire cache une douleur réelle.
    \item \textbf{Tension centrale :} \cle{Rire / Horreur} — \cle{Forme satirique / Fond tragique} — \cle{Dénonciation politique / Empathie humaine}.
\end{itemize}

\textbf{3. Formulation de la problématique (Validée)}
\begin{quote}
\textbf{« En quoi \emph{Le Vieux nègre et la médaille}, sous les apparences de la satire de l'assimilation, révèle-t-elle la dimension tragique d'une aliénation qui coûte la vie et l'identité ? »}
\end{quote}
\textit{Validation :} Question ouverte (« En quoi »), lie les deux pôles (satire / tragique, assimilation / aliénation/mort), engage l'analyse textuelle précise, permet un plan dialectique en 3 parties.

\textbf{4. Annonce de plan (Sommaire pour l'introduction)}
\begin{quote}
« Nous analyserons dans un premier temps comment l'ironie et la caricature font de l'œuvre une satire féroce de l'assimilation (I). Nous montrerons ensuite que cette satire laisse place à une tragédie humaine où l'assimilation se révèle mortifère (II). Enfin, nous verrons que l'œuvre dépasse le cadre strictement politique pour interroger la dignité africaine face à l'Histoire, notamment par le personnage de Kelara (III). »
\end{quote}
\end{exoresolu}

\courssub{Rechercher et organiser les idées}

\begin{methode}[Recherche, sélection et hiérarchisation des idées (Méthode du « Mind-Mapping » structuré)]
\textbf{Étape 1 : Le Brainstorming / Mind-Mapping (Divergence)}
\begin{enumerate}
    \item Au centre : \textbf{Problématique} (ex: \emph{Satire vs Tragique / Aliénation mortifère}).
    \item Branches principales = \textbf{Grosses parties potentielles} (I, II, III).
    \item Sous-branches = \textbf{Arguments} (idées forces).
    \item Feuilles = \textbf{Preuves textuelles} (citations précises, numéros de page/chapitre, scènes précises, procédés stylistiques).
    \item \textbf{Consigne :} Ne pas se censurer. Tout noter : thèmes, personnages, style, contexte, critiques connus.
\end{enumerate}

\textbf{Étape 2 : La Sélection / Tri (Convergence)}
\begin{enumerate}
    \item \textbf{Garder :} Idées directement liées à la problématique, étayables par le texte, originales.
    \item \textbf{Éliminer :} Hors-sujet (biographie d'Oyono sans lien), généralités non prouvées, répétitions, idées trop faibles.
    \item \textbf{Regrouper :} Fusionner les idées proches sous une même étiquette d'argument.
\end{enumerate}

\textbf{Étape 3 : L'Organisation en Plan Dialectique (Structure standard Bac)}
\begin{itemize}
    \item \textbf{Plan en 2 parties (Thèse / Antithèse) :} Rare au Bac, souvent trop binaire. Réservé aux sujets « Pour ou Contre » simples.
    \item \textbf{Plan en 3 parties (Thèse / Antithèse / Synthèse ou Dépassement) :} \textbf{Standard attendu.}
    \begin{itemize}
        \item \textbf{I. Thèse :} Réponse « Oui, mais... » à la problématique (valide l'affirmation du sujet). Ex : \emph{Les mécanismes de la satire.}
        \item \textbf{II. Antithèse / Nuance :} Réponse « Non, c'est plus complexe » (limites de la thèse). Ex : \emph{La tragédie sous la satire.}
        \item \textbf{III. Synthèse / Dépassement :} Réponse « En réalité, l'œuvre fait... » (vision unifiée, portée universelle). Ex : \emph{La reconquête de la dignité / L'écriture comme résistance.}
    \end{itemize}
\end{itemize}

\textbf{Étape 4 : Le Détail des Sous-parties (Niveau 2 : A, B / 1, 2)}
Chaque grande partie (I, II, III) = 2 ou 3 sous-parties (A, B ou 1, 2).
\begin{itemize}
    \item \textbf{Règle d'or :} 1 Sous-partie = 1 Argument principal + 1 ou 2 Preuves textuelles (Citations) + 1 Analyse (Commentaire de la citation).
    \item \textbf{Équilibre :} Même nombre de sous-parties par grande partie (ex: 2/2/2 ou 3/3/3). Même volume approximatif.
\end{itemize}

\textbf{Étape 5 : Les Transitions}
Préparer les phrases de transition entre I $\to$ II et II $\to$ III. Elles doivent reprendre l'idée finale de la partie précédente et annoncer le renversement ou l'approfondissement de la suivante.
\end{methode}

\begin{exoresolu}[Recherche et organisation des idées pour le plan]
\textbf{Énoncé :} À partir de la problématique validée : \emph{« En quoi Le Vieux nègre et la médaille, sous les apparences de la satire de l'assimilation, révèle-t-elle la dimension tragique d'une aliénation qui coûte la vie et l'identité ? »}, effectuez le brainstorming guidé, la sélection et proposez l'architecture détaillée des 3 parties avec arguments et citations précises (références chapitres/scènes).
\tcblower
\textbf{Solution détaillée :}

\textbf{1. Brainstorming guidé (Extrait synthétique)}

\begin{center}
\begin{tabularx}{\linewidth}{|l|X|X|}\hline
\textbf{Partie Potentielle} & \textbf{Arguments (Idées forces)} & \textbf{Preuves Textuelles (Citations / Scènes)} \\ \hline
\textbf{I. Satire de l'assimilation} & A. Caricature du « évolué » / collaborateur. & Costume trop grand (Ch. 1), discours stéréotypés du Commandant (Ch. 5), médaille « breloque ». \\
& B. Ironie du rituel colonial (théâtre grotesque). & Cérémonie : fanfare qui rate, prêtre complice, foule manipulée (Ch. 5-6). \\
& C. Dénonciation de l'Église (assimilation spirituelle). & Prêtre : « Meka est un bon chrétien » (Ch. 3), opposition Kelara/Prêtre. \\ \hline
\textbf{II. Tragédie de l'aliénation} & A. Sacrifices vains (fils, terre, jeunesse). & Fils aîné mort à Bir Hakeim (Ch. 2), fils cadet prêtre « blanc » (Kelara : « Il n'est plus mon fils »), terres données. \\
& B. Le corps usé / La mort comme seule issue. & Crise cardiaque finale (Ch. 6), « fer froid », fatigue, vieillesse prématurée. \\
& C. Perte d'identité / Vide de la récompense. & Meka ne sait pas lire la citation, silence final, médaille vendue par Kelara. \\ \hline
\textbf{III. Dépassement : Résistance / Dignité} & A. Kelara : conscience lucide, garde-fou. & Refus de la messe (Ch. 1), « On a mangé ton argent », vente finale de la médaille (survie). \\
& B. L'écriture d'Oyono : langue française détournée (oralité, proverbes). & Style : « Le blanc est comme le fromage... », proverbes ewondo, ironie narrative. \\
& C. Portée universelle / Historique. & Allégorie du Cameroun/Indépendance : la médaille = drapeau ? La mort = fin colonisation ? \\ \hline
\end{tabularx}
\end{center}

\textbf{2. Sélection et Architecture du Plan Détaillé (Niveau 1 et 2)}

\textbf{I. Une satire féroce des mécanismes de l'assimilation coloniale}
\begin{itemize}
    \item \textbf{A. La caricature du « bon indigène » : Meka, pantin décoré.}
    \begin{itemize}
        \item Argument : L'assimilation fabrique un clone ridicule, vidé de sa substance.
        \item Preuves : Costume « trop grand » (Ch.1) $\to$ Métaphore visuelle de l'inadéquation. Discours du Commandant : « mon vieux » (Ch.5) $\to$ Paternalisme méprisant.
    \end{itemize}
    \item \textbf{B. Le théâtre de la cérémonie : une mascarade institutionnalisée.}
    \begin{itemize}
        \item Argument : Le pouvoir colonial met en scène sa propre légitimité par le ridicule.
        \item Preuves : Fanfare « qui souffle dans le vide » (Ch.5), discours creux du Prêtre (complicité), foule acquise à la cause par la bière/riz.
    \end{itemize}
    \item \textbf{C. La langue comme instrument de domination (et de satire).}
    \begin{itemize}
        \item Argument : Le français imposé devient langue de bois.
        \item Preuves : Citations latines du Prêtre incomprises, formulaires administratifs, opposition français / ewondo (proverbes de Kelara).
    \end{itemize}
\end{itemize}

\textbf{II. Le basculement dans le tragique : l'assimilation comme aliénation mortifère}
\begin{itemize}
    \item \textbf{A. Le prix du sang : une famille sacrifiée sur l'autel de la France.}
    \begin{itemize}
        \item Argument : L'assimilation exige le don total (vie des fils, terre, santé) pour un retour nul.
        \item Preuves : Fils aîné « mort pour la France » (Bir Hakeim, Ch.2), Fils cadet (prêtre) coupé de sa mère « Il n'est plus mon fils » (Kelara, Ch.3), Meka « usé » avant l'âge.
    \end{itemize}
    \item \textbf{B. Le corps du vieux nègre : site de l'exploitation et de la mort.}
    \begin{itemize}
        \item Argument : La médaille pèse sur le corps comme une condamnation.
        \item Preuves : Sensation « fer froid contre sa peau » (Ch.6), crise cardiaque immédiate après la remise, « Meka est mort » (dernière phrase). La médaille ne prolonge pas la vie, elle la termine.
    \end{itemize}
    \item \textbf{C. Le vide sémantique de la récompense.}
    \begin{itemize}
        \item Argument : L'objet (médaille) ne signifie rien pour le récipiendaire (analphabète, culture orale).
        \item Preuves : Meka ne lit pas la citation, Kelara : « Ça ne mange pas » (Ch.6), vente finale pour acheter du riz $\to$ Valeur d'échange > Valeur symbolique.
    \end{itemize}
\end{itemize}

\textbf{III. Au-delà de la satire et du tragique : l'affirmation d'une dignité résistante}
\begin{itemize}
    \item \textbf{A. Kelara, figure de la lucidité et de la survie (contre-pouvoir féminin).}
    \begin{itemize}
        \item Argument : Elle incarne le refus lucide, l'ancrage dans le réel (nourrir, soigner), l'avenir.
        \item Preuves : Refus de la messe d'action de grâce (Ch.1), critique acerbe « On a mangé ton argent », récupération finale de la médaille (Ch.6) $\to$ Pragmatisme vital.
    \end{itemize}
    \item \textbf{B. L'écriture d'Oyono : une « satire armée » (détournement de la langue du colon).}
    \begin{itemize}
        \item Argument : L'auteur utilise la langue française pour la retourner contre l'idéologie coloniale.
        \item Preuves : Ironie narrative (narrateur > personnages), proverbes ewondo intraduisibles (résistance culturelle), style oralisé (rythme, répétitions).
    \end{itemize}
    \item \textbf{C. Une allégorie de l'Histoire : la fin d'un cycle.}
    \begin{itemize}
        \item Argument : La mort de Meka = fin de l'ordre colonial ; Kelara = avenir africain.
        \item Preuves : Date publication (1956), contexte loi-cadre, « Le Vieux nègre » = figure du colonisé historique. La médaille vendue = récupération du symbole par le peuple.
    \end{itemize}
\end{itemize}

\textbf{3. Transitions préparées}
\begin{itemize}
    \item \textbf{I $\to$ II :} « Mais si le rire de la satire dénonce le ridicule du système, il ne doit pas masquer l'horreur du prix payé par les corps et les âmes. Derrière le faste de la cérémonie se profile la tragédie intime d'une vie donnée pour rien. »
    \item \textbf{II $\to$ III :} « Pourtant, l'œuvre ne s'arrête pas à la mort de Meka. Par la figure de Kelara et par l'écriture elle-même, Oyono ouvre une brèche : la dignité ne meurt pas avec le vieux chef, elle se déplace vers celle qui sait que « ça ne mange pas » une médaille. »
\end{itemize}
\end{exoresolu}

\courssub{Élaborer et produire le plan détaillé}

\begin{methode}[Rédiger le plan détaillé : Normes de présentation et exigences du Bac]
\textbf{Rappel :} Le plan détaillé est \textbf{obligatoire} en brouillon et souvent demandé en début de copie (ou évalué implicitement par la structure). Il se présente sous forme \textbf{hiérarchisée et indentée}.

\textbf{1. Structure formelle (Modèle canonique)}
\begin{flushleft}
\texttt{INTRODUCTION}\\
\texttt{\ \ [Amorce] ...}\\
\texttt{\ \ [Présentation sujet/œuvre] ...}\\
\texttt{\ \ [Définitions] ...}\\
\texttt{\ \ [Problématique] ...}\\
\texttt{\ \ [Annonce de plan] ...}\\
\texttt{DÉVELOPPEMENT}\\
\texttt{I. [Titre Partie I - Majuscules, nominale ou verbale]}\\
\texttt{\ \ A. [Titre Sous-partie A]}\\
\texttt{\ \ \ \ 1. Argument principal (Idée force)}\\
\texttt{\ \ \ \ 2. Preuve textuelle (Citation précise + Référence)}\\
\texttt{\ \ \ \ 3. Analyse / Commentaire (Lien argument-preuve-problématique)}\\
\texttt{\ \ B. [Titre Sous-partie B]}\\
\texttt{\ \ \ \ 1. Argument ...}\\
\texttt{\ \ \ \ 2. Preuve ...}\\
\texttt{\ \ \ \ 3. Analyse ...}\\
\texttt{\ \ [Transition I $\to$ II rédigée]}\\
\texttt{II. [Titre Partie II]}\\
\texttt{\ \ A. [Titre Sous-partie A]}\\
\texttt{\ \ \ \ 1. Argument ...}\\
\texttt{\ \ \ \ 2. Preuve ...}\\
\texttt{\ \ \ \ 3. Analyse ...}\\
\texttt{\ \ B. [Titre Sous-partie B]}\\
\texttt{\ \ \ \ 1. Argument ...}\\
\texttt{\ \ \ \ 2. Preuve ...}\\
\texttt{\ \ \ \ 3. Analyse ...}\\
\texttt{\ \ [Transition II $\to$ III rédigée]}\\
\texttt{III. [Titre Partie III - Synthèse / Dépassement]}\\
\texttt{\ \ A. [Titre Sous-partie A]}\\
\texttt{\ \ \ \ 1. Argument ...}\\
\texttt{\ \ \ \ 2. Preuve ...}\\
\texttt{\ \ \ \ 3. Analyse ...}\\
\texttt{\ \ B. [Titre Sous-partie B]}\\
\texttt{\ \ \ \ 1. Argument ...}\\
\texttt{\ \ \ \ 2. Preuve ...}\\
\texttt{\ \ \ \ 3. Analyse ...}\\
\texttt{CONCLUSION}\\
\texttt{\ \ [Bilan / Réponse à la problématique] ...}\\
\texttt{\ \ [Ouverture] ...}
\end{flushleft}

\textbf{2. Règles d'or de rédaction}
\begin{enumerate}
    \item \textbf{Titres (I, II, III, A, B) :} Phrases nominales (pas de verbe conjugué), concises, parlantes, \textbf{parallèles} (même structure syntaxique si possible). Pas de numérotation 1.1.1 (réservé au corps du devoir).
    \item \textbf{Niveau 3 (1, 2) :} Phrases verbales courtes (sujet + verbe + complément). Écrire l'\textbf{argument} (pas juste un mot-clé). Intégrer la \textbf{citation} (ou référence exacte : Ch., p.) dans le point 2.
    \item \textbf{Citations :} Courtes (max 1 vers ou 1 segment de phrase), intégrées dans la phrase ou isolées entre guillemets. Indiquer la référence (Chapitre / Page édition de référence).
    \item \textbf{Transitions :} Rédigées \textbf{en toutes lettres} dans le plan détaillé (elles seront recopiées en début de partie II et III dans la copie finale).
    \item \textbf{Introduction / Conclusion :} Noter les \textbf{mots-clés} de chaque étape (Amorce, Problématique, Annonce plan / Bilan, Ouverture). Ne pas rédiger le texte intégral au brouillon (gain de temps), mais avoir les formules exactes prêtes.
\end{enumerate}

\textbf{3. Check-list de validation du plan (Auto-évaluation avant rédaction)}
\begin{itemize}
    \item[\textbullet] Le plan répond-il \textbf{exactement} à la problématique ?
    \item[\textbullet] Y a-t-il une \textbf{progression logique} (pas de répétition I/II, III va plus loin) ?
    \item[\textbullet] L'\textbf{équilibre} est-il respecté (2 ou 3 sous-parties par partie, volume égal) ?
    \item[\textbullet] Les \textbf{preuves} sont-elles \textbf{précises, variées} (début, milieu, fin, narration, dialogue, description) et \textbf{analysables} ?
    \item[\textbullet] Les \textbf{titres} sont-ils \textbf{clairs} et \textbf{accrocheurs} ?
    \item[\textbullet] Les \textbf{transitions} font-elles le lien logique ?
\end{itemize}
\end{methode}

\begin{exoresolu}[Production du plan détaillé complet (Prêt pour l'examen)]
\textbf{Énoncé :} Produisez le \textbf{plan détaillé final, prêt à être utilisé en examen}, pour le sujet : \emph{« "Le Vieux nègre et la médaille est une satire de l'assimilation." Discutez cette affirmation. »} avec la problématique : \emph{« En quoi Le Vieux nègre et la médaille, sous les apparences de la satire de l'assimilation, révèle-t-elle la dimension tragique d'une aliénation qui coûte la vie et l'identité ? »}. Respectez scrupuleusement la mise en forme, les citations intégrées et les transitions rédigées.
\tcblower
\textbf{Solution détaillée :}

\textbf{PLAN DÉTAILLÉ — DISSERTATION LITTÉRAIRE}

\textbf{INTRODUCTION}
\begin{itemize}
    \item \textbf{Amorce :} Citation de Frantz Fanon (\emph{Peau noire, masques blancs}) : « Le Nègre n'est pas. Le Nègre devient. » $\to$ L'assimilation comme aliénation.
    \item \textbf{Présentation :} \emph{Le Vieux nègre et la médaille}, roman de Ferdinand Oyono (1956), Cameroun. Contexte : fin colonisation, loi-cadre Defferre.
    \item \textbf{Définitions :} Satire (rire correcteur, caricature) ; Assimilation (politique d'absorption culturelle/francisation).
    \item \textbf{Problématique :} \textbf{En quoi \emph{Le Vieux nègre et la médaille}, sous les apparences de la satire de l'assimilation, révèle-t-elle la dimension tragique d'une aliénation qui coûte la vie et l'identité ?}
    \item \textbf{Annonce de plan :} Nous analyserons d'abord les mécanismes satiriques qui dénoncent l'absurdité de l'assimilation (I). Nous montrerons ensuite que la satire laisse place à une tragédie où l'assimilation se révèle mortifère (II). Enfin, nous verrons comment l'œuvre dépasse le constat pour affirmer une dignité résistante par l'écriture et le féminin (III).
\end{itemize}

\textbf{DÉVELOPPEMENT}

\textbf{I. Une satire féroce des mécanismes de l'assimilation coloniale}
\begin{itemize}
    \item \textbf{A. La caricature du « bon indigène » : Meka, pantin décoré.}
    \begin{enumerate}
        \item L'assimilation fabrique un clone ridicule, vidé de sa substance africaine.
        \item Preuve : « \emph{Son costume était trop grand pour lui} » (Ch. 1) / « \emph{Voilà, mon vieux, tu es décoré} » (Commandant, Ch. 5).
        \item Analyse : L'hyperbole vestimentaire (costume/casque) figure l'inadéquation ontologique ; le tutoiement paternaliste (« mon vieux ») masque le mépris racial.
    \end{enumerate}
    \item \textbf{B. Le théâtre de la cérémonie : une mascarade institutionnalisée.}
    \begin{enumerate}
        \item Le pouvoir colonial met en scène sa propre légitimité par le grotesque.
        \item Preuve : « \emph{La fanfare soufflait dans le vide} » (Ch. 5) ; Discours du Prêtre : « \emph{Meka est un bon chrétien} » (Ch. 5).
        \item Analyse : La fanfare incompétente = administration coloniale déconnectée ; le prêtre valide l'aliénation spirituelle (complicité Église/État).
    \end{enumerate}
    \item \textbf{C. La langue française détournée en langue de bois.}
    \begin{enumerate}
        \item Le français imposé devient instrument de vide sémantique.
        \item Preuve : Citations latines du Prêtre non traduites (Ch. 5) ; Formulaires administratifs incompris de Meka.
        \item Analyse : La langue du colonisateur (hypéronyme « Langue de civilisation ») se réduit à un rituel magique vide de sens pour le sujet colonisé.
    \end{enumerate}
\end{itemize}

\textbf{\textit{Transition I $\to$ II :} Mais si le rire de la satire dénonce le ridicule du système, il ne doit pas masquer l'horreur du prix payé par les corps et les âmes. Derrière le faste de la cérémonie se profile la tragédie intime d'une vie donnée pour rien.}

\textbf{II. Le basculement dans le tragique : l'assimilation comme aliénation mortifère}
\begin{itemize}
    \item \textbf{A. Le prix du sang : une famille sacrifiée sur l'autel de la France.}
    \begin{enumerate}
        \item L'assimilation exige le don total (vie des fils, terre, santé) pour un retour nul.
        \item Preuve : Fils aîné « \emph{mort pour la France à Bir Hakeim} » (Ch. 2) ; Kelara : « \emph{Il n'est plus mon fils} » (à propos du cadet prêtre, Ch. 3).
        \item Analyse : L'antonymie « Mort pour la France / Vivant pour rien » ; la filiation brisée (prêtre « blanc ») = mort symbolique de la lignée.
    \end{enumerate}
    \item \textbf{B. Le corps du vieux nègre : site de l'exploitation et de la mort.}
    \begin{enumerate}
        \item La médaille pèse sur le corps comme une condamnation physique.
        \item Preuve : « \emph{Il sentit le fer froid contre sa peau} » (Ch. 6) ; « \emph{Meka est mort} » (Dernière phrase).
        \item Analyse : Oxymore « fer froid » (métal inerte / vie chaude) ; la remise de médaille = déclencheur de la mort (causalité ironique). Le corps usé est le vrai texte de l'Histoire.
    \end{enumerate}
    \item \textbf{C. Le vide sémantique de la récompense : l'objet vs la vie.}
    \begin{enumerate}
        \item L'objet médaille ne signifie rien pour le récipiendaire (analphabète, culture orale/utile).
        \item Preuve : Meka ne lit pas la citation (Ch. 5) ; Kelara : « \emph{Cela ne mange pas} » (Ch. 6) ; Vente finale de la médaille.
        \item Analyse : Antonymie \emph{Valeur symbolique coloniale} / \emph{Valeur d'usage vitale}. La vente par Kelara = réappropriation matérielle du symbole.
    \end{enumerate}
\end{itemize}

\textbf{\textit{Transition II $\to$ III :} Pourtant, l'œuvre ne s'arrête pas à la mort de Meka. Par la figure de Kelara et par l'écriture elle-même, Oyono ouvre une brèche : la dignité ne meurt pas avec le vieux chef, elle se déplace vers celle qui sait que « ça ne mange pas » une médaille.}

\textbf{III. Au-delà de la satire et du tragique : l'affirmation d'une dignité résistante}
\begin{itemize}
    \item \textbf{A. Kelara, figure de la lucidité et de la survie (contre-pouvoir féminin).}
    \begin{enumerate}
        \item Elle incarne le refus lucide, l'ancrage dans le réel (nourrir, soigner), l'avenir.
        \item Preuve : Refus de la messe d'action de grâce (Ch. 1) ; « \emph{On a mangé ton argent} » (Ch. 1) ; Récupération finale : « \emph{Kelara prit la médaille... pour aller la vendre} » (Ch. 6).
        \item Analyse : Kelara = hypéronyme « Mère Afrique » / « Peuple ». Son pragmatisme (vente) > Symbolique coloniale. Elle « lit » la vérité là où Meka subissait le symbole.
    \end{enumerate}
    \item \textbf{B. L'écriture d'Oyono : une « satire armée » (détournement de la langue du colon).}
    \begin{enumerate}
        \item L'auteur utilise la langue française pour la retourner contre l'idéologie coloniale.
        \item Preuve : Ironie narrative (narrateur omniscient jugeant le Commandant) ; Proverbes ewondo : « \emph{Le blanc est comme le fromage, il a une croûte dure et un cœur mou} » (Ch. 2) ; Style oralisé (rythme binaire, répétitions).
        \item Analyse : L'hyponymie « Proverbe / Sagacité » vs « Discours officiel / Mensonge ». L'écriture devient acte de résistance (décolonisation du langage).
    \end{enumerate}
    \item \textbf{C. Une allégorie de l'Histoire : la fin d'un cycle colonial.}
    \begin{enumerate}
        \item La mort de Meka = fin de l'ordre colonial ; Kelara = avenir africain.
        \item Preuve : Date 1956 (Loi-cadre, autonomie) ; « \emph{Le Vieux nègre} » (figure allégorique du colonisé historique).
        \item Analyse : Le roman prophétise l'indépendance : la médaille (symbole colonial) redevient métal brut (matière première) entre les mains de Kelara (peuple souverain).
    \end{enumerate}
\end{itemize}

\textbf{CONCLUSION}
\begin{itemize}
    \item \textbf{Bilan :} L'œuvre n'est pas \emph{seulement} une satire. Elle utilise la satire (I) pour mieux révéler la tragédie (II) et fonder une espérance résistante (III). L'assimilation est démasquée comme crime contre l'humanité (aliénation, mort).
    \item \textbf{Réponse à la problématique :} La satire est le \emph{moyen} (forme), la tragédie le \emph{contenu} (réalité), la résistance la \emph{finalité} (avenir).
    \item \textbf{Ouverture :} Comparer avec \emph{Une vie de boy} (Oyono) ou \emph{L'Aventure ambiguë} (Kane) : l'assimilation comme déchirure identitaire universelle. Actualité : question des restitutions d'objets/symboles coloniaux (médiations muséales vs besoins vitaux des peuples).
\end{itemize}
\end{exoresolu}

\courssub{Synthèse transversale : Du texte à la dissertation (Parcours intégré)}

\begin{methode}[Parcours intégré : De la lecture de l'œuvre à la copie finale (Récapitulatif opératoire)]
Cette fiche résume l'enchaînement logique des 6 séquences du programme pour le jour de l'épreuve.
\begin{enumerate}
    \item \textbf{J1-J2 (Réception) :} \textbf{Paratexte + Lecture ouvroir.} $\to$ Fiche identité + Résumé + Hypothèses thèmes. \textit{Objectif : Maîtriser l'intrigue, les personnages, le contexte.}
    \item \textbf{J3-J4 (Langue I) :} \textbf{Référents/Substituts + Synonymie/Antonymie.} $\to$ Exercices sur extraits. \textit{Objectif : Outils pour l'analyse stylistique (cohérence, champs lexicaux) et la contraction.}
    \item \textbf{J5-J6 (Langue II) :} \textbf{Hypéronymie/Hyponymie + Manifestes/Latents + Argumentatif.} $\to$ Analyse de passages clés (ironie, implicite). \textit{Objectif : Lire en profondeur (sous-texte), repérer la thèse narrative.}
    \item \textbf{J7-J8 (Production I) :} \textbf{Analyse sujet + Problématique.} $\to$ Entraînement sur 3 sujets types. \textit{Objectif : Automatiser la méthode en 5 temps (Exégèse $\to$ Tension $\to$ Problématique validée).}
    \item \textbf{J9-J10 (Production II) :} \textbf{Recherche idées + Organisation (Mind-map $\to$ Plan 3 parties).} $\to$ Fiches « Banque d'arguments/Preuves » par thème (Satire, Tragique, Femme, Langue, Histoire). \textit{Objectif : Avoir du « stock » textuel prêt à l'emploi.}
    \item \textbf{J11-J12 (Production III) :} \textbf{Plan détaillé + Transitions + Introduction/Conclusion types.} $\to$ Rédaction de 2 plans détaillés complets en conditions chronométrées (1h30). \textit{Objectif : Maîtriser la norme formelle (indentation, citations, transitions rédigées).}
    \item \textbf{J13 (Remédiation) :} \textbf{Correction type Bac.} $\to$ Grille d'auto-évaluation (Hors-sujet ? Plan binaire ? Preuves manquantes ? Langue ?). \textit{Objectif : Corriger ses erreurs récurrentes.}
    \item \textbf{J14 (Évaluation) :} \textbf{Devoir sur table type Probatoire (4h).} Sujet inédit. Application intégrale.
\end{enumerate}
\textbf{Reflexe « Bac Blanc » :} Chronométrage type : 1h Analyse sujet/Problématique/Plan détaillé $\to$ 2h30 Rédaction $\to$ 30min Relecture (Orthographe, Citations exactes, Structure, Transitions).
\end{methode}

\begin{exoresolu}[Simulation d'examen : Sujet inédit - Plan détaillé express]
\textbf{Énoncé :} Sujet Probatoire 2025 (Simulé) : \emph{« Dans Le Vieux nègre et la médaille, la parole des humbles (Kelara, les villageois) déconstruit-elle le discours des puissants (Commandant, Prêtre, Meka) ? »}
Vous disposez de 45 minutes pour produire : 1. L'exégèse rapide, 2. La problématique, 3. Le plan détaillé (Titres I, II, III, A, B + 1 argument + 1 citation par sous-partie + Transitions).
\tcblower
\textbf{Solution détaillée (Modèle de copie « Brouillon propre » - 45 min)}

\textbf{1. EXÉGÈSE RAPIDE (5 min)}
\begin{itemize}
    \item \textbf{Mots-clés :} \cle{Parole des humbles} (Kelara, peuple, oralité, proverbes, silence, actes) vs \cle{Discours des puissants} (Commandant/Administration, Prêtre/Église, Meka/Collaborateur — \textit{nuance : Meka est puissant local mais dominé global}). \cle{Déconstruire} = démonter, montrer le mensonge, opposer une vérité contraire.
    \item \textbf{Instruction :} Question fermée « Déconstruit-elle ? » $\to$ Transformer en « Comment / En quoi / Dans quelle mesure » pour dissertation.
    \item \textbf{Tension :} Discours officiel (français, écrit, solennel, mensonger, paternaliste) vs Parole vraie (ewondo/argot, orale, prosaïque, lucide, résistante). Meka est le point de bascule : il parle le discours du puissant mais son corps/silence rejoignent les humbles.
\end{itemize}

\textbf{2. PROBLÉMATIQUE (3 min)}
\begin{quote}
\textbf{« Comment la contre-parole des humbles (Kelara, villageois), par son ancrage dans le réel et l'oralité, déconstruit-elle l'illusion du discours dominant (colonial, religieux, collaborateur) pour restaurer une vérité africaine ? »}
\end{quote}

\textbf{3. PLAN DÉTAILLÉ (37 min)}

\textbf{INTRODUCTION} : Amorce (Proverbe ewondo : « La parole est un œuf... ») $\to$ Présentation œuvre/contexte $\to$ Définitions (Humbles/Puissants, Déconstruire) $\to$ Problématique $\to$ Annonce : I. Discours des puissants : théâtre de l'illusion. II. Parole des humbles : lucidité du réel. III. Le basculement : la parole vraie l'emporte par l'écriture.

\textbf{DÉVELOPPEMENT}

\textbf{I. Le discours des puissants : une parole vide, théâtrale et aliénante}
\begin{itemize}
    \item \textbf{A. Le Commandant : la langue coloniale, instrument de domination symbolique.}
    \begin{enumerate}
        \item Argument : Discours stéréotypé, paternaliste, vide de sens réel pour l'auditoire.
        \item Citation : « \emph{Voilà, mon vieux, tu es décoré. La France te remercie.} » (Ch. 5).
    \end{enumerate}
    \item \textbf{B. Le Prêtre : la parole religieuse détournée en caution de l'ordre colonial.}
    \begin{enumerate}
        \item Argument : Latin incompris, bénédiction de l'injustice, complicité active.
        \item Citation : « \emph{Meka est un bon chrétien...} » (Ch. 5) + opposition avec Kelara (Ch. 1).
    \end{enumerate}
    \item \textbf{C. Meka : l'aliénation par la parole empruntée (le collaborateur muet).}
    \begin{enumerate}
        \item Argument : Meka intériorise le discours du maître, perd sa propre voix (silence final).
        \item Citation : « \emph{Meka ne répondit rien} » (Ch. 6) — Silence = échec de la parole assimilée.
    \end{enumerate}
\end{itemize}

\textit{Transition I $\to$ II :} Face à cette parole officielle qui masque le néant, s'élève une autre parole, celle du quotidien, du corps et de la mère, qui refuse le mensonge.

\textbf{II. La parole des humbles : lucidité, oralité et résistance concrète}
\begin{itemize}
    \item \textbf{A. Kelara : la parole vraie, pragmatique et prophétique.}
    \begin{enumerate}
        \item Argument : Elle nomme les choses (argent mangé, médaille inutile), refuse le rituel, assure la survie.
        \item Citation : « \emph{On a mangé ton argent} » (Ch. 1) / « \emph{Cela ne mange pas} » (Ch. 6) / Vente finale.
    \end{enumerate}
    \item \textbf{B. Le peuple (villageois) : l'oralité comme contre-pouvoir (chants, rires, proverbes).}
    \begin{enumerate}
        \item Argument : La foule n'est pas dupe ; chants ironiques, proverbes ewondo intraduisibles = espace de liberté.
        \item Citation : Proverbe : « \emph{Le blanc est comme le fromage...} » (Ch. 2) ; Rires pendant la cérémonie (Ch. 5).
    \end{enumerate}
    \item \textbf{C. Le corps et le silence : parole non-verbale de la vérité.}
    \begin{enumerate}
        \item Argument : Le corps de Meka (fer froid, mort) et le geste de Kelara (serrer pagne) disent l'horreur que les mots officiels taisent.
        \item Citation : « \emph{Il sentit le fer froid contre sa peau} » (Ch. 6) ; « \emph{Kelara, derrière la porte, serra son pagne} ».
    \end{enumerate}
\end{itemize}

\textit{Transition II $\to$ III :} Cette déconstruction par les humbles ne reste pas au stade de la critique ; elle s'incarne dans l'écriture même du roman, qui leur donne le dernier mot.

\textbf{III. L'écriture d'Oyono : la déconstruction achevée par la littérature (Le dernier mot aux humbles)}
\begin{itemize}
    \item \textbf{A. Le récit comme réappropriation de la parole (Ironie narrative / Point de vue).}
    \begin{enumerate}
        \item Argument : Le narrateur adopte le point de vue des humbles (ironie, focalisation sur Kelara/Meka) pour juger les puissants.
        \item Citation : Description du Commandant « \emph{ventre bedonnant} » (Ch. 1) vs description Meka « \emph{usure} ».
    \end{enumerate}
    \item \textbf{B. La langue française « africainisée » : victoire de l'hyponyme culturel sur l'hypéronyme colonial.}
    \begin{enumerate}
        \item Argument : Syntaxes ewondo, proverbes, rythme oral $\to$ Le français devient langue africaine.
        \item Citation : Structure phrase : « \emph{Meka, il est vieux...} » (Topicalisation) ; Proverbes intraduits.
    \end{enumerate}
    \item \textbf{C. La fin ouverte : Kelara héritière du sens (Métal vs Vie).}
    \begin{enumerate}
        \item Argument : La dernière action (vente médaille) et la dernière image (Kelara partie) scellent la victoire du concret sur le symbole.
        \item Citation : Dernière phrase : « \emph{Kelara prit la médaille... pour aller la vendre.} » (Ch. 6).
    \end{enumerate}
\end{itemize}

\textbf{CONCLUSION} : Bilan : Oui, déconstruction totale. Discours puissants = vide/menace. Parole humbles = plein/vie. Ouverture : Oyono précurseur de la « décolonisation des esprits » (Ngũgĩ wa Thiong'o) — La parole orale/littéraire comme arme de libération.
\end{exoresolu}

\begin{attention}[Les 5 erreurs qui coûtent des points au Bac (Dissertation Littéraire)]
\begin{enumerate}
    \item \textbf{Hors-sujet / Mauvaise problématique :} Ne pas analyser \textbf{tous} les termes du sujet (ex: oublier « satire » ou « assimilation » ou « discuter »). Formuler une problématique fermée (OUI/NON) ou hors tension. \textit{Parade :} Exégèse rigoureuse (surligner, définir, opposer) avant d'écrire.
    \item \textbf{Plan binaire (I/II seulement) ou déséquilibré :} Faire un plan « Pour / Contre » sans synthèse (III). Avoir 3 sous-parties en I et 1 seule en II. \textit{Parade :} Imposer la structure dialectique 3 parties (Thèse/Antithèse/Synthèse) et l'équilibre 2/2/2 ou 3/3/3.
    \item \textbf{Absence de preuves textuelles précises (Le « vide » textuel) :} Dire « Meka est ridicule » sans citer le costume trop grand. Citer de mémoire de façon approximative (guillemets obligatoires, référence chapitre). \textit{Parade :} Constituer une « Banque de citations » (10-15 citations clés classées par thème) et les apprendre par cœur.
    \item \textbf{Paraphrase / Résumé d'intrigue au lieu d'argumentation :} Raconter l'histoire (Meka attend, la cérémonie a lieu, il meurt) au lieu de répondre à la problématique par des arguments analysés. \textit{Parade :} Chaque paragraphe (sous-partie) doit commencer par une \textbf{phrase-argument} (idée force) et non par « Au chapitre 1... ».
    \item \textbf{Négligence de la forme et de la langue :} Pas de transitions (rupture brutale I/II). Orthographe/grammaire défaillante (accords, temps). Présentation illisible (pas d'indentation, pas de saut de ligne pour Intro/Dev/Concl). \textit{Parade :} 10 min de relecture ciblée : 1. Structure (Titres, Transitions) 2. Citations (Guillemets, Exactitude) 3. Langue (Accords, Ponctuation, Vocabulaire précis).
\end{enumerate}
\end{attention}

\newpage

\coursec{Exercices}

\coursec{Leçon 01 : Dissertation — Produire le plan détaillé d'une dissertation littéraire}

\courssub{Lecture de l'œuvre support : paratexte et lecture ouvroir}

\begin{aretenir}[Le Vieux nègre et la médaille — Repères essentiels]
\emph{Le Vieux nègre et la médaille} (1956) est un roman de \textbf{Ferdinand Oyono} (Cameroun, 1929–2010), figure majeure de la littérature africaine d'expression française. C'est une satire féroce de l'administration coloniale française au Cameroun.
\begin{itemize}
\item \textbf{Paratexte} : Le titre associe \og vieux nègre \fg{} (déshumanisation, mépris colonial) et \og médaille \fg{} (leurre, symbole vide). La couverture, la préface et la quatrième de couverture orientent la lecture vers la critique de l'assimilation et de la collaboration.
\item \textbf{Lecture ouvroir} : Découverte du personnage principal, Meka, vieux chef traditionnel naïf qui attend la cérémonie de remise de médaille comme une consécration. L'ironie dramatique s'installe dès l'ouverture : le lecteur perçoit le piège colonial que Meka ne voit pas.
\item \textbf{Thèmes centraux} : La collaboration, l'illusion, la spoliation des terres, la déchéance du chef traditionnel, l'ironie comme arme littéraire.
\end{itemize}
\end{aretenir}

\courssub{Outils de langue I : référent, substituts, synonymie et antonymie}

\begin{definition}[Référent et substituts]
Le \cle{référent} est l'être, l'objet ou l'idée désigné(e) dans l'univers du discours (souvent introduit par un groupe nominal). Pour éviter la répétition lourde, la langue utilise des \cle{substituts} :
\begin{itemize}
\item \textbf{Pronoms} : personnels (\emph{il, elle, le, la, lui, y, en}), relatifs (\emph{qui, que, dont}), démonstratifs (\emph{celui, celle, ce}), possessifs (\emph{le mien, la tienne}), indéfinis (\emph{un, une, certains}).
\item \textbf{Groupes nominaux synonymiques (GNS)} : expressions nominales reprenant le référent par une qualité, une fonction ou une périphrase (\emph{le vieillard, cet homme, le récipiendaire, l'ancien combattant}).
\item \textbf{Nominalisation / Ellipse} : \emph{le premier, l'heureux élu, ce dernier}.
\end{itemize}
La chaîne référentielle assure la cohésion textuelle.
\end{definition}

\begin{definition}[Synonymie et antonymie]
\begin{itemize}
\item \cle{Synonymes} : mots ou expressions de sens très proche, interchangeables dans un contexte donné (ex. \emph{médaille / distinction / décoration / récompense}). Attention aux nuances de registre, de connotation ou de valeur expressive.
\item \cle{Antonymes} : mots de sens opposé. Trois types : \emph{contradictoires} (mort/vivant), \emph{contraires} (chaud/froid — graduels), \emph{inverses} (acheter/vendre — relation réciproque).
\end{itemize}
\end{definition}

\begin{methode}[Repérer les substituts d'un référent dans un texte]
\begin{enumerate}
\item Identifier le nom (ou GN) qui introduit le référent pour la première fois.
\item Relever tous les éléments (pronoms, GNS) qui reprennent ce référent par la suite.
\item Classer chaque substitut selon sa nature grammaticale (pronom personnel, relatif, démonstratif ; GNS).
\item Vérifier l'accord (genre, nombre, personne) entre le référent et ses substituts.
\end{enumerate}
\end{methode}

\courssub{Outils de langue II : hypéronymie/hyponymie, contenus manifestes et latents, texte argumentatif}

\begin{definition}[Hyponymie et hyperonymie]
Relation de \cle{inclusion} sémantique entre un terme général (\cle{hyperonyme} / superordonné) et des termes plus précis (\cle{hyponymes} / subordonnés).
\begin{itemize}
\item \emph{Livre} (hyperonyme) $\rightarrow$ \emph{roman, essai, poème, théâtre} (hyponymes).
\item \emph{Récompense} (hyperonyme) $\rightarrow$ \emph{médaille, décorations, prix, titre} (hyponymes).
\item Test : \og Un \emph{hyponyme} \textbf{est un} \emph{hyperonyme} \fg{} (ex. \og Un roman est un livre \fg{}).
\end{itemize}
\end{definition}

\begin{definition}[Contenus manifestes et contenus latents]
\begin{itemize}
\item \cle{Contenu manifeste} (explicite) : ce qui est dit littéralement, l'information factuelle accessible à la lecture de surface (qui fait quoi, où, quand).
\item \cle{Contenu latent} (implicite) : ce qui est suggéré, sous-entendu, que le lecteur doit inférer (ironie, critique, émotion, idéologie, intention profonde de l'auteur). L'analyse du latent repose sur les indices textuels (choix lexicaux, syntaxe, figures de style, contexte).
\end{itemize}
\end{definition}

\begin{propriete}[Caractéristiques du texte argumentatif]
Un texte argumentatif vise à \textbf{convaincre} (par la raison) ou \textbf{persuader} (par l'affect). Il se repère à :
\begin{itemize}
\item Une \cle{thèse} (position défendue, explicite ou implicite).
\item Des \cle{arguments} (raisons logiques) appuyés sur des \cle{exemples} (faits, citations, références littéraires, historiques, statistiques).
\item Une organisation structurée par des \cle{connecteurs logiques} (\emph{donc, car, en effet, cependant, par conséquent, en outre}).
\item Des \cle{modalisateurs} marquant l'engagement de l'énonciateur : verbes d'opinion (\emph{je pense que, il est évident que}), adverbes (\emph{certainement, sans doute}), temps/modalités, figures de style (ironie, hyperbole).
\end{itemize}
\end{propriete}

\courssub{Analyser le sujet et formuler la problématique}

\begin{methode}[Analyse d'un sujet de dissertation — 5 étapes]
\begin{enumerate}
\item \textbf{Délimiter le thème} : identifier le domaine général (ex. littérature engagée, injustice, fidélité, école, citoyenneté).
\item \textbf{Repérer les mots-clés} : termes forts à définir et à problématiser (ex. \emph{arme, efficace, littérature engagée, citoyen, fidélité, récompense}).
\item \textbf{Dégager les présupposés} : ce que le sujet suppose acquis (ex. \og la littérature engagée existe \fg{}, \og l'injustice est une réalité \fg{}, \og la fidélité mérite récompense \fg{}).
\item \textbf{Reformuler la question} : en ses propres mots pour vérifier la compréhension et éviter le hors-sujet.
\item \textbf{Formuler la problématique} : question \textbf{centrale, ouverte, débattue}, qui structure toute la réflexion. Elle n'est ni une simple reformulation, ni la conclusion, mais le fil conducteur du devoir.
\end{enumerate}
\end{methode}

\courssub{Rechercher et organiser les idées}

\begin{methode}[Brainstorming, tri et hiérarchisation des idées]
\begin{enumerate}
\item \textbf{Lister} sans censure toutes les idées, arguments, exemples, citations qui viennent à l'esprit (mind-mapping, colonnes \og Pour \fg{} / \og Contre \fg{}).
\item \textbf{Regrouper} par affinités thématiques (arguments sociaux, politiques, psychologiques, esthétiques, historiques...).
\item \textbf{Sélectionner} les plus pertinents, variés et étayables par l'œuvre au programme (\emph{Le Vieux nègre et la médaille}) et d'autres références (programme ou culture personnelle).
\item \textbf{Hiéarchiser} : du plus faible (exemple unique, évidence, lieu commun) au plus fort (argument structurel, référence majeure, analyse fine).
\item \textbf{Structurer} selon le type de plan choisi (voir étape 6).
\end{enumerate}
\end{methode}

\courssub{Élaborer et produire le plan détaillé}

\begin{methode}[Types de plans et structure du plan détaillé]
\textbf{Trois types principaux adaptés au sujet :}
\begin{itemize}
\item \cle{Plan thématique} (ou par aspects) : pour sujet \og Montrez que... \fg{}, \og Quels sont les... \fg{}, \og Quelles leçons... \fg{}. Parties = grands thèmes distincts (ex. I. Social, II. Politique, III. Psychologique).
\item \cle{Plan dialectique} (Thèse / Antithèse / Synthèse) : pour sujet question fermée (\og Est-ce que... ? \fg{}, \og Faut-il... ? \fg{}). I. Thèse (arguments pour), II. Antithèse (arguments contre), III. Synthèse (dépassement, nuance, conditionnalité).
\item \cle{Plan analytique} (Causes / Conséquences / Solutions) : pour sujet \og Pourquoi... ? \fg{}, \og Quelles causes / conséquences ? \fg{}. I. Causes, II. Conséquences, III. Perspectives / Remèdes.
\end{itemize}
\textbf{Plan détaillé =} Pour chaque grande partie (I, II, III) :
\begin{itemize}
\item Titre de la partie (affirmatif, nuancé).
\item Au moins 2 sous-parties (A, B) avec titre.
\item Pour chaque sous-partie : \textbf{Argument} + \textbf{Exemple précis} (citation, scène, référence textuelle exacte).
\end{itemize}
\textbf{Introduction} : Amorce $\rightarrow$ Présentation sujet/œuvre $\rightarrow$ Problématique $\rightarrow$ Annonce du plan.
\textbf{Conclusion} : Bilan synthétique $\rightarrow$ Élargissement (autre œuvre, actualité, question philosophique).
\end{methode}

\begin{attention}[Pièges fréquents sanctionnés au Probatoire/Bac]
\begin{itemize}
\item Plan déséquilibré (partie I très développée, II maigre, III absente ou répétitive).
\item Sous-parties qui se chevauchent, se répètent ou ne sont que des exemples sans argument.
\item Synthèse qui répète la thèse au lieu de dépasser l'opposition thèse/antithèse.
\item Conclusion sans élargissement (simple résumé).
\item Absence de problématique ou problématique hors-sujet / fermée.
\item Exemples vagues (\og dans le roman \fg{}) au lieu de références précises (chapitre, scène, citation).
\item Confusion entre plan thématique et plan dialectique (choix inadapté au libellé du sujet).
\end{itemize}
\end{attention}

\courssub{Je vérifie mes connaissances}

\begin{exercice}[QCM : les étapes de la dissertation]
Pour chaque question, entoure la (ou les) bonne(s) réponse(s).
\begin{enumerate}
\item La problématique d'un sujet de dissertation est :
\begin{enumerate}
\item une simple reformulation du sujet ;
\item la question centrale qui oriente toute la réflexion ;
\item la conclusion du devoir.
\end{enumerate}
\item Le plan détaillé comporte :
\begin{enumerate}
\item uniquement les titres des grandes parties ;
\item les parties, les sous-parties, les arguments et les exemples ;
\item la rédaction complète du devoir.
\end{enumerate}
\item Dans la phrase \og Meka reçut une médaille ; \emph{elle} brillait au soleil \fg{}, le mot \emph{elle} est :
\begin{enumerate}
\item le référent ;
\item un substitut du référent ;
\item un hyponyme.
\end{enumerate}
\item \og Médaille \fg{} et \og distinction \fg{} sont :
\begin{enumerate}
\item antonymes ;
\item synonymes ;
\item sans rapport de sens.
\end{enumerate}
\end{enumerate}
\end{exercice}

\begin{exercice}[Vrai ou faux, en justifiant]
Réponds par vrai ou faux, puis justifie ta réponse en une ou deux phrases.
\begin{enumerate}
\item Un plan dialectique (thèse -- antithèse -- synthèse) convient à un sujet formulé sous forme de question fermée appelant une discussion.
\item Dans un texte argumentatif, le contenu latent est ce qui est dit explicitement par l'auteur.
\item \og Roman \fg{} est un hyponyme de \og livre \fg{}.
\item La conclusion d'une dissertation littéraire doit obligatoirement contenir un élargissement.
\end{enumerate}
\end{exercice}

\begin{exercice}[Définitions à compléter]
Recopie et complète chaque définition avec les mots suivants : \emph{hyponyme, problématique, substitut, antonyme, amorce, élargissement}.
\begin{enumerate}
\item Le \ldots\ldots\ldots{} est un mot ou un groupe de mots qui reprend le référent pour éviter sa répétition.
\item La \ldots\ldots\ldots{} est la question essentielle dégagée de l'analyse du sujet.
\item Un \ldots\ldots\ldots{} est un terme dont le sens est contraire à celui d'un autre terme.
\item L'\ldots\ldots\ldots{} est la première phrase de l'introduction, qui met le lecteur en contact avec le sujet.
\item L'\ldots\ldots\ldots{} est l'ouverture finale de la conclusion vers une question plus large.
\item \og Médaille \fg{} est un \ldots\ldots\ldots{} de \og récompense \fg{}.
\end{enumerate}
\end{exercice}

\begin{exercice}[Repérage rapide dans le texte]
On donne l'extrait suivant, inspiré de \emph{Le Vieux nègre et la médaille} : \og Le vieux Meka attendait la cérémonie. Le vieillard souriait. Cet homme croyait enfin à la justice des Blancs. \fg{}
\begin{enumerate}
\item Identifie le référent du passage.
\item Releve les deux substituts qui le remplacent et précise leur nature (pronom, groupe nominal synonymique).
\item Releve un contenu manifeste et un contenu latent du passage.
\end{enumerate}
\end{exercice}

\courssub{Je m'entraîne}

\begin{exercice}[Éviter les répétitions]
Le paragraphe suivant est volontairement répétitif : \og La médaille de Meka est remise à Meka par l'administration. La médaille récompense la fidélité de Meka. Mais la médaille ne change rien à la condition de Meka : la médaille est un leurre. \fg{}
\begin{enumerate}
\item Réécris ce paragraphe en remplaçant les répétitions par des substituts du référent (pronoms, groupes nominaux) et des synonymes appropriés.
\item Souligne dans ta version chaque substitut utilisé et indique son référent.
\end{enumerate}
\end{exercice}

\begin{exercice}[Trier et hiérarchiser des idées]
On donne 12 idées en vrac sur le sujet : \og L'école forme-t-elle encore des citoyens ? \fg{}
\begin{center}
\begin{tabularx}{\linewidth}{|X|X|}\hline
1. L'école transmet le savoir et le savoir-faire. & 2. L'école coûte cher aux familles pauvres. \\ \hline
3. Elle enseigne le respect des règles communes. & 4. Certains diplômés restent sans emploi. \\ \hline
5. Elle apprend à vivre ensemble. & 6. Elle développe l'esprit critique. \\ \hline
7. L'école détourne parfois des métiers manuels. & 8. Elle forme les futurs électeurs. \\ \hline
9. Les programmes négligent parfois la culture locale. & 10. Elle lutte contre l'obscurantisme. \\ \hline
11. La discipline scolaire prépare à la vie civique. & 12. L'école peut reproduire les inégalités sociales. \\ \hline
\end{tabularx}
\end{center}
\begin{enumerate}
\item Regroupe ces idées en deux grands ensembles : les arguments \emph{pour} et les arguments \emph{contre}.
\item Trie chaque ensemble et hiérarchise les arguments du plus faible au plus fort.
\item Déduis-en un plan thématique en trois parties (avec sous-parties) adapté au sujet.
\end{enumerate}
\end{exercice}

\begin{exercice}[Choisir le bon type de plan]
Pour chacun des trois sujets suivants, indique le type de plan le plus adapté (plan thématique, plan dialectique, plan analytique) et justifie ton choix en deux phrases.
\begin{enumerate}
\item \og La littérature engagée est-elle une arme efficace contre l'injustice ? \fg{}
\item \og Montrez que \emph{Le Vieux nègre et la médaille} est une dénonciation de la colonisation. \fg{}
\item \og Quelles leçons le personnage de Meka nous offre-t-il sur les illusions ? \fg{}
\end{enumerate}
\end{exercice}

\begin{exercice}[Contenus manifestes et latents]
On donne cette citation adaptée d'Oyono : \og On lui avait dit que la médaille était en or. Meka la regarda longtemps, puis il sourit : il avait donné ses terres, ses fils et sa vie pour ce petit cercle de métal. \fg{}
\begin{enumerate}
\item Dégage les contenus manifestes de ce passage (ce qui est dit explicitement).
\item Dégage les contenus latents (ce qui est suggéré : ironie, critique de la colonisation).
\item Tire de cette analyse deux arguments utilisables dans une dissertation sur la littérature engagée, en précisant pour chacun l'exemple qui l'illustre.
\end{enumerate}
\end{exercice}

\courssub{Je me prépare au Bac}

\begin{exercice}[Sujet type Baccalauréat : dissertation complète]
Sujet : \og La littérature engagée est-elle une arme efficace contre l'injustice ? \fg{}
\begin{enumerate}
\item Analyse le sujet : délimite le thème, repère les mots-clés, dégage les présupposés et reformule la question en tes propres termes.
\item Formule la problématique sous forme d'une question principale accompagnée de deux questions secondaires.
\item Recherche les idées : propose au moins six arguments (trois pour la thèse, trois pour l'antithèse), chacun accompagné d'un exemple tiré de \emph{Le Vieux nègre et la médaille} ou d'une autre œuvre de ton programme.
\item Élabore le plan détaillé complet : introduction (amorce, présentation du sujet, problématique, annonce du plan), développement en trois parties avec sous-parties, arguments et exemples, conclusion (bilan et élargissement).
\end{enumerate}
\end{exercice}

\begin{exercice}[Introduction rédigée]
Sujet : \og Faut-il récompenser la fidélité ? \fg{}, en t'appuyant sur \emph{Le Vieux nègre et la médaille} de Ferdinand Oyono.
\begin{enumerate}
\item Analyse le sujet et formule la problématique.
\item Propose deux types de plan possibles pour ce sujet, puis choisis celui que tu défendras en justifiant ton choix.
\item Rédige l'introduction complète du devoir : amorce, présentation du sujet et de l'œuvre, problématique, annonce claire du plan en trois parties.
\item Rédige ensuite la conclusion : bilan de la réflexion et élargissement vers une autre œuvre ou une situation de la vie courante.
\end{enumerate}
\end{exercice}

\begin{exercice}[Corriger un plan fautif]
Un élève propose le plan détaillé suivant pour le sujet : \og Le roman africain doit-il dénoncer les injustices sociales ? \fg{}
\begin{center}
\begin{tabularx}{\linewidth}{|l|X|}\hline
\textbf{Partie} & \textbf{Contenu annoncé} \\ \hline
I. Le roman africain dénonce les injustices. & Sous-partie 1 : il critique la colonisation (ex. Oyono). Sous-partie 2 : il dénonce les injustices sociales (ex. la pauvreté). Sous-partie 3 : il critique les colons. \\ \hline
II. Le roman africain ne dénonce pas toujours. & Sous-partie 1 : il peut divertir. \\ \hline
III. Synthèse. & Le roman africain dénonce les injustices. \\ \hline
Conclusion. & Le roman dénonce les injustices sociales. \\ \hline
\end{tabularx}
\end{center}
\begin{enumerate}
\item Releve au moins quatre défauts de ce plan (déséquilibre des parties, sous-parties qui se chevauchent, synthèse sans apport nouveau, conclusion sans élargissement).
\item Explique pour chaque défaut pourquoi il serait pénalisé à l'examen.
\item Propose une version remédiée du plan détaillé : parties équilibrées, sous-parties distinctes, arguments et exemples précis, synthèse réelle et conclusion avec élargissement.
\end{enumerate}
\end{exercice}

\newpage

\coursec{Corrigés des exercices}

\begin{corrige}[Exercice 1 — QCM : les étapes de la dissertation]
\textbf{1. Réponse b)} La problématique est la \cle{question centrale} qui oriente toute la réflexion. Elle n'est ni une simple reformulation du sujet (réponse a, qui resterait au niveau de la paraphrase), ni la conclusion du devoir (réponse c, qui livre une réponse alors que la problématique pose une question ouverte).

\textbf{2. Réponse b)} Le plan détaillé comporte les \cle{parties}, les \cle{sous-parties}, les \cle{arguments} et les \cle{exemples}. Les seuls titres des grandes parties (réponse a) constituent un plan sommaire, insuffisant ; la rédaction complète (réponse c) est l'étape qui suit le plan, non le plan lui-même.

\textbf{3. Réponse b)} \emph{Elle} est un \cle{substitut} du référent : c'est un pronom personnel qui reprend le groupe nominal \og une médaille \fg{} (féminin singulier, d'où l'accord). Le référent est \emph{médaille} ; \emph{elle} n'est ni le référent lui-même (a), ni un hyponyme (c), qui serait un terme plus précis inclus dans un terme général.

\textbf{4. Réponse b)} \og Médaille \fg{} et \og distinction \fg{} sont des \cle{synonymes} : ils désignent tous deux une récompense honorifique et sont interchangeables dans ce contexte (avec une nuance : \emph{distinction} est plus général et plus soutenu). Ils ne sont pas antonymes (a), car leurs sens ne s'opposent pas.

\textbf{Points de vigilance :} bien distinguer \emph{référent} (la réalité désignée) et \emph{substitut} (le mot qui la reprend) ; ne pas confondre synonymie (proximité de sens) et hyponymie (inclusion de sens).
\end{corrige}

\begin{corrige}[Exercice 2 — Vrai ou faux, en justifiant]
\textbf{1. Vrai.} Le plan dialectique (thèse -- antithèse -- synthèse) convient aux sujets formulés comme une \cle{question fermée} appelant une discussion (\og Est-ce que\ldots ? \fg{}, \og Faut-il\ldots ? \fg{}) : il permet d'examiner les arguments pour, puis contre, avant de dépasser l'opposition dans la synthèse.

\textbf{2. Faux.} C'est le \cle{contenu manifeste} qui est dit explicitement (l'information littérale, factuelle). Le \cle{contenu latent} est l'implicite : ce qui est suggéré, sous-entendu, et que le lecteur doit inférer à partir d'indices textuels (ironie, choix lexicaux, contexte).

\textbf{3. Vrai.} \og Roman \fg{} est un \cle{hyponyme} de \og livre \fg{}, car le sens de \emph{roman} est inclus dans celui de \emph{livre}. Le test le confirme : \og Un roman \textbf{est un} livre \fg{} est une phrase correcte ; l'inverse (\og Un livre est un roman \fg{}) est faux.

\textbf{4. Vrai.} Au Probatoire comme au Baccalauréat technique, la conclusion doit comporter un \cle{bilan} synthétique de la réflexion \emph{et} un \cle{élargissement} (ouverture vers une autre œuvre, une situation de la vie courante ou une question plus large). Une conclusion qui se contente de répéter la thèse est pénalisée.

\textbf{Points de vigilance :} la justification est exigée : une réponse \og vrai/faux \fg{} non justifiée ne rapporte pas tous les points.
\end{corrige}

\begin{corrige}[Exercice 3 — Définitions à compléter]
\begin{enumerate}
\item Le \textbf{substitut} est un mot ou un groupe de mots qui reprend le référent pour éviter sa répétition.
\item La \textbf{problématique} est la question essentielle dégagée de l'analyse du sujet.
\item Un \textbf{antonyme} est un terme dont le sens est contraire à celui d'un autre terme.
\item L'\textbf{amorce} est la première phrase de l'introduction, qui met le lecteur en contact avec le sujet.
\item L'\textbf{élargissement} est l'ouverture finale de la conclusion vers une question plus large.
\item \og Médaille \fg{} est un \textbf{hyponyme} de \og récompense \fg{} (test : \og une médaille est une récompense \fg{}).
\end{enumerate}
\textbf{Points de vigilance :} ne pas confondre \emph{hyponyme} (terme précis, subordonné) et \emph{hyperonyme} (terme général, superordonné) ; ne pas écrire \og synonyme \fg{} à la question 6 : médaille et récompense ne sont pas interchangeables, il y a inclusion de sens.
\end{corrige}

\begin{corrige}[Exercice 4 — Repérage rapide dans le texte]
\textbf{1. Le référent} est \emph{Meka}, introduit par le groupe nominal \og Le vieux Meka \fg{}.

\textbf{2. Les deux substituts :}
\begin{itemize}
\item \og \textbf{Le vieillard} \fg{} : \cle{groupe nominal synonymique} (GNS), qui reprend Meka par une de ses qualités (l'âge).
\item \og \textbf{Cet homme} \fg{} : \cle{groupe nominal synonymique} construit avec un adjectif démonstratif, qui reprend le référent par sa nature (être humain de sexe masculin).
\end{itemize}
Aucun pronom n'est employé ici ; les deux substituts sont donc tous deux des GNS.

\textbf{3. Contenus :}
\begin{itemize}
\item \cle{Contenu manifeste} : un vieil homme nommé Meka attend une cérémonie ; il sourit ; il croit à la justice des Blancs. Ce sont des faits littéralement énoncés.
\item \cle{Contenu latent} : l'adverbe \emph{enfin} et le verbe \emph{croyait} suggèrent que cette foi est naïve et sera déçue : le passage installe une \cle{ironie dramatique} — le lecteur devine que la \og justice des Blancs \fg{} est une illusion coloniale. C'est le contenu implicite, à inférer.
\end{itemize}
\textbf{Points de vigilance :} citer précisément les mots du texte ; pour le latent, s'appuyer sur un indice textuel identifiable (\emph{enfin}, \emph{croyait}) et non sur une impression vague.
\end{corrige}

\begin{corrige}[Exercice 5 — Éviter les répétitions]
\textbf{1. Paragraphe réécrit (version de référence) :}

\og La médaille de Meka lui est remise par l'administration. Cette \textbf{distinction} récompense sa fidélité. Mais \textbf{elle} ne change rien à la condition du \textbf{vieillard} : cette \textbf{décoration} n'est qu'un leurre. \fg{}

\textbf{2. Substituts et synonymes utilisés :}
\begin{center}
\begin{tabularx}{\linewidth}{|l|X|l|}\hline
\rowcolor{popblueL}\textbf{Substitut} & \textbf{Nature} & \textbf{Référent} \\ \hline
\emph{lui} & pronom personnel (COI) & Meka \\ \hline
\emph{Cette distinction} & GNS + synonyme & la médaille \\ \hline
\emph{sa} & adjectif possessif (renvoi à Meka) & Meka \\ \hline
\emph{elle} & pronom personnel (sujet) & la médaille \\ \hline
\emph{le vieillard} & groupe nominal synonymique & Meka \\ \hline
\emph{cette décoration} & GNS + synonyme & la médaille \\ \hline
\end{tabularx}
\end{center}
\textbf{Justification :} les synonymes choisis (\emph{distinction}, \emph{décoration}) appartiennent au champ lexical de la récompense honorifique et conviennent au contexte ; les pronoms respectent l'accord en genre et en nombre avec leur référent (\emph{elle} reprend \emph{médaille}, féminin singulier). Le texte gagne en fluidité sans perdre en clarté : la chaîne référentielle reste cohérente.

\textbf{Points de vigilance :} vérifier que chaque substitut renvoie sans ambiguïté à un seul référent ; éviter d'alterner trop de synonymes au point de brouiller la lecture ; accorder correctement les pronoms.
\end{corrige}

\begin{corrige}[Exercice 6 — Trier et hiérarchiser des idées]
\textbf{1. Regroupement :}
\begin{itemize}
\item \textbf{Arguments POUR} (l'école forme des citoyens) : 1 (savoir et savoir-faire), 3 (respect des règles), 5 (vivre ensemble), 6 (esprit critique), 8 (futurs électeurs), 10 (lutte contre l'obscurantisme), 11 (discipline et vie civique).
\item \textbf{Arguments CONTRE} (limites de l'école) : 2 (coût pour les pauvres), 4 (diplômés sans emploi), 7 (détournement des métiers manuels), 9 (négligence de la culture locale), 12 (reproduction des inégalités).
\end{itemize}

\textbf{2. Hiérarchisation (du plus faible au plus fort) :}
\begin{itemize}
\item \emph{Pour} : 1 (évidence, lieu commun : transmettre des savoirs n'est pas proprement former un citoyen) $\rightarrow$ 3 et 11 (socialisation aux règles) $\rightarrow$ 5 (vivre ensemble) $\rightarrow$ 8 (formation civique de l'électeur) $\rightarrow$ 6 et 10 (esprit critique, combat contre l'obscurantisme : le cœur de la citoyenneté éclairée).
\item \emph{Contre} : 2 et 4 (obstacles matériels : coût, chômage — limites conjoncturelles) $\rightarrow$ 7 et 9 (défauts d'orientation et de contenu) $\rightarrow$ 12 (reproduction des inégalités : objection structurelle, la plus forte car elle contredit la mission même de l'école).
\end{itemize}

\textbf{3. Plan thématique en trois parties (proposition) :}
\begin{itemize}
\item \textbf{I. L'école transmet les fondements de la vie citoyenne.} A. Elle dispense savoirs et règles communes (idées 1, 3, 11). B. Elle apprend le vivre ensemble (idée 5).
\item \textbf{II. L'école forme la conscience critique du citoyen.} A. Elle prépare aux responsabilités civiques (idée 8). B. Elle développe l'esprit critique contre l'obscurantisme (idées 6, 10).
\item \textbf{III. Mais l'école réelle présente des limites qui affaiblissent sa mission.} A. Limites matérielles et sociales (idées 2, 4, 12). B. Limites pédagogiques et culturelles (idées 7, 9).
\end{itemize}
\textbf{Points de vigilance :} chaque sous-partie doit reposer sur un \emph{argument} distinct, non sur un simple exemple ; vérifier que les parties ne se chevauchent pas et que la troisième nuance sans détruire les deux premières. \textbf{Note :} si le sujet est une question fermée (\og L'école forme-t-elle le citoyen ? \fg{}), un plan dialectique (Thèse/Antithèse/Synthèse) est souvent plus attendu ; le plan thématique ci-dessus reste valide s'il répond à une commande \og Quels sont les rôles et les limites de l'école dans la formation du citoyen ? \fg{}.
\end{corrige}

\begin{corrige}[Exercice 7 — Choisir le bon type de plan]
\textbf{1. \og La littérature engagée est-elle une arme efficace contre l'injustice ? \fg{} — Plan dialectique.} Le sujet est une \cle{question fermée} (\og est-elle\ldots ? \fg{}) qui appelle une discussion : on peut répondre oui (la littérature dénonce, éveille, mobilise) et non (elle reste sans effet immédiat, censurée, lue par une minorité). La synthèse permettra de dépasser l'opposition (efficacité réelle mais indirecte et à long terme).

\textbf{2. \og Montrez que \emph{Le Vieux nègre et la médaille} est une dénonciation de la colonisation. \fg{} — Plan thématique.} Le sujet est un \cle{impératif démonstratif} (\og Montrez que\ldots \fg{}) : la thèse est imposée, il n'y a pas à discuter sa validité mais à l'illustrer par des aspects distincts (dénonciation politique : la spoliation des terres ; dénonciation sociale : le mépris et l'humiliation ; dénonciation psychologique : l'illusion et la naïveté de Meka).

\textbf{3. \og Quelles leçons le personnage de Meka nous offre-t-il sur les illusions ? \fg{} — Plan thématique (par aspects).} La question ouverte en \og quelles \fg{} invite à inventorier des leçons distinctes (l'illusion du pouvoir, l'illusion de la récompense, l'illusion de la justice coloniale), chacune formant une partie. Un plan dialectique serait inadapté car le sujet ne demande pas de discuter une opinion mais d'analyser un personnage.

\textbf{Points de vigilance :} c'est le \emph{libellé} du sujet qui commande le type de plan ; repérer les marques : question fermée $\rightarrow$ dialectique ; \og montrez que / quelles / quels \fg{} $\rightarrow$ thématique ; \og pourquoi / quelles causes, quelles conséquences \fg{} $\rightarrow$ analytique.
\end{corrige}

\begin{corrige}[Exercice 8 — Contenus manifestes et latents]
\textbf{1. Contenus manifestes (explicites) :} on avait annoncé à Meka que la médaille était en or ; Meka la regarde longuement et sourit ; il a donné ses terres, ses fils et sa vie en échange de ce petit objet de métal. Ce sont les faits directement énoncés.

\textbf{2. Contenus latents (implicites) :}
\begin{itemize}
\item \og On lui avait dit que\ldots \fg{} : le discours rapporté suggère le \cle{mensonge} — rien ne garantit que la médaille soit en or ; l'administration trompe.
\item \og \textbf{petit} cercle de \textbf{métal} \fg{} : la dégradation (or promis $\rightarrow$ simple métal) et la petitesse de l'objet dévalorisent la récompense face à l'énormité du sacrifice (\emph{terres, fils, vie} : gradation ascendante).
\item Le sourire de Meka, contrastant avec la réalité de la spoliation, installe une \cle{ironie amère} : la colonisation repose sur l'illusion consentie par ses victimes.
\end{itemize}

\textbf{3. Deux arguments utilisables en dissertation :}
\begin{itemize}
\item \textbf{Argument 1 :} la littérature engagée dénonce l'injustice en révélant la spoliation coloniale. \emph{Exemple :} Meka a cédé ses terres et perdu ses fils pour une médaille sans valeur — Oyono montre que la colonisation échange des symboles vides contre des biens réels.
\item \textbf{Argument 2 :} la littérature engagée éveille les consciences par l'ironie. \emph{Exemple :} le sourire de Meka devant le \og petit cercle de métal \fg{} provoque chez le lecteur la pitié et la révolte, et l'amène à juger le système colonial.
\end{itemize}
\textbf{Points de vigilance :} distinguer rigoureusement ce qui est \emph{écrit} (manifeste) de ce qui est \emph{inféré} (latent) ; chaque argument doit être accompagné d'un exemple textuel précis.
\end{corrige}

\begin{corrige}[Exercice 9 — Sujet type Baccalauréat : dissertation complète]
\textbf{1. Analyse du sujet.}
\begin{itemize}
\item \textbf{Thème :} la fonction sociale et politique de la littérature (littérature engagée).
\item \textbf{Mots-clés :} \emph{littérature engagée} (littérature qui prend parti, défend une cause), \emph{arme} (métaphore : moyen de combat), \emph{efficace} (qui produit un effet réel), \emph{injustice} (violation du droit et de l'équité).
\item \textbf{Présupposés :} la littérature engagée existe ; l'injustice est une réalité contre laquelle on peut lutter ; une œuvre littéraire peut agir sur le réel.
\item \textbf{Reformulation :} une œuvre qui dénonce les injustices parvient-elle réellement à les combattre et à transformer la société ?
\end{itemize}

\textbf{2. Problématique.} Question principale : \emph{La littérature engagée dispose-t-elle d'une efficacité réelle dans la lutte contre l'injustice, ou son pouvoir reste-t-il symbolique ?} Questions secondaires : par quels moyens propres la littérature combat-elle l'injustice ? Quelles sont les limites de son action sur le réel ?

\textbf{3. Recherche des idées (six arguments).}
\begin{itemize}
\item \emph{Thèse (efficacité) :} A1. Elle dénonce et révèle les injustices cachées (ex. Oyono révèle la spoliation des terres de Meka derrière la cérémonie de la médaille). A2. Elle éveille et mobilise les consciences par l'ironie et la satire (ex. le rire amer provoqué par la naïveté de Meka pousse le lecteur à condamner le colonialisme). A3. Elle a une portée historique durable (ex. \emph{Le Vieux nègre et la médaille} (1956) a nourri la conscience anticoloniale africaine ; on peut citer aussi \emph{Les Misérables} de Hugo contre la misère).
\item \emph{Antithèse (limites) :} A4. Son effet est indirect et différé : un roman ne change pas immédiatement une loi (ex. la médaille de Meka n'a rendu ni ses terres ni sa dignité ; l'écriture seule n'a pas libéré les colonies). A5. Elle touche un public limité, souvent déjà convaincu (ex. les œuvres engagées circulent surtout parmi les lettrés). A6. Elle est exposée à la censure et à la récupération (ex. les écrivains engagés ont été surveillés, interdits ou exilés sous les régimes autoritaires).
\end{itemize}

\textbf{4. Plan détaillé complet.}
\begin{itemize}
\item \textbf{Introduction.} Amorce : depuis toujours, l'écrivain intervient dans les débats de son temps. Présentation : la littérature engagée, illustrée au Cameroun par \emph{Le Vieux nègre et la médaille} de Ferdinand Oyono (1956), prétend combattre l'injustice. Problématique (reprise de la question principale). Annonce du plan : nous verrons d'abord que la littérature engagée est une arme réelle, puis que son efficacité connaît des limites, enfin qu'elle agit autrement et à long terme.
\item \textbf{I. La littérature engagée est une arme efficace contre l'injustice.} A. Elle dénonce : elle rend visibles les mécanismes de l'oppression (ex. la satire de l'administration coloniale chez Oyono). B. Elle mobilise : par l'ironie et l'émotion, elle transforme le lecteur en témoin révolté (ex. le sourire de Meka devant le \og petit cercle de métal \fg{}).
\item \textbf{II. Cependant, cette arme connaît des limites.} A. Son action est indirecte et différée : l'écrit ne supprime pas l'injustice immédiate (ex. la médaille ne rend rien à Meka). B. Son audience est restreinte et elle subit la censure (ex. les œuvres anticoloniales longtemps marginalisées).
\item \textbf{III. Son efficacité, réelle mais spécifique, est celle des consciences.} A. Elle agit en profondeur et durablement, en formant le jugement des générations (ex. la postérité d'Oyono dans la conscience africaine). B. Elle complète, sans les remplacer, les autres formes de lutte (politique, juridique, sociale).
\item \textbf{Conclusion.} Bilan : la littérature engagée est une arme efficace, non au sens d'une action immédiate, mais parce qu'elle éveille et durablement transforme les consciences. Élargissement : aujourd'hui encore, romans, chansons et réseaux numériques poursuivent ce combat — l'art peut-il se passer de l'engagement ?
\end{itemize}

\textbf{Barème indicatif (sur 20) :} analyse du sujet et problématique : 4 pts ; pertinence et précision des arguments et exemples : 6 pts ; cohérence et équilibre du plan : 6 pts ; qualité de l'introduction et de la conclusion (annonce du plan, élargissement) : 4 pts.

\textbf{Points de vigilance :} problématique ouverte et liée aux mots-clés ; trois parties équilibrées ; synthèse qui dépasse l'opposition (ni répétition de la thèse, ni simple compromis) ; exemples précis (auteur, titre, scène) ; conclusion avec élargissement obligatoire.
\end{corrige}

\begin{corrige}[Exercice 10 — Introduction rédigée]
\textbf{1. Analyse et problématique.}
\begin{itemize}
\item \textbf{Thème :} la récompense et ses valeurs. \textbf{Mots-clés :} \emph{récompenser} (sanctionner positivement un mérite), \emph{fidélité} (constance dans l'attachement et le service).
\item \textbf{Présupposé :} la fidélité est un mérite qui appelle une récompense.
\item \textbf{Problématique :} \emph{Toute fidélité mérite-t-elle d'être récompensée, et toute récompense honore-t-elle réellement celui qui la reçoit ?}
\end{itemize}

\textbf{2. Deux types de plan possibles.}
\begin{itemize}
\item \emph{Plan dialectique :} I. Il faut récompenser la fidélité (la récompense reconnaît le mérite, encourage l'engagement). II. Mais certaines récompenses sont des leurres ou des instruments de domination (la médaille de Meka). III. Seule une récompense juste et désintéressée honore la fidélité.
\item \emph{Plan thématique :} I. Les valeurs de la fidélité. II. Les fonctions de la récompense. III. Les dérives de la récompense.
\item \textbf{Choix justifié :} le plan \cle{dialectique} est le plus adapté, car le sujet est une question fermée (\og Faut-il\ldots ? \fg{}) qui appelle une discussion ; il permet d'exploiter pleinement le cas paradoxal de Meka, récompensé et pourtant spolié.
\end{itemize}

\textbf{3. Introduction rédigée (modèle).}

\og Dans toutes les sociétés, les hommes ont institué des récompenses pour honorer ceux qui servent avec constance. Mais que vaut une distinction lorsque celui qui la décerne est aussi celui qui spolie ? C'est toute la question que pose Ferdinand Oyono dans \emph{Le Vieux nègre et la médaille} (1956), roman satirique où le vieux Meka, chef traditionnel camerounais, reçoit une médaille de l'administration coloniale en échange d'une vie de fidélité — et de ses terres. On peut alors se demander si toute fidélité mérite d'être récompensée et si toute récompense honore véritablement son récipiendaire. Nous montrerons d'abord que la fidélité appelle légitimement la récompense, puis que certaines récompenses ne sont que des leurres au service du pouvoir, enfin que seule une récompense juste et désintéressée donne à la fidélité son vrai prix. \fg{}

\textbf{4. Conclusion rédigée (modèle).}

\og La fidélité, parce qu'elle engage toute une vie, semble naturellement appeler la reconnaissance. Mais l'exemple de Meka nous a montré que la récompense peut devenir un instrument de mystification lorsqu'elle masque une spoliation : la médaille ne rend ni les terres ni la dignité. Seule une récompense sincère, proportionnée au sacrifice et libre de tout calcul de domination, honore réellement la fidélité. Cette leçon dépasse le roman d'Oyono : dans nos sociétés contemporaines, décorations officielles et distinctions de façade n'invitent-elles pas à toujours interroger la valeur réelle des honneurs et la sincérité de ceux qui les décernent ? \fg{}

\textbf{Barème indicatif (sur 20) :} analyse et problématique : 4 pts ; choix du plan justifié : 3 pts ; introduction (amorce, présentation, problématique, annonce du plan) : 7 pts ; conclusion (bilan, élargissement) : 6 pts.

\textbf{Points de vigilance :} l'introduction doit contenir les quatre étapes dans l'ordre ; l'annonce du plan doit reprendre fidèlement les trois parties ; l'élargissement ne doit pas introduire un nouveau développement mais ouvrir une question ; citer correctement l'œuvre (auteur, titre en italique, date).
\end{corrige}

\begin{corrige}[Exercice 11 — Corriger un plan fautif]
\textbf{1. et 2. Défauts relevés et sanctions encourues :}
\begin{itemize}
\item \textbf{Déséquilibre des parties :} la partie I compte trois sous-parties, la partie II une seule, la partie III aucune. À l'examen, un plan déséquilibré traduit une réflexion inaboutie : l'antithèse et la synthèse sont sacrifiées, ce qui coûte des points de cohérence.
\item \textbf{Sous-parties qui se chevauchent :} dans la partie I, \og critiquer la colonisation \fg{}, \og dénoncer les injustices sociales \fg{} et \og critiquer les colons \fg{} disent trois fois la même chose. La répétition révèle l'absence d'arguments distincts et est pénalisée comme du remplissage.
\item \textbf{Synthèse sans apport nouveau :} la partie III répète mot pour mot la thèse de la partie I. Or la synthèse doit \emph{dépasser} l'opposition thèse/antithèse ; une synthèse copiée-collée est sanctionnée comme une absence de synthèse.
\item \textbf{Conclusion sans élargissement :} la conclusion se contente de répéter la thèse ; elle omet l'ouverture exigée (autre œuvre, actualité, question plus large), ce qui fait perdre les points réservés à l'élargissement.
\item On peut ajouter : les exemples sont vagues (\og la pauvreté \fg{} n'est ni une œuvre ni une référence précise).
\end{itemize}

\textbf{3. Version remédiée du plan détaillé.}
\begin{itemize}
\item \textbf{I. Le roman africain assume une fonction de dénonciation des injustices sociales.} A. Il dénonce la spoliation coloniale : Oyono montre dans \emph{Le Vieux nègre et la médaille} Meka dépossédé de ses terres contre une médaille sans valeur. B. Il dénonce l'humiliation et le mépris : la satire de l'administration et la déchéance du chef traditionnel exposent l'injustice vécue au quotidien.
\item \textbf{II. Cependant, le roman africain ne se réduit pas à la dénonciation.} A. Il peut divertir et célébrer : le conte, le roman d'amour ou d'aventures visent le plaisir et la transmission des valeurs culturelles (ex. les récits merveilleux de la tradition orale romancée). B. La dénonciation a ses limites : l'écriture seule ne supprime pas l'injustice et touche un public restreint.
\item \textbf{III. Le roman africain dépasse la dénonciation pour construire une conscience nouvelle.} A. En révélant l'injustice, il éveille durablement les lecteurs : l'ironie d'Oyono transforme le rire en révolte et en jugement. B. Il ouvre des perspectives : proposer des modèles de dignité et de résistance vaut autant que dénoncer (ex. les figures de résistance dans le roman africain contemporain).
\item \textbf{Conclusion.} Bilan : le roman africain dénonce effectivement les injustices, mais sa mission est plus large : éveiller, célébrer et construire. Élargissement : à l'ère des réseaux sociaux, la littérature partage aujourd'hui cette fonction critique avec de nouveaux médias — reste à savoir si elle garde un pouvoir propre.
\end{itemize}

\textbf{Barème indicatif (sur 20) :} repérage d'au moins quatre défauts justifiés : 8 pts ; qualité de la remédiation (équilibre, distinction des sous-parties, synthèse réelle) : 8 pts ; conclusion avec élargissement : 4 pts.

\textbf{Points de vigilance :} nommer précisément chaque défaut (ne pas écrire \og le plan est mauvais \fg{}) ; dans la version remédiée, vérifier que chaque sous-partie contient un argument \emph{et} un exemple précis ; la synthèse doit apporter une idée nouvelle qui dépasse l'opposition.
\end{corrige}

\newpage

\coursec{Fiche bilan}

\begin{popfigure}
\centering
\begin{tikzpicture}[mindmap, grow cyclic,
  every node/.style={concept, text=white, font=\small},
  root concept/.append style={concept color=popdark, font=\bfseries},
  level 1/.append style={level distance=3.6cm, sibling angle=60},
  level 1 concept/.append style={font=\footnotesize}]
\node[root concept]{Dissertation littéraire : le plan détaillé}
  child[concept color=popblue]{node[align=center]{Lecture de l'œuvre\\ \emph{Le Vieux nègre et la médaille}\\ paratexte, lecture ouvroir}}
  child[concept color=poporange]{node[align=center]{Langue I\\ référent et substituts\\ synonymie / antonymie}}
  child[concept color=popgreen]{node[align=center]{Langue II\\ hypéronymie / hyponymie\\ contenus manifestes et latents}}
  child[concept color=poppurple]{node[align=center]{Texte argumentatif\\ thèse, arguments, exemples\\ connecteurs logiques}}
  child[concept color=popgold]{node[align=center]{Analyse du sujet\\ mots-clés, type de sujet\\ problématique}}
  child[concept color=poppink]{node[align=center]{Recherche des idées\\ remue-méninges\\ classement, tri}};
\end{tikzpicture}
\legende{Carte mentale de la leçon : les six étapes qui mènent de la lecture de l'œuvre à la production du plan détaillé de la dissertation.}
\end{popfigure}

\begin{bilan}
\textbf{Définitions clés}
\begin{itemize}
  \item \cle{Dissertation} : exercice écrit de réflexion personnelle et argumentée sur un sujet littéraire, organisé en introduction, développement et conclusion.
  \item \cle{Problématique} : question centrale (souvent en deux ou trois interrogations liées) qui naît de l'analyse du sujet et dirige tout le devoir.
  \item \cle{Plan détaillé} : présentation écrite des parties, sous-parties, arguments et exemples du devoir, avec les transitions.
  \item \cle{Référent} : réalité (personne, objet, notion) désignée par un mot ; ses \cle{substituts} sont les reprises qui le désignent à nouveau (pronoms, synonymes, hyperonymes).
  \item \cle{Synonymie} : relation entre mots de sens proche (\emph{vieillard} / \emph{ancien}) ; \cle{antonymie} : relation entre mots de sens contraire (\emph{justice} / \emph{injustice}).
  \item \cle{Hyperonyme} : terme général (\emph{fruit}) ; \cle{hyponyme} : terme particulier (\emph{mangue}).
  \item \cle{Contenu manifeste} : ce que le texte dit explicitement ; \cle{contenu latent} : ce qu'il suggère ou sous-entend.
  \item \cle{Texte argumentatif} : texte qui défend une thèse à l'aide d'arguments (raisons) et d'exemples (illustrations), reliés par des connecteurs logiques.
\end{itemize}

\textbf{Repères à connaître par cœur}
\begin{center}
\begin{tabularx}{\linewidth}{|l|X|}
\hline
\rowcolor{popblueL} \textbf{Notion} & \textbf{À retenir} \\
\hline
Types de sujet & Sujet de réflexion (question ouverte), sujet de discussion (affirmation à discuter), sujet de commentaire (citation à expliquer). \\
\hline
Types de plan & Plan thématique (aspects du sujet), plan dialectique (thèse / antithèse / synthèse), plan analytique (constat / causes / conséquences-solutions). \\
\hline
Structure du devoir & Introduction (amorce, présentation, problématique, annonce du plan) ; développement (2 ou 3 parties, chacune : idée, argument, exemple) ; conclusion (bilan, ouverture). \\
\hline
Connecteurs utiles & D'abord, ensuite, en effet, cependant, par conséquent, enfin, ainsi. \\
\hline
\end{tabularx}
\end{center}

\textbf{Méthodes express}
\begin{enumerate}
  \item Analyser le sujet : souligner les mots-clés, définir chacun, repérer le type de sujet.
  \item Formuler la problématique : transformer la tension du sujet en une ou deux questions.
  \item Rechercher les idées : remue-méninges, puis trier, regrouper et hiérarchiser.
  \item Choisir le type de plan adapté au sujet et répartir les idées en 2 ou 3 parties équilibrées.
  \item Rédiger le plan détaillé : intitulés de parties et sous-parties, arguments, exemples tirés de l'œuvre, transitions.
\end{enumerate}
\end{bilan}

\begin{aretenir}[Checklist avant le Bac]
\begin{itemize}
  \item $\square$ Je sais définir la dissertation, la problématique et le plan détaillé.
  \item $\square$ Je connais le paratexte du \emph{Vieux nègre et la médaille} (titre, auteur, quatrième de couverture) et j'ai achevé la lecture ouvroir.
  \item $\square$ Je sais repérer un référent et ses substituts dans un texte.
  \item $\square$ Je distingue synonymie et antonymie, hyperonymie et hyponymie.
  \item $\square$ Je sais distinguer contenu manifeste et contenu latent d'un énoncé.
  \item $\square$ Je reconnais les caractéristiques du texte argumentatif : thèse, arguments, exemples, connecteurs.
  \item $\square$ Je sais dégager les mots-clés d'un sujet et identifier son type.
  \item $\square$ Je sais formuler une problématique claire en une ou deux questions.
  \item $\square$ Je sais choisir entre plan thématique, dialectique et analytique.
  \item $\square$ Mon plan détaillé comporte des parties équilibrées, avec arguments et exemples.
  \item $\square$ Je prévois les transitions entre les parties et une conclusion avec ouverture.
  \item $\square$ Je gère mon temps : analyse du sujet et plan avant toute rédaction.
\end{itemize}
\end{aretenir}

\findecours

\end{document}
